\documentclass[10pt,a4paper,draft]{article}

\usepackage{csquotes}
\usepackage[a4paper, margin=2.5cm]{geometry}
\usepackage{enumitem}                                  % list config
\usepackage{tocloft}                                   % ToC
\usepackage{amssymb,amsmath,amsthm,stmaryrd,mathtools} % fonts, symbols
\usepackage{xcolor}                                    % \textcolor
\usepackage[framemethod=tikz]{mdframed}                % theorem styling
\usepackage{unicode-math}                              % fonts
\usepackage{memoize}                                   
\usepackage{wrapfig}
\usepackage{booktabs}                                  % cmidrules
\usepackage{tikz-cd}                                   % diagrams
\usepackage[backend=biber,style=alphabetic]{biblatex}  % citations
\usepackage{imakeidx}                                  % index
\usepackage{hyperref}
\usepackage{newunicodechar}

\usepackage{todonotes} %TODO: enlever une fois fini

\makeindex

% Fonts
\setmainfont{fbb}
\setmathfont{Libertinus Math}
%FIXME: Trouver mieux que STIX Two Math
\setmathfont{STIX Two Math}[range={\obslash, \oslash, \sharp, \flat}]

% Shortcuts for math fonts
\newcommand{\bb}[1]{\mathbb{#1}}
\newcommand{\ca}[1]{\mathcal{#1}}
\newcommand{\fk}[1]{\mathfrak{#1}}
\newcommand{\sr}[1]{\mathscr{#1}}
\newcommand{\cat}[1]{\textcat{#1}}
\newcommand{\textcat}[1]{{\normalfont\textsf{#1}}}

% Utilities
\newcommand{\flip}[1]{\text{\rotatebox[origin=c]{-180}{$#1$}}}

% Shortcuts
% FIXME: blacktriangleleft and right too short, cf setmathfont
\newcommand{\PP}{{\bb P}}
\newcommand{\To}{\Rightarrow}
\newcommand{\coce}[2]{{\hat{#1}}_{#2}}
\newcommand{\coc}[2]{{\widehat{#1}}(#2)}
\newcommand{\colimeq}{\mathrel{\hat\approx}}
\newcommand{\colimleq}{\mathrel{\hat\lesssim}}
\newcommand{\copsh}[2]{{\widebreve{#1}}[{#2}]}
\newcommand{\dfrom}{\nvleftarrow}
\newcommand{\from}{\leftarrow}
\newcommand{\ip}{\flip\pi}
\newcommand{\iset}[1]{\llbracket #1 \rrbracket}
\newcommand{\lext}{\mathbin{\rhd}}
\newcommand{\lift}{\mathbin{\lhd}}
\newcommand{\nvequal}{=\!\!\!|\!\!\!=}
\newcommand{\op}{\mathrm{op}}
\newcommand{\oy}{\flip\yo}
\newcommand{\plext}{\mathbin{{\rhd}\mathclap{\mspace{-17.5mu}\cdot}}}
\newcommand{\plift}{\mathbin{\flip{{\rhd}\mathclap{\mspace{-17.5mu}\cdot}}}}
\newcommand{\prext}{\mathbin{{\blacktriangleright}\mathclap{\mspace{-10.5mu}\raisebox{-0.45pt}{\textcolor{white}{$\cdot$}}}}}
\newcommand{\prift}{\mathbin{\flip{{\blacktriangleright}\mathclap{\mspace{-10.5mu}\raisebox{-0.45pt}{\textcolor{white}{$\cdot$}}}}}}
\newcommand{\psh}[2]{{\widehat{#1}}[{#2}]}
\newcommand{\rext}{\mathbin{\blacktriangleright}}
\newcommand{\rift}{\mathbin{\blacktriangleleft}}
\newcommand{\set}[1]{\left\{#1\right\}}
\newcommand{\simeqq}{\cong}
\newcommand{\tp}[1]{\langle#1\rangle}
\newcommand{\unit}{{\tp{}}}

% Operators
\DeclareMathOperator{\Alg}{\cat Alg}
\DeclareMathOperator{\Opalg}{\cat Opalg}
\DeclareMathOperator{\bSpan}{\ca{S}\mathit{pan}}
\DeclareMathOperator{\Dist}{\bb{D}\textbf{ist}}
\DeclareMathOperator{\Fun}{Fun}
\DeclareMathOperator{\Hom}{Hom}
\DeclareMathOperator{\Lan}{Lan}
\DeclareMathOperator{\Mod}{\bb{M}\textbf{od}}
\DeclareMathOperator{\Span}{\bb{S}\textbf{pan}}
\DeclareMathOperator{\Spec}{Spec}
\DeclareMathOperator{\cart}{cart}
\DeclareMathOperator{\cod}{cod}
\DeclareMathOperator{\colim}{colim}
\DeclareMathOperator{\co}{co}
\DeclareMathOperator{\id}{id}
\DeclareMathOperator{\ob}{ob}
\DeclareMathOperator{\opcart}{opcart}

%FIXME: pullbacks not looking good
% Adjunctions, horizontal arrows, pullback
\usetikzlibrary{decorations.pathmorphing}
\usetikzlibrary{decorations.markings}
\tikzset{%
    adj/.style={%
        /tikz/commutative diagrams/phantom,
        "\dashv"{rotate=-90}
    },
    loose/.default=0.5ex,
    loose/.style={%
        /utils/exec=\tikzset{every node/.append style={outer sep=0.8ex}},
        postaction=decorate,
        decoration={markings, mark=at position 0.5 with {\draw[-] (0,#1) -- (0,-#1);}}
    },
    pullback/.style={
        /tikz/commutative diagrams/phantom,
        "\scalebox{1.5}{$\lrcorner$}"{pos=0.125}
    }
}

% `よ`
\DeclareFontFamily{U}{min}{}
\DeclareFontShape{U}{min}{m}{n}{<-> udmj30}{}
\newcommand\yo{\!\text{\usefont{U}{min}{m}{n}\symbol{'207}}\!}

% widearc and widebreve
\makeatletter
\def\widebreve{\mathpalette\wide@breve}
\def\wide@breve#1#2{\sbox\z@{$#1#2$}%
     \mathop{\vbox{\m@th\ialign{##\crcr
\kern0.08em\brevefill#1{0.8\wd\z@}\crcr\noalign{\nointerlineskip}%
                    $\hss#1#2\hss$\crcr}}}\limits}
\def\brevefill#1#2{$\m@th\sbox\tw@{$#1($}%
  \hss\resizebox{#2}{\wd\tw@}{\rotatebox[origin=c]{90}{\upshape(}}\hss$}
\makeatother
\makeatletter
\def\widearc{\mathpalette\wide@arc}
\def\wide@arc#1#2{\sbox\z@{$#1#2$}%
     \mathop{\vbox{\m@th\ialign{##\crcr
\kern0.08em\arcfill#1{0.8\wd\z@}\crcr\noalign{\nointerlineskip}%
                    $\hss#1#2\hss$\crcr}}}\limits}
\def\arcfill#1#2{$\m@th\sbox\tw@{$#1($}%
  \hss\resizebox{#2}{\wd\tw@}{\rotatebox[origin=c]{-90}{\upshape(}}\hss$}
\makeatother

% Bijection of cells
\newenvironment{iffseq}{%
    \global\let\externaldblbackslash\\
    \begin{array}{cl}
    \ifundef{\internaldblbackslash}{%
        \global\let\internaldblbackslash\\%
        \gdef\\{\internaldblbackslash\cmidrule{1-1}\morecmidrules\cmidrule{1-1}}%
    }{}
}{%
    \end{array}
    \global\undef\internaldblbackslash
    \global\let\\\externaldblbackslash
}

% memoize config
\mmzset{prefix=diagrams/,mkdir,auto={tikzcd}{memoize,verbatim}}

% Frames for math env
\usetikzlibrary{calc}
\tikzset{
    envframe/.style={%
        line width=.5pt,
        draw=black,
    },
}
\mdfdefinestyle{defstyle}{% 
    linecolor=black,
    linewidth=1pt,
    frametitlerule=true,
    frametitlerulewidth=.5pt,
    frametitlebelowskip=0pt,
    innertopmargin=2pt,
    frametitlealignment={\noindent\hspace*{-11pt}},
    hidealllines=true,
    singleextra={%
      \path let \p1=(P), \p2=(O) in (\x2,\y1-\mdfframetitleboxtotalheight) coordinate (D);
      \path [envframe] ($(O)+(25pt,0)$) -| ($(D) + (5pt,0)$);
      % \path [envframe] (D) -- ($(D)+(\mdfframetitleboxdepth,0)$);
  },
  firstextra={% %FIXME: Bancal...
      \path let \p1=(P), \p2=(O) in (\x2,\y1-\mdfframetitleboxtotalheight) coordinate (D);
      \path [envframe] ($(O)+(5pt,0)$) -- ($(D) + (5pt,-18pt)$);
  },
  middleextra={%
      \path let \p1=(P), \p2=(O) in (\x2,\y1) coordinate (D);
      \path [envframe] ($(O)+(5pt,0)$) -- ($(D) + (5pt,0)$);
  },
  secondextra={%
      \path let \p1=(P), \p2=(O) in (\x2,\y1) coordinate (D);
      \path [envframe] ($(O)+(25pt,0)$) -| ($(D) + (5pt,0)$);
  },
} 

% Environements
\theoremstyle{plain}
\newtheorem{theorem}{Theorem}[subsection]
\newtheorem{proposition}[theorem]{Proposition}
\newtheorem{corollary}[theorem]{Corollary}
\newtheorem{lemma}[theorem]{Lemma}

\theoremstyle{definition}
\mdtheorem[style=defstyle]{definition}[theorem]{Definition}

\theoremstyle{remark}
\newtheorem{remark}[theorem]{Remark}
\newtheorem{example}[theorem]{Example}
\newtheorem{notation}[theorem]{Notation}

% Table of contents
\renewcommand{\contentsname}%
    {Contents~\hfill\normalsize{\textbf{Page}}}%
\renewcommand\cftaftertoctitle{\par\noindent\hrulefill\par\vskip-.3em}
\renewcommand{\cftsecleader}{\cftdotfill{\cftdotsep}}


% hyperlinks
\definecolor{ocre}{RGB}{0, 113, 188}
\hypersetup{
    colorlinks,
    citecolor=ocre,
    filecolor=ocre,
    linkcolor=ocre,
    urlcolor=ocre
}

% lists (fix babel)
\setlist[itemize]{label=\textbullet}
\setlist[enumerate]{label=(\alph*)}

% citations
\addbibresource{main.bib}



\begin{document}

\title{Formal Category Theory and Relative Monads}
\author{Luc Saccoccio--Le Guennec}
\date{\today}

\maketitle

\tableofcontents
\noindent\hrulefill

\section{Introduction}

\subsection{Outline of the thesis}

We introduce in a self-contained way the methods of formal category theory from \cite{arkor2024formaltheoryrelativemonads}, \cite{arkor2024nervetheoremrelativemonads}, \cite{arkor2025idempotencerelativemonads}, \cite{arkor2025relativemonadicity} and \cite{arkor2026presheavescocompletionsformalcategory}; and apply the theory of formal preasheaves to the study of the Isbell adjunction, in particular generalizing most of the results in \cite{averyleinster}.

No prerequisites on formal category theory are needed, as the framework is entirely described in the following pages. However, a solid knowledge of category theory (as studied in \cite{riehl_context} for example) is advised, and a basic knowledge of enriched categories (as studied in \cite{kelly}) and $\infty$-categories (an informal introduction suffices, we mention \cite{groth2015shortcourseinftycategories} and the first chapters of \cite{cnossen} but many more exist) is a plus for Section \ref{sec:instances}.

\underline{What is not in this thesis:} We avoid the discussion of the many flavours of skew multicategory/multiactegory used in \cite{arkor2024formaltheoryrelativemonads} and enhanced in \cite{rogers26monoid} as they are not needed for most of the results (and this thesis was already too long). The central results in \cite{arkor2024nervetheoremrelativemonads} do not appear as they require exact virtual equipments. The formal theory of fibered categories also do not appear, due to a lack of time. Finally, (co)ends do not appear: when I started working on the formalization of coends, I stumbled upon \cite{coend_nickel}. I chose to pursue other results, and avidly await Jana Nickel and Nathanael Arkor's article on the subject. %FIXME

\subsection{Conventions}

We will use Sans Serif for 1-categories: \cat{A}, \cat{B}, \cat{C}..., calligraphy for 2-categories or bicategories: $\ca{H}, \ca{J}, \ca{K}$..., blackboard bold for virtual double categories: $\bb{X}, \bb{Y}$... Usually, objects will be in capital letters $A, B, C$, lower case for arrows $f, g, h$ (although horizontal arrows/loose-cells will use letters $p, q, r, s$...) and Greek letters for cells $\alpha, \beta, \phi, \psi$...

As this thesis was first written in French, tight and loose cells are called vertical arrows (or just arrows) and horizontal arrows (or distributors) as in a double category, and we reserve the term "cell" for 2-cells. Every occurrence of "vertical" (resp. "horizontal") can be replaced with tight (resp. "loose") to recover the vocabulary of Arkor and McDermott.

Also, a greater attention to detail than usual will be paid to the dualization of definitions and results in this thesis, as they usually are key notions.

In accordance to N. Arkor and D. McDermott's articles, horizontal arrows will be denoted from right to left ($A \dfrom B$) in diagrams \textit{and in text} for consistency.

\subsection{Acknowledgments}

Arkor et McDermott for articles (+ discussion on Zulip with Arkor), Jana K. Nickel for handwritten notes on conjoint/companion, Morgan Rogers as a tutor, Jonas Frey for discussion on collage

Justin Carel, Hadrien Chalandon-Goskrzynski, Ugo Fialo

Marta Vigano

\section{Formal Category Theory}

We start with the basic notions, in particular "where" we do formal category, followed by "how" or "with what". This section takes from \cite{arkor2024formaltheoryrelativemonads} modernized with \cite{arkor2026presheavescocompletionsformalcategory}, in particular adapting the results for arbitrary weights. A few notions related to (co)limits are also from \cite{arkor2025relativemonadicity} as they are needed to establish a monadicity theorem in Subsection \ref{subsec:rel_monadicity}.

\subsection{Virtual Double Categories}

\begin{definition}
    \index{virtual double category}
    \label{def:vdc}
    A \textbf{virtual double category} $\bb{X}$ comprises of the following data:
    \begin{enumerate}[ref=\ref{def:vdc} (\alph*)]
        \item A category of objects $A, B, C, \dots$ and vertical arrows $f : A \to B, \cdots$. Vertical arrows are also called arrows.
        \item For every pair of objects $A$ and $B$, a class of horizontal arrows $p : A \nvrightarrow B, \cdots$. Horizontal arrows are also called distributors.
        \item \label{def:vdc:cell} For every compatible chain of horizontal arrows $p_1, \cdots, p_n$ ($n \geq 0$), two vertical arrows $f, g$ and $q$ a horizontal arrow, a class of cells:
            \[\begin{tikzcd}[row sep=tiny]
                \bullet \ar[dd, "f"] & \ar[l, "p_1"', loose] \bullet
                                     & \ar[l, "p_2"', loose] \cdots
                                     & \ar[l, "p_{n-1}"', loose] \bullet
                                     & \ar[l, "p_n"', loose] \bullet \ar[dd, "g"]\\
                && \phi && \\
                \bullet &&&& \ar[llll, "p", loose] \bullet
            \end{tikzcd}\]
        \item \label{def:vdc:comp} For every configuration of cells of the right shape with $n\geq 0$ and $\forall i\in\iset{0, n}, a_i\in\bb{N}$:
            \[\begin{tikzcd}[row sep=tiny]
                \bullet \ar[dd, "f"'] & \ar[l, "p_1^1"', loose] \cdots & \ar[l, "p_{a_1}^1"', loose]
                    \bullet \ar[dd] & \cdots & \bullet \ar[dd]
                                    & \cdots \ar[l, "p_1^n"', loose] \cdots & \ar[l, "p_{a_n}^n"', loose] \bullet \ar[dd, "g"]\\
                & \phi_1 & && & \phi_n & \\
                \bullet \ar[dd, "h"'] && \ar[ll, loose] \bullet & \cdots & \bullet && \ar[ll, loose] \bullet \ar[dd, "i"]\\
                &&& \psi &&&\\
                \bullet &&&&&& \ar[llllll, "q", loose] \bullet
            \end{tikzcd}\]
            A cell called the \textbf{composite}:
            \[\begin{tikzcd}[row sep=tiny]
                \bullet \ar[dd, "hf"'] & \ar[l, "p_1^1"', loose] \cdots & \ar[l, "p_{a_n}^n"', loose] \bullet \ar[dd, "ig"]\\
                & \psi(\phi_1, \cdots, \phi_n) & \\
                \bullet && \ar[ll, "q", loose] \bullet
            \end{tikzcd}\]
        \item For each horizontal arrow $p : A \dfrom B$, an identity cell $1_p : p \Rightarrow p$
            \[\begin{tikzcd}[sep=tiny]
                A \ar[dd, equal] && \ar[ll, "p"', loose] B \ar[dd, equal]\\
                                 & = &\\
                A && \ar[ll, "p", loose] B
            \end{tikzcd}\]
    \end{enumerate}
    We ask for associativity:
    \[
        (\psi(\phi_1, \cdots, \phi_n))(\alpha^1_1, ..., \alpha^m_{k_m}) = \psi(\phi_1(\alpha^1_1, \cdots, \alpha^1_{k_1}), \cdots, \phi_m(\alpha^m_1, \cdots, \alpha^m_{k_m}))
    \]
    And unitality:
    \begin{align*}
        \alpha(1_{p_1}, \cdots, 1_{p_n}) & = \alpha & 1_q(\alpha) = \alpha
    \end{align*}
\end{definition}

\begin{remark}
    This notion generalises multiple structures:
    \begin{enumerate}
        \item A \textbf{category} is a virtual double category without horizontal arrows nor cells.
        \item A \textbf{multicategory} (or nonsymmetric colored operad) is a virtual double category with a single object and arrow (the identity)
        \item A \textbf{double category} is a virtual double category where all cells have a single horizontal arrow as domain, and where we can compose horizontal arrows.
    \end{enumerate}
\end{remark}

We choose to think of a virtual double category as a "category" of categories, objects as categories, vertical arrows as functors and horizontal arrows as distributors. We can construct many virtual double category of this type:

\begin{example}
    \index{$\Dist$}
    The virtual double category $\Dist$ has:
    \begin{enumerate}
        \item As underlying category, the (big) category of small categories and functors.
        \item As horizontal arrows, the distributors.
        \item As cells
        \item For cells
            \[\begin{tikzcd}[row sep=tiny]
                \cat{A}_0 \ar[dd, "f"'] & \ar[l, "p_1"', loose] \cdots & \ar[l, "p_n"', loose] \cat{A}_n \ar[dd, "g"]\\
                                       & \phi & \\
                \cat{B} && \ar[ll, "q", loose] \cat{B'}
            \end{tikzcd}\]
            A family of functions $\phi_{a_0, \dots, a_n} : p_1(a_0, a_1) \times \cdots \times p_n(a_{n-1}, a_n) \to q(fa_0, ga_n)$ satisfying extra-naturality conditions that can be found in Subsection \ref{subsec:enriched}.
    \end{enumerate}
\end{example}

Note that $\cat{V}$-categories, $\cat{V}$-functors, $\cat{V}$-distributors and $\cat{V}$-natural transformations also form a virtual double category, although we have to be careful with the definition of $\cat{V}$-distributors: we can ask for a condition of symmetry on $\cat{V}$, or a condition of co and contravariance on each object as in \cite[§8.1]{arkor2024formaltheoryrelativemonads}. They will be studied in Section \ref{subsec:enriched}

\begin{example}
    \label{ex:span}
    For $\cat{E}$ a category with pullbacks, we can construct $\Span(\cat{C})$ with :
    \begin{enumerate}
        \item Underlying category $\cat{E}$.
        \item Horizontal arrows $B \dfrom C$ the spans $B \leftarrow A \rightarrow C$.
        \item Cells
            \[\begin{tikzcd}
                A_0 \ar[d, "f"'] & \ar[l, "p_1"] X_1 \ar[r, "q_1"] & A_1 & \ar[l, "p_2"] \cdots \ar[r, "q_n"] & A_n \ar[d, "g"]\\
                B && \ar[ll, "r"] Y \ar[rr, "s"] && C
            \end{tikzcd}\]
            A span morphism
            \[\begin{tikzcd}
                A_0 \ar[d, "f"'] & \ar[l, "p_1"'] X_1 & \ar[l, "\pi_1"'] X_1 \times_{A_1} \cdots \times_{A_n} X_n \ar[d, "h"] \ar[r, "\pi_n"]
                                 & X_n \ar[r, "q_n"] & A_n \ar[d, "g"]\\
                B && \ar[ll, "r"] Y \ar[rr, "s"] && C
            \end{tikzcd}\]
    \end{enumerate}
\end{example}

\begin{remark}
    \index{virtual double category!dual}
    Contrary to double categories or 2-categories, there's a unique notion of dual fir virtual double categories because of the asymmetry between the horizontal domain and codomain: we will denote $\bb{X}^{\co}$ the dual of a virtual double category, having the same objects and arrows but distributors reversed, and a cell with frame
    \[\begin{tikzcd}
        A_0 \ar[d, "f_0"'] & \ar[l, loose, "p_1"'] \cdots & \ar[l, loose, "p_n"'] A_n \ar[d, "f_n"]\\
        B_0 && \ar[ll, loose, "q"] B_n
    \end{tikzcd}\]
    Are the cells in $\bb{X}$ with frame
    \[\begin{tikzcd}
        A_n \ar[d, "f_n"'] & \ar[l, loose, "p_n"'] \cdots & \ar[l, loose, "p_1"'] A_0 \ar[d, "f_0"]\\
        B_n && \ar[ll, loose, "q"] B_0
    \end{tikzcd}\]
\end{remark}

% En bons catégoriciens, on assemble les catégories virtuelles doubles, on assemble les catégories virtuelles doubles en une 2-catégorie :
In good categoricians, virtual double categories are assembled into a 2-category:

\begin{definition}\hfill
    \begin{itemize}
        \item A \textbf{functor} $F : \bb{X} \to \bb{Y}$ of virtual double categories comprises of the data of a functor on underlying categories, an assignment on horizontal arrows and cells strictly preserving identities and cell composition.
        \item For $F, G : \bb{X} \to \bb{Y}$ two virtual double category functors, a \textbf{natural transformation} $\varphi : F \To G$ comprises of a natural transformation on the underlying functor of underlying categories and, for each horizontal arrow $p : X \dfrom Y$ of $\bb{X}$, a cell
            \[\begin{tikzcd}[sep=tiny]
                FX \ar[dd, "\varphi_X"'] && \ar[ll, "Fp"', loose] FY \ar[dd, "\varphi_Y"]\\
                                        & \varphi_p &\\
                GX && \ar[ll, "Gp", loose] GY
            \end{tikzcd}\]
            Verifying a naturality condition: for each cell
            \[\begin{tikzcd}[sep=small]
                X_0 \ar[dd, "f"'] & \ar[l, "p_1"', loose] \cdots & \ar[l, "p_n"', loose] X_n \ar[dd, "g"] \\
                                 & \xi & \\
                X && \ar[ll, "p", loose] Y
            \end{tikzcd}\]
            We have
            \[\begin{tikzcd}[row sep=tiny]
                FX_0 \ar[dd, "Ff"'] & \ar[l, "Fp_1"', loose] \cdots & \ar[l, "Fp_n"', loose] FX_n \ar[dd, "Fg"] \\
                                 & F\xi & \\
                FX \ar[dd, "\phi_X"'] && \ar[ll, "Fp", loose] FY \ar[dd, "\phi_Y"]\\
                                 & \phi_p & \\
                GX && \ar[ll, "Gp", loose] GY
            \end{tikzcd}
            \hspace{1em}=\hspace{1em}
            \begin{tikzcd}[row sep=tiny, column sep=small]
                FX_0 \ar[dd, "\phi_{X_0}"'] && \ar[ll, "Fp_1"', loose] X_1 \ar[dd, "\phi_{X_1}"] & \cdots & FX_{n-1} \ar[dd, "\phi_{X_{n-1}}"'] && \ar[ll, "Fp_n"', loose] FX_n \ar[dd, "\phi_{X_n}"] \\
                                            & \phi_{p_1} &&&& \phi_{p_n} \\
                GX_0 \ar[dd, "Gf"'] && \ar[ll, "Gp_1", loose] GX_1 & \cdots & GX_{n-1} && \ar[ll, "Gp_n", loose] GX_n \ar[dd, "Gg"]\\
                                    &&& G\xi && \\
                GX &&&&&& \ar[llllll, "Gp", loose] GY
            \end{tikzcd}\]
    \end{itemize}
\end{definition}

\begin{proposition}
    \index{$\ca{VDbl}$}
    The virtual double categories, their functors and natural transformations form a 2-category $\ca{VDbl}$.
\end{proposition}

Moreover, we know distributors on small categories can be composed using a coend. We can define composition through its universal property in a virtual double category:

\begin{definition}
    \index{cell!opcartesian}
    A cell
    \[\begin{tikzcd}[row sep=tiny]
        \cdot \ar[dd, equal] & \ar[l, "q_1"', loose] \cdots & \ar[l, "q_m"', loose] \cdot \ar[dd, equal] \\
                             & \opcart &\\
        \cdot && \ar[ll, "q", loose] \cdot
    \end{tikzcd}\]
    Is \textbf{opcartesian} if every cell with frame
    \[\begin{tikzcd}[row sep=tiny]
        \cdot \ar[dd, "f"'] & \ar[l, "p_1"', loose] \cdots & \ar[l, "p_l"', loose] \cdot
                           & \ar[l, "q_1"', loose] \cdots & \ar[l, "q_m"', loose] \cdot
                           & \ar[l, "r_1"', loose] \cdots & \ar[l, "r_n"', loose] \cdot \ar[dd, "g"] \\
                           &&& \phi &&&\\
        \cdot &&&&&& \ar[llllll, "s", loose] \cdot
    \end{tikzcd}\]
    Factorizes uniquely as
    \[\begin{tikzcd}[row sep=tiny]
        \cdot \ar[dd, equal] & \ar[l, "p_1"', loose] \cdots & \ar[l, "p_l"', loose] \cdot \ar[dd, equal]
                            & \ar[l, "q_1"', loose] \cdots & \ar[l, "q_m"', loose] \cdot \ar[dd, equal]
                           & \ar[l, "r_1"', loose] \cdots & \ar[l, "r_n"', loose] \cdot \ar[dd, equal] \\
                           & = && \opcart && = & \\
        \cdot \ar[dd, "f"'] & \ar[l, "p_1", loose] \cdots & \ar[l, "p_l", loose] \cdot
                           && \ar[ll, "q", loose] \cdot
                           & \ar[l, "r_1", loose] \cdots & \ar[l, "r_n", loose] \cdot \ar[dd, "g"] \\
                           &&& \tilde{\phi} &&& \\
        \cdot &&&&&& \ar[llllll, "s", loose] \cdot
    \end{tikzcd}\]
    $q$ is the \textbf{composite} $q_1 \odot \cdots \odot q_m$ of $q_1, \dots, q_m$. For $m = 0$, $q : A \dfrom A$ is the \textbf{horizontal identity} of $A$, noted $A(1, 1)$, or $\nvequal$ in diagrams.

    In the cases where we only verify the property for $l = 0$ and $f = \id$ (resp. $n = 0$ and $g = \id$), we speak of \textbf{left-opcartesian cell} (resp. right-) and \textbf{left-composite} (resp. right-), noted $\odot_L$ (resp. $\odot_R$).
\end{definition}

The universal property gives an identity or a composite (eventually left- or right-biased) and associativity (idem) up to an inversible cell; when the formers exists. In the case of the composite, asking for equality is too strict. In the case of identity however, there's often a 'canonical choice'.

\begin{definition}
    \index{virtual double category!normal}
    A virtual double category is \textbf{normal} if it is equipped with a choice of horizontal identities.
\end{definition}

\begin{notation}
    We will write $e_A$ the cell associated to the horizontal identity $A(1, 1)$ when it exists:
    \[\begin{tikzcd}[column sep=tiny]
        & \ar[dl, equal] A \ar[dr, equal] &\\
        A & {} & \ar[ll, equal, loose] A
        \ar[from=1-2, to=2-2, phantom, "e_A"{anchor=center}]
    \end{tikzcd}\]
    When the vertical arrows of a cell are the identity, we will write $\varphi : p_1, \cdots, p_n \To q$ and when the horizontal domain is empty, $\varphi : \To q$ (e.g. $e_A : \To A(1, 1)$).
\end{notation}

\begin{remark}
    \label{rem:id_f}
    Let $f : A \to B$ be a vertical arrow. If $B(1, 1)$ exists, and taking $n = 0$ in \ref{def:vdc:comp}, we can compose the diagram on the left to get the one on the right:
    \[\begin{tikzcd}[column sep=tiny]
        & A \ar[d, "f"] &\\
        & B \ar[dl, equal] \ar[dr, equal] & \\
        B & {} & \ar[ll, equal, loose] B
        \ar[from=2-2, to=3-2, phantom, "e_B"{anchor=center}]
    \end{tikzcd}
    \hspace{1em} = \hspace{1em}
    \begin{tikzcd}[column sep=tiny]
        & A \ar[dl, "f"'] \ar[dr, "f"] &\\
        B & {} & \ar[ll, equal, loose] B
        \ar[from=1-2, to=2-2, phantom, "e_f"{anchor=center}]
    \end{tikzcd}\]
    We think of $e_f$ as an identity cell for $f$. If $A(1, 1)$ exists, we can also apply the universal property of $A(1, 1)$ to get:
    \[\begin{tikzcd}[column sep=tiny]
        & A \ar[dl, equal] \ar[dr, equal] &\\
        A \ar[d, "f"'] & {} & \ar[ll, equal, loose] A \ar[d, "f"]\\
        B & {} & \ar[ll, equal, loose] B
        \ar[from=1-2, to=2-2, phantom, "e_A"{anchor=center}]
        \ar[from=2-2, to=3-2, phantom, "\widetilde{e_f}"{anchor=center}]
    \end{tikzcd}\]
\end{remark}

\begin{notation}
    When there's no risk of confusion, we will write $\varphi : f \To g$ a cell with frame
    \[\begin{tikzcd}[column sep=tiny]
        & A \ar[dl, "f"'] \ar[dr, "g"] &\\
        B & {} & \ar[ll, equal, loose] B
        \ar[from=1-2, to=2-2, phantom, "\varphi"{anchor=center}]
    \end{tikzcd}\]
\end{notation}

\begin{proposition}
    \label{prop:uX-2cat}
    When $\bb{X}$ is normal, we can form a 2-category $\underline{\bb{X}}$ from the underlying category of $\bb{X}$ and with 2-morphisms the cells $f \To g$ of $\bb{X}$ and identities $\id_f : f \To f$ the cell $e_f : f \To f$ of Remark \ref{rem:id_f}.
\end{proposition}

Moreover, when considering distributors, it is possible to restrict them: from a distributor $P : \cat{A} \dfrom \cat{B}$ and two functors $F : \cat{C} \to \cat{A}, G : \cat{D} \to \cat{B}$, we can form $P(F-, G-)$. Once again we define restriction through its universal property:

\begin{definition}
    \index{cell!cartesian}\index{restriction}\index{horizontal arrow!representable}\index{companion}\index{conjoint}
    A cell
    \[\begin{tikzcd}[sep=small]
        \cdot \ar[dd, "f"'] && \ar[ll, "p"', loose] \cdot \ar[dd, "g"]\\
                            & \text{cart} & \\
        \cdot && \ar[ll, "q", loose] \cdot
    \end{tikzcd}\]
    is \textbf{cartesian} if every cell with frame as in left factors as in the right:
    \[\begin{tikzcd}[row sep=tiny]
        \cdot \ar[dd, "gf"'] & \ar[l, "r_1"', loose] \cdots & \ar[l, "r_n"', loose] \cdot \ar[dd, "ih"]\\
                            & \phi & \\
        \cdot && \ar[ll, "q", loose] \cdot
    \end{tikzcd}
    \hspace{1em}=\hspace{1em}
    \begin{tikzcd}[row sep=tiny]
        \cdot \ar[dd, "f"'] & \ar[l, "r_1"', loose] \cdots & \ar[l, "r_n"', loose] \cdot \ar[dd, "h"]\\
                           & \tilde{\phi} & \\
        \cdot \ar[dd, "g"'] && \ar[ll, "p", loose] \cdot \ar[dd, "i"]\\
                            & \text{cart} &\\
        \cdot && \ar[ll, "q", loose] \cdot
    \end{tikzcd}\]
    We write $q(f, g) := p$ and call it the \textbf{restriction}. When it exists, the restriction $A(1, g)$ (resp. $A(f, 1)$) is called the \textbf{companion} (resp. \textbf{conjoint}) of $f$ (resp. $g$). A horizontal arrow is \textbf{representable} (resp. \textbf{corepresentable}) when it is the companion (resp. conjoint) or a vertical arrow.
\end{definition}

\begin{remark}
    Once again, the universal property gives restriction up to an inversible cell.
\end{remark}

When the restrictions exists, we have a useful result:

\begin{theorem}[{\cite[Theorem 7.16]{cs10}}]\label{th:comp_rest}
    Let $p : B \dfrom C$, $f : A \to B$ and $g : D \to C$. Then $B(f, 1) \odot p \odot C(1, g)$ is isomorph to $p(f, g)$ when one of the sides exists.
\end{theorem}

\begin{notation}
    Let $f : A \to B$, we suppose that $A(1, 1), B(1, 1), B(1, f)$ et $B(f, 1)$ exist. Using the universal property of the restriction, we can factors $e_f$ through $B(1, f)$ and $B(f, 1)$:
    \[\begin{tikzcd}
        A \ar[d, "f"'] & \ar[l, equal] A \ar[d, equal]\\
        B \ar[d, equal] & \ar[l, loose, "{B(1, f)}"', ""{anchor=center,name=t}] A \ar[d, "f"]\\
        B & \ar[l, equal, loose, ""{anchor=center, name=b}] B
        \ar[from=t, to=b, phantom, "\cart"{anchor=center}]
    \end{tikzcd}
    \hspace{1em}=\hspace{1em}
    \begin{tikzcd}
        \ar[d, "f"'] A \ar[r, equal] & A \ar[d, "f"]\\
        B \ar[r, equal, loose] & B
    \end{tikzcd}
    \hspace{1em}=\hspace{1em}
    \begin{tikzcd}
        A \ar[d, equal] & \ar[l, equal] A \ar[d, "f"]\\
        A \ar[d, "f"'] & \ar[l, loose, "{B(f, 1)}"', ""{anchor=center,name=t}] B \ar[d, equal]\\
        B & \ar[l, equal, loose, ""{anchor=center, name=b}] B
        \ar[from=t, to=b, phantom, "\cart"{anchor=center}]
    \end{tikzcd}\]
    We define the cells $\smallsmile_f$ as the composite
    \[\begin{tikzcd}
        B \ar[d, equal] & \ar[l, loose, "{B(1, f)}"'] A \ar[d, "f"] & \ar[l, loose, "{B(f, 1)}"'] B \ar[d, equal]\\
        B & \ar[l, equal, loose] B & \ar[l, equal, loose] B
    \end{tikzcd}
    \hspace{1em} = \hspace{1em}
    \begin{tikzcd}
        B \ar[d, equal] & \ar[l, loose, "{B(1, f)}"'] A & \ar[l, loose, "{B(f, 1)}"'] B \ar[d, equal]\\
        B & {} & \ar[ll, equal, loose] B
        \ar[from=1-2, to=2-2, phantom, "\smallsmile_f"{anchor=center}]
    \end{tikzcd}\]
    Furthermore, knowing that $B(f, f) \simeq B(f, 1) \odot B(1, f)$ when one of the sides exists, we define $\smallfrown_f$:
    \[\begin{tikzcd}
        A \ar[d, equal] & \ar[l, loose, equal] A \ar[d, "f"] & \ar[l, loose, equal] A \ar[d, equal]\\
        A \ar[d, equal] & \ar[l, "{B(f, 1)}"', loose] B & \ar[l, "{B(1, f)}"', loose] A \ar[d, equal]\\
        A & {} & \ar[ll, loose, "{B(f, f)}"] A
        \ar[from=2-2, to=3-2, phantom, "\opcart"{anchor=center}]
    \end{tikzcd}
    =
    \begin{tikzcd}
        A \ar[d, equal] & \ar[l, loose, equal] A & \ar[l, loose, equal] A \ar[d, equal]\\
        A & {} & \ar[ll, loose, "{B(f, f)}"] A
        \ar[from=1-2, to=2-2, phantom, "\smallfrown_f"{anchor=center}]
    \end{tikzcd}\]
\end{notation}

\begin{theorem}[{\cite[Theorem 7.20]{cs10}}]
    From those two cells and factorization, we get a bijection between cells with frames
    \[\begin{tikzcd}
        A \ar[d, "f"'] & \ar[l, loose, "p"', ""{anchor=center,name=t}] C \ar[d, "g"]\\
        B & \ar[l, loose, "q", ""{anchor=center,name=b}] D
        \ar[from=t, to=b, Rightarrow, shorten=2mm]
    \end{tikzcd}
    \hspace{1em}\begin{tikzcd}{}&\ar[l,leftrightsquigarrow]{}\end{tikzcd}\hspace{1em}
    \begin{tikzcd}
        \ar[d, equal] B & \ar[l, loose, "{B(1, f)}"'] A & \ar[l, loose, "p"', ""{anchor=center,name=t}] C & \ar[l, loose, "{D(g, 1)}"'] D\ar[d, equal]\\
        B &&& \ar[lll, loose, "q", ""{anchor=center,name=b}] D
        \ar[from=t, to=b, Rightarrow, shorten=2mm]
    \end{tikzcd}\]
\end{theorem}

\begin{corollary}[{\cite[Corollaries 7.21 and 7.22]{cs10}}]
    \label{cor:bending}
    There's a bijection between cells with frames
    \[\begin{tikzcd}
        A \ar[d, "k"'] & \ar[l, loose, "p"'] D & \ar[l, loose, "{E(f, 1)}"'] E \ar[d, "k"]\\
        B && \ar[ll, loose, "{q(g, 1)}"] F
    \end{tikzcd}
    \hspace{1em}\begin{tikzcd}{}&\ar[l,leftrightsquigarrow]{}\end{tikzcd}\hspace{1em}
    \begin{tikzcd}
        A \ar[d, "k"'] & \ar[l, loose, "p"'] D \ar[d, "f"]\\
        B \ar[d, "g"'] & E \ar[d, "k"]\\
        C & \ar[l, loose, "q"] F
    \end{tikzcd}\]
    Taking $p, q, h, k$ to be units and identities, we get bijection between $f\To g$ and cells with frame:
    \[\begin{tikzcd}
        \ar[d, equal]A & \ar[l, loose, "{B(f, 1)}"'] B\ar[d, equal]\\
        A & \ar[l, loose, "{B(g, 1)}"] B
    \end{tikzcd}
    \hspace{3em}
    \begin{tikzcd}
        \ar[d, equal]B & \ar[l, loose, "{B(1, f)}"'] A\ar[d, equal]\\
        B & \ar[l, loose, "{B(1, g)}"] A
    \end{tikzcd}\]
\end{corollary}

\begin{corollary}
    \label{cor:res_ff}
    In particular, restriction is fully faithful and $f\simeqq g$ if and only if $B(1, f) \simeqq B(1, g)$ if and only if $B(f, 1) \simeqq B(g, 1)$.
\end{corollary}

\begin{definition}
    \index{virtual equipment}\index{virtual equipment!strict}
    \label{def:ve}
    A \textbf{virtual equipment} is a virtual double category in which all horizontal identities and restrictions exist. It is moreover \textbf{strict} if it is equipped with a functorial choice of restrictions : $p(1, 1) = p$ et $p(f, g)(f', g') = p(ff', gg')$.
\end{definition}

\begin{example}[Back to Example \ref{ex:span}]
    For $E$ a category with pullbacks, $\Span(\cat{E})$ is always a normal virtual equipment, where horizontal identities are
    \[\begin{tikzcd}
        & \ar[dl, "\id_A"'] A \ar[dr, "\id_A"] &\\
        A && A
    \end{tikzcd}\]
    And restrictions are given by the pullback
    \[\begin{tikzcd}
        & \ar[dl] X \times_B A \times_C Y \ar[dr] &\\
        X \ar[d] && Y \ar[d]\\
        B & \ar[l] A \ar[r] & C
    \end{tikzcd}\]
\end{example}

\begin{theorem}[{\cite[Theorem A.1]{arkor2024nervetheoremrelativemonads}}]
    Every virtual equipment is equivalent to a strict virtual equipment.
\end{theorem}

We can then, up to equivalence, work in a strict virtual equipment.

\subsection{Extensions, liftings and (co)limits}

We work in a virtual equipment $\bb{X}$. Right extensions and lifts will be a very important, probably the most important tool in a virtual equipment, subsuming every notion. A lot of results develop a \textit{calculus of extensions} (and lifts!), replacing both reasoning internal to categories and coend calculus to an extent. A lot of results in Section \ref{sec:presheaves} were proven using coend calculus, and can be proven in a virtual equipment similarly with extension calculus.

\begin{definition}
    \index{weight}\index{weight!unitary}
    A \textbf{weight} is a chain of horizontal arrows:
    \[\begin{tikzcd}[sep=small]
        X_0 & \ar[l, loose, "p_1"'] X_1 & \ar[l, loose, "p_2"'] \cdots & \ar[l, loose, "p_m"'] X_m
    \end{tikzcd}\]
    We will write $\vec{p} := (p_1, \cdots, p_m)$. A weight is \textbf{unitary} if $m=1$.
\end{definition}

The notion of weight generalizes the one in the enriched case and the one appearing in \cite{arkor2024formaltheoryrelativemonads}.

\begin{definition}
    \index{lift!right}
    Let $\vec{p} : Z \dfrom \cdots \dfrom Y$, $q : Z \dfrom X$. The data of a horizontal arrow $q \rift \vec{p} : Y \dfrom X$ and a cell $\omega : p_1, \cdots, p_m, q \rift \vec{p} \To q$ (the \textbf{counit}) is a \textbf{right lift} of $q$ through $\vec{p}$ if every cell factors uniquely:
    \[\begin{tikzcd}
        Z \ar[dd, equal] & \ar[l, loose, "p_1"'] \cdots & \ar[l, loose, "p_m"'] Y & \ar[l, loose, "r_1"'] \cdots & \ar[l, loose, "r_n"'] X \ar[dd, equal]\\
          && \phi &&\\
        Z &&&& \ar[llll, loose, "q"] X
    \end{tikzcd}
    \hspace{1em} = \hspace{1em}
    \begin{tikzcd}[row sep=tiny]
        Z \ar[dd, equal] & \ar[l, loose, "p_1"'] \cdots & \ar[l, loose, "p_m"'] Y \ar[dd, equal] & \ar[l, loose, "r_1"'] \cdots & \ar[l, loose, "r_n"'] X \ar[dd, equal]\\
                         & = && \tilde{\phi} &\\
        Z \ar[dd, equal] & \ar[l, "p_1"', loose] \cdots & \ar[l, loose, "p_m"'] Y && \ar[ll, loose, "q \rift {\vec{p}}"'] X \ar[dd, equal]\\
                         && \omega && \\
        Z &&&& \ar[llll, loose, "q"] X
    \end{tikzcd}\]
\end{definition}

\begin{remark}
    In other words, we have a diagram and a bijection between cells
    \[\begin{tikzcd}
        Y \ar[dr, "\vec{p}"', ""{name=p, anchor=center}, loose] & \ar[l, loose, "q \rift \vec{p}"'] X \ar[d, loose, "q", ""{name=q, anchor=center}]\\
        \ar[from=p, to=q, Rightarrow, shift left=2, xshift=-1mm]
                                                      & Z
    \end{tikzcd}
    \hspace{10pt}
    \begin{iffseq}
        r_1, \cdots, r_n \To q \rift \vec{p}\\
        \vec{p}, r_1, \cdots, r_n \To q
    \end{iffseq}\]
\end{remark}

Because of the asymmetry between the horizontal domain and codomain of cells, we cannot define left lift, neither can we define left extension of horizontal arrows. Nonetheless, we define right extension as right lift in the dual.

\begin{definition}
    \index{extension!right}
    let $\vec{p} : Y \dfrom \cdots \dfrom X$, $q : Z \dfrom X$. The data of a horizontal arrow $\vec{p} \rext q : Z \dfrom Y$ and a cell $\omega : \vec{p} \rext q, p_1, \cdots, p_m \To q$ (the \textbf{counit}) is a \textbf{right extension} of $q$ along $\vec{p}$ if every cell factors uniquely:
    \[\begin{tikzcd}
        Z \ar[dd, equal] & \ar[l, loose, "r_1"'] \cdots & \ar[l, loose, "r_n"'] Y & \ar[l, loose, "p_1"'] \cdots & \ar[l, loose, "p_m"'] X \ar[dd, equal]\\
          && \phi &&\\
        Z &&&& \ar[llll, loose, "q"] X
    \end{tikzcd}
    \hspace{1em} = \hspace{1em}
    \begin{tikzcd}[row sep=tiny]
        Z \ar[dd, equal] & \ar[l, loose, "r_1"'] \cdots & \ar[l, loose, "r_n"'] Y \ar[dd, equal] & \ar[l, loose, "p_1"'] \cdots & \ar[l, loose, "p_m"'] X \ar[dd, equal]\\
                         & \tilde{\phi} && = &\\
        Z \ar[dd, equal] && \ar[ll, loose, "\vec{p} \rext q"'] Y & \ar[l, "p_1"', loose] \cdots & \ar[l, loose, "p_m"'] X \ar[dd, equal]\\
                         && \omega && \\
        Z &&&& \ar[llll, loose, "q"] X
    \end{tikzcd}\]
\end{definition}

\begin{remark}
    Once again, we have a diagram and a bijection between cells
    \[\begin{tikzcd}
        Y \ar[dr, "\vec{p}"', ""{name=p, anchor=center}, loose] & \ar[l, loose, "\vec{p} \rext q"'] X \ar[d, loose, "q", ""{name=q, anchor=center}]\\
        \ar[from=p, to=q, Rightarrow, shift left=2, xshift=-1mm]
                                                      & Z
    \end{tikzcd}
    \hspace{10pt}
    \begin{iffseq}
        r_1, \cdots, r_n \To \vec{p} \rext q\\
        r_1, \cdots, r_n, \vec{p} \To q
    \end{iffseq}\]
\end{remark}

Juggling with universal properties of extensions and lifts, and with cartesian and opcartesian cells, we find the following properties:

\begin{lemma}
    \label{lem:ext_lift}
    We fix a weight $\vec{p} := (p_1, \cdots, p_m) : Z \dfrom \cdots \dfrom Y$ and $q : Z \dfrom X$. For the dual propositions, we rather consider $\vec{p} : Y \dfrom \cdots \dfrom X$.
    \begin{enumerate}[ref=\ref{lem:ext_lift} (\alph*)]
        \item \label{lem:ext_lift:stab_wcomp} If $p_1 \odot_L \cdots \odot_L p_i$ exists for some $i\in\iset{1, m}$, and if either side exists, then:
            $$q \rift (p_1, \cdots, p_m) \simeq q \rift (p_1 \odot_L \cdots \odot_L p_i ,\cdots, p_m)$$
            Dually, if $p_i \odot_R \cdots \odot_R p_m$ exists, and if either side exists, then:
            $$(p_1, \cdots, p_m) \rext q \simeq (p_1, \cdots, p_i \odot_R \cdots \odot_R p_m) \rext q$$
        \item \label{lem:ext_lift:stab_comp} If $p_i \odot \cdots \odot p_j$ exists for some $i \leq j \in \iset{i, m}$, and if either side exists, then:
            $$q \rift (p_1, \cdots, p_m) \simeq q \rift (p_1, \cdots, p_i \odot \cdots \odot p_j, \cdots, p_m)$$
            Dually, mutatis mutandis :
            $$(p_1, \cdots, p_m) \rext q \simeq (p_1, \cdots, p_i \odot \cdots \odot p_j, \cdots, p_m) \rext q$$
        \item \label{lem:ext_lift:stab_split} If $q \rift (p_1, \cdots, p_i)$ exists for some $i\in\iset{1, m}$, and if either side exists, then:
            $$(q \rift (p_1, \cdots, p_i)) \rift (p_{i+1}, \cdots, p_m) \simeq q \rift(p_1, \cdots, p_m)$$
            Dually :
            $$(p_1, \cdots, p_i) \rext ((p_{i+1}, \cdots, p_m) \rext q) \simeq (p_1, \cdots, p_m) \rext q$$
        \item \label{lem:ext_lift:res} Let $x : X' \to X, y : Y' \to Y$. If $q \rift \vec{p}$ exists, then
            $$(q \rift \vec{p})(y, x) \simeq q(1, x) \rift (p_1, \cdots, p_m(1, y))$$
            Dually, for $z : Z' \to Z$:
            $$(\vec{p} \rext q)(z, y) \simeq (p_1(y, 1), \cdots, p_m)\rext q(z, 1)$$
        \item \label{lem:ext_lift:assoc_rift_rext} Let $\vec{r} : Y \dfrom \cdots \dfrom X$ be a weight. Suppose that $\vec{r} \rext q$ and $q \rift \vec{p}$ exist. and if either side exists, then:
            $$(\vec{p} \rext q) \rift \vec{r} \simeq \vec{p} \rext (q \rift \vec{r})$$
    \end{enumerate}
\end{lemma}

\begin{remark}
    \label{rem:ext_lift:comp}
    From the properties \ref{lem:ext_lift:stab_comp} and \ref{lem:ext_lift:stab_split}, and up to compatibility of horizontal arrows and existence of either side, we have:
    \begin{align*}
        q \rift (p\odot_Lp') &= (q \rift p) \rift p' & (p\odot_Rp') \rext q &= p' \rext (p \rext q)
    \end{align*}
\end{remark}

Using right lift and extensions, we can express the universal property of weighted (co)limits in a virtual equipment:

\begin{definition}
    \index{$\vec{p}$-!cocone}\index{$\vec{p}$-weighted!cocone}\index{$\vec{p}$-!colimit}\index{$\vec{p}$-weighted!colimit}
    \index{$\vec{p}$-!cone}\index{$\vec{p}$-weighted!cone}\index{$\vec{p}$-!limit}\index{$\vec{p}$-weighted!limit}
    \label{def:colim}\label{def:lim}
For $\vec{p} := (p_1, \cdots, p_m) : Y \dfrom \cdots \dfrom J$ a weight and $f : J \to X$ a vertical arrow, a \textbf{$\vec{p}$-cocone} (or \textbf{$\vec{p}$-weighted cocone}) is a pair of a vertical arrow $c : Y \to X$ and a cell $\lambda : \vec{p} \To X(f, c)$. A $\vec{p}$-cocone $(c, \lambda)$ is \textbf{colimiting}, or is a $\vec{p}$-\textbf{colimit} if the following diagram exhibits $X(c, 1)$ as the right lift $X(f, 1) \rift \vec{p}$:
    \[\begin{tikzcd}[sep=huge]
        J \ar[d, equal] & \ar[l, "p_1"', loose] \cdots & \ar[l, "p_m"', loose] Y \ar[d, equal] & \ar[l, "{X(c, 1)}"'{name=t}, loose] X \ar[d, equal]\\
        J \ar[d, equal] & {} & \ar[ll, "{X(f, c)}"{name=0}, loose] Y & \ar[l, "{X(c, 1)}"{name=b}, loose] X \ar[d, equal]\\
        J &&& \ar[lll, "{X(f, 1)}"{name=2}, loose] X
        \ar[from=1-2, to=0, phantom, "\lambda"{anchor=center}] \ar[from=t, to=b, phantom, "="{anchor=center}]
        \ar[from=2-3, to=2-2, phantom, ""{name=1}] \ar[from=1, to=2, phantom, "{\smallsmile_{c}(f, 1)}"{anchor=center}]
    \end{tikzcd}\]
    Dually, for $\vec{p} := (p_1, \cdots, p_m) : J \dfrom \cdots \dfrom Y$, we can define \textbf{$\vec{p}-$cones} as a pair of a vertical arrow $c : Y \to X$ and a cell $\mu : \vec{p} \To X(c, f)$, and a $\vec{p}$-cone $(c, \mu)$ is \textbf{limiting} (or a $\vec{p}$-limit) if the following diagram exhibits $X(1, c)$ as the right extension $\vec{p}\rext X(1, f)$.
    \[\begin{tikzcd}[sep=huge]
        X \ar[d, equal] & \ar[l, "{X(1, c)}"'{name=t}, loose] Y \ar[d, equal] & \ar[l, "p_1"', loose] \cdots & \ar[l, "p_m"', loose] J \ar[d, equal]\\
        X \ar[d, equal] & \ar[l, "{X(1, c)}"{name=b}, loose] Y & {} & \ar[ll, "{X(c, f)}"{name=0}, loose] J \ar[d, equal]\\
        X &&& \ar[lll, "{X(1, f)}"{name=2}, loose] J
        \ar[from=1-3, to=0, phantom, "\mu"{anchor=center}] \ar[from=t, to=b, phantom, "="{anchor=center}]
        \ar[from=2-3, to=2-2, phantom, ""{name=1}] \ar[from=1, to=2, phantom, "{\smallsmile_{c}(1, f)}"{anchor=center}]
    \end{tikzcd}\]
    A colimiting cone is usually written $\colim^{\vec{p}}f$ and a limiting cone is written $\lim^{\vec{p}}f$.
\end{definition}

We can apply the lemmas \ref{lem:ext_lift:stab_wcomp} and \ref{lem:ext_lift:res} to get the following properties:

\begin{lemma}
    \label{lem:co_lim}
    Let $f : J \to X$ be a vertical arrow.
    \begin{enumerate}[ref=\ref{lem:co_lim} (\alph*)]
        \item \label{lem:co_lim:comp} Let $(p, p') : J \dfrom Y' \dfrom Y$ be a weight such that $p \odot_L p' : J \dfrom Y$ and $\colim^pf$ exists. If either side exists, then:
            $$\colim^{p\odot_Lp'}f \simeq \colim^p(\colim^{p'}f)$$
            Dually, for $(p, p') : Y \dfrom Y' \dfrom J$ a weight such that $p \odot_R p' : Y \dfrom J$ and $\lim^{p'}f$ exist. If either side exists, then:
            $$\lim{}^{p\odot_Rp'}f \simeq \lim{}^p(\lim{}^{p'}f)$$
        \item \label{lem:co_lim:res} Let $p : Y \dfrom \cdots \dfrom J$ be a weight such that $\colim^{\vec{p}}f$ exists. Then for all $x : X' \to X, y : Y' \to Y$
            $$X((\colim^{\vec{p}}f)y, x) \simeq X(f, x) \rift (p_1, \cdots, p_m(1, y))$$
            Dually, for $p : J \dfrom \cdots \dfrom Y$ such that $\lim^{\vec{p}}f$ exists and $x : X' \to X, y: Y' \to Y$
            $$X(x, (\lim{}^{\vec{p}}f)y) \simeq (p_1(y, 1), \cdots, p_m) \rext X(x, \lim{}^{\vec{p}}f)$$
    \end{enumerate}
\end{lemma}

\begin{corollary}
    Let $\vec{p} : Z \dfrom \cdots \dfrom Y$ be a weight, $f : Z \to X$ and $g : Y \to X$ vertical arrows such that $\colim^{\vec{p}}f$ et $\lim^{\vec{p}}g$ exist. Then we have a bijection of cells in $\underline{\bb{X}}$:
    \[\begin{iffseq}
        \colim^{\vec{p}}f\To g \\
        f \To \lim^{\vec{p}}g
    \end{iffseq}\]
\end{corollary}

We also define a notion of respect for (co)limits by a horizontal arrow, akin to the definition of preservation of (co)limits by functors. We will in fact deduce the preservation of (co)limits by vertical arrows from it.

\begin{definition}
    \index{respect!of limits}\index{respect!of colimits}
    A horizontal arrow $q : X \dfrom W$ \textbf{respects a colimit} $\colim^{\vec{p}}f$ when the cell $\vec{p}, q(\colim^{\vec{p}}f, 1) \To q(f, 1)$ exhibits $q(\colim^{\vec{p}}f, 1)$ as the right lift $q(f, 1) \rift \vec{p}$.

    \noindent
    Dually, $q = W \dfrom X$ \textbf{respects a limit} $\lim^{\vec{p}}f$ when the cell $q(1, \lim^{\vec{p}}f), \vec{p} \To q(1, f)$ exhibits $q(1, \lim^{\vec{p}}f)$ as the right extension $\vec{p} \rext q(1, f)$.
\end{definition}

\begin{definition}
    \index{limit!preservation}\index{colimit!preservation}
    A vertical arrow $f : X \to Y$ \textbf{preserves a colimit} if its conjoint $X(f, 1) : X \dfrom Y$ respects it. Dually, it \textbf{preserves a limit} if its companion $X(1, f) : Y \dfrom X$ respects it.
\end{definition}

\begin{remark}
    Consequently, a vertical arrow $g : X \to Y$ preserves a colimit $(\colim^{\vec{p}}f, \lambda)$ when the $\vec{p}$-cocone $(g\circ\colim^{\vec{p}}f, (g\lambda))$ is colimiting. Dually, $g : X \to Y$ preserves a limit $(\lim^{\vec{p}}f, \mu)$ when the $\vec{p}$-cone ($g\circ\lim^{\vec{p}}f, (g\mu))$ is limiting.
\end{remark}

However, this is not the only interaction between arrows and (co)limits:

\begin{definition}
    \index{limit!reflection}\index{colimit!reflection}
    An arrow $g : X \to Z$ \textbf{reflects} $\vec p$-colimits if a $\vec p$-cocone $(c, \lambda)$ for $f$ is colimiting only if the $\vec p$-cocone $(gc, g\gamma)$ for $gf$ is colimiting. Dually, it reflects $\vec p$-limits if a $\vec p$-cone for $f$ is limiting only if the induced $\vec p$-cone for $gf$ is limiting.
\end{definition}

\begin{definition}
    \index{limit!lifting}\index{colimit!lifting}
    An arrow $g : X \to Z$ \textbf{lifts} $\vec p$-colimits when, if $(d, \lambda)$ is a colimiting $\vec p$-cocone for $gf$, there exists a colimiting $\vec p$-cocone $(c, \alpha)$ for $f$ such that $(gc, g\alpha)$ is colimiting. Dually, it lifts limits when $(d, \lambda)$ is a limiting $\vec p$-cone for $gf$, there exists a limiting $\vec p$-cone $(c, \alpha)$ for $f$ such that $(gc, g\alpha)$ is limiting.

    It \textbf{stricly lifts} $\vec p$-colimits (reps. $\vec p$-limits) when $(gf, g\alpha) = (d, \lambda)$, and \textbf{uniquely} if the $\vec p$-cocone $(c, \alpha$ is unique with this property.
\end{definition}

\begin{definition}
    \index{limits!creation}\index{colimits!creation}
    An arrow $g : X \to Z$ (resp. strictly) \textbf{creates} $\vec p$-colimits if it both (resp. strictly) lifts and reflects $\vec p$-colimits. Dually, an arrow $g : X \to Z$ (resp. strictly) \textbf{creates} $\vec p$-limits if it both (resp. strictly) lifts and reflects $\vec p$-limits.
\end{definition}

\begin{remark}
    Unravelling the definition, an arrow $g : X \to Z$ non-strictly creates $\vec p$-colimits if, for $f : Z \to W, \vec p : Z \dfrom \cdots \dfrom Y$ such that $\colim^{\vec p}gf$ exists we have:
    \begin{enumerate}
        \item For every $\vec p$-cocone $(c, \lambda)$ for $f$, the $\vec p$-cocone $(gc, g\lambda)$ is a $\vec p$-colimit of $gf$ if and only if $(c, \lambda)$ is colimiting.
        \item $W$ admits the colimit $\colim^{\vec p}f$.
    \end{enumerate}
    The property is unique and strict if for a colimiting $\vec p$-cocone $(d, \lambda)$, there is a unique colimiting $\vec p$-cocone $(c, \alpha)$ such that $d = gc$ and $\lambda = g\alpha$.
\end{remark}

\begin{definition}
    \index{colimit!$j$-absolute} \index{limit!$j$-absolute}
    Let $\vec p : Z \dfrom \cdots \dfrom Y, j : A \dfrom E, f : Z \to E$ such that $(\colim^{\vec p} f, \lambda)$ exists. This colimit is \textbf{$j$-absolute} when the following cell is left opcartesian:
    \[\begin{tikzcd}[sep=huge]
        \ar[d, equal] A & \ar[l, loose, "{j(1, f)}"', ""{anchor=center,name=t}] Z \ar[d, equal] & \ar[l, loose, "p_1"'] \cdots & \ar[l, loose, "p_m"'] Y \ar[d, equal]\\
        \ar[d, equal] A & \ar[l, loose, "{j(1, f)}"', ""{anchor=center,name=m}] Z & {} & \ar[ll, loose, "{E(f, \colim^{\vec p}f)}"'] Y \ar[d, equal]\\
        A &&& \ar[lll, loose, "{j(1, \colim^{\vec p}f)}", ""{anchor=center,name=b}] Y
        \ar[from=t, to=m, phantom, "="{anchor=center}] \ar[from=1-3, to=2-3, phantom, "\lambda"{anchor=center}]
        \ar[from=2-1, to=2-4, phantom, ""{name=1}] \ar[from=1, to=b, phantom, "{j\odot\smallsmile_f(1, \colim^{\vec p}f)}"{anchor=center}]
    \end{tikzcd}\]
    For $j$ an arrow, a colimit is \textbf{$j$-absolute} if it is $E(j, 1)$-absolute. It is \textbf{absolute} if it is $\id_E$-absolute.
\end{definition}

The notion of $j$-absolute colimits is usually expressed as preservation by any functor. We get a similar property relative to $j$ and get the usual characterization as a special case.

\begin{lemma}
    \label{lem:rift_absolute}
    Every right lift through $j : A \dfrom E$ respects $j$-absolute colimits.
\end{lemma}

\begin{proof}
    Let $\colim^{\vec p}f$ a $j$-absolute colimit, we have the following isomorphisms:
    \begin{align*}
        (q\rift j)(\colim^{\vec p}f, 1) &\simeqq q\rift j(1, \colim^{\vec p}f) & (\ref{lem:ext_lift:res})\\
                                        &\simeqq q\rift (j(1, f)\odot_L\vec p) & (j-\text{absoluteness})\\
                                        &\simeqq (q\rift (j(1, f))\rift\vec p & (\ref{lem:ext_lift:stab_wcomp})\\
                                        &\simeqq (q\rift j)(f, 1)\rift \vec p & (\ref{lem:ext_lift:res})
    \end{align*}
\end{proof}

\begin{corollary}
    In particular for $j$ an arrow, left pointwise extension along $j$ preserves $j$-absolute colimits. Taking $j=\id_E$, any functor preserves $\id_E$-absolute colimits.
\end{corollary}

\begin{lemma}
    \label{lem:lifting_abs_iff_lift_respect}
    Let $j : A \dfrom E$ be a lifting distributor. A colimit in E is $j$-absolute if and only if every right lift through $j$ respects it.
\end{lemma}

\begin{proof}
    One direction is the Lemma \ref{lem:rift_absolute}. For the other, let %FIXME: finish/repair proof
\end{proof}

<++>

Finally, we define extensions and lift of vertical arrows, that will act like pointwise Kan extension and lifts. It is possible to define them as weighted colimit (as it is done in \cite[Chapter 4]{kelly} or \cite[Remark 4.1.13]{coend_calculus}), but unfolding the definition, it is possible to define four notions directly using right extensions and lifts.

\begin{definition}
    \index{pointiwse lift!left-}\index{pointwise lift!right-}\index{pointwise extension!right-}\index{pointwise extension!left-}
    \label{def:pointwise}
    Let $f : Z \to Y, j : Z \to X$ be vertical arrows.
    \begin{itemize}
        \item The \textbf{left pointwise extension of $f$ along $j$} is, if it exists, the vertical arrow $j\plext f : X \to Y$ such that $Y(j \plext f, 1) \simeq Y(f, 1) \rift X(j, 1)$.
        \item The \textbf{right pointwise extension of $f$ along $j$} is, if it exists, the vertical arrow $j \prext f : X' \to Y$ such that $Y(1, j\prext f) \simeq X(1, j) \rext Y(1, f)$.
    \end{itemize}
    Soit $j : X \to Z, f : Y \to Z$.
    \begin{itemize}
        \item The \textbf{left pointwise lift of $f$ through $j$} is, if it exists, the vertical arrow $f \plift j : X \to Y$ such that $Y(f\plift j, 1) \simeq Z(f, 1) \rext Z(j, 1)$ which is $Z(j, f)$ by Lemma \ref{lem:ext_lift:res}.
        \item The \textbf{right pointwise lift of $f$ through $j$} is, if it exists, the vertical arrow $f \prift j : X \to Y$ such that $Y(1, f\prift j) \simeq Z(1, j) \rift Z(1, f)$ which is $Z(f, j)$ by Lemma \ref{lem:ext_lift:res}.
    \end{itemize}
\end{definition}

In particular, the left pointwise extension of $f$ through $j$ is the $X(j, 1)$-colimit of $f$ and the right pointwise extension of $f$ through $j$ is the $X(1, j)$-limit of $f$. We can then deduce the following properties:

\begin{lemma}
    \label{lem:pt_ext_props}
    Let $f : Z \to Y, j : Z \to X$ be vertical arrows such that $j \plext f : X \to Y$ exists.
    \begin{enumerate}
        \item $1_Z \plext f$ exists and is isomorph to $f$.
        \item Let $j' : X \to X'$ be another vertical arrow. If either side exists, then:
            $$(j'\circ j)\plext f \simeq j' \plext (j\plext f)$$
        \item For all $x : X' \to X, y : Y' \to Y$
            $$Y((j\plext f)x, y) \simeq Y(f, y) \rift X(j, x)$$
    \end{enumerate}
\end{lemma}

\begin{proof}
    The first item is a consequence of the fact that right lift along identities are identities. The second comes from Lemma \ref{lem:co_lim:comp} and the fact that $Y(j', 1) \odot Y(j, 1) \simeqq Y(j'j, 1)$. The last one comes from Lemma \ref{lem:co_lim:res}.
\end{proof}

As a consequence, pointwise extensions satisfy the 2-categorical property of usual Kan extensions.

\begin{corollary}
    \label{cor:kan_2cat}
    Let $j : Z \to X, f : Z \to Y$ such that $j\plext f : X \to Y$ exists. Then we have a bijection of cells:
    \[\begin{iffseq}
        \vec{r}\To X((j\plext f)y, x)\\
        Y(j, y), \vec{r}\To X(f, x)
    \end{iffseq}\]
    And in particular, in $\bb{X}$:
    \[\begin{iffseq}
        j\plext f\To x\\
        f\To x\circ j
    \end{iffseq}\]
\end{corollary}

\begin{proof}
    The first bijection comes from the universal property of the right lift, while the second one follow by taking $y = \id_Y$ and $n=0$.
\end{proof}

\subsection{Density and full faithfulness}

\begin{definition}
    \index{arrow!fully faithful}
    An arrow $j : A \to E$ is \textbf{fully faithful} if it satisfies one of the following equivalent characterizations:
    \begin{enumerate}
        \item $e_f : f \To f$ is cartesian.
        \item $\smallfrown_f : \To E(j, j)$ is opcartesian.
        \item The cell $A(1, 1) \To E(j, j)$ induced by $\smallfrown_j$ is invertible.
    \end{enumerate}
\end{definition}

Having \textit{any} isomorphism suffices:

\begin{corollary}
    $j : A \to E$ is fully faithful if and only if $A(1, 1) \simeqq E(j, j)$.
\end{corollary}

\begin{proof}
    The proof uses horizontal adjunctions, introduced in Subsection \ref{subsec:hadj}. Postcomposition by $\smallfrown_j : \To E(j, j)$ is the unit of $E(1, j) \dashv E(j, 1)$, and by \ref{lem:hadj_opcart} it is opcartesian if and only if $A(1, 1) \simeqq E(j, j)$.
\end{proof}

\begin{lemma}
    \label{lem:ff_refl_co_lim}
    Fully faithful arrows reflect weighted limits and colimits.
\end{lemma}

\begin{proof}
    Let $\vec p : Z \dfrom \cdots \dfrom Y$, $f : Z \to W$ and $g : W \hookrightarrow X$ such that there exists a $\vec p$-cocone $(c : Y \to X, \lambda)$ for $f$ with $(gc, g\lambda)$ exhibiting a colimit $\colim^{\vec p} gf$. Then $W(c, 1) \simeqq X(gc, g) \simeqq X(gf, g)\rext \vec p \simeqq X(f, 1) \rext \vec p$. The dual is similar.
\end{proof}

\begin{lemma}
    \label{lem:ff_res}
    Let $j : A \to E$ be a fully faithful arrow. Then for every $p : A \dfrom X$, the right lift $p \rift E(j, 1)$ exists and the associated cell $(p \rift E(j, 1))(j, 1) \To p$ is an isomorphism.
\end{lemma}

\begin{proof}
    The cell is equal to the composite
    $$p \simeqq p\rift A(1, 1) \simeqq p\rift E(j, j) \simeqq (p\rift E(j, 1))(j, 1)$$
\end{proof}

\begin{lemma}
    \label{lem:ff_lext_comp}
    Let $j : W \to Y, f : W \to X, g : W \to Z$ be arrows such that $j$ is fully faithful, $f\plext j$ and $j\plext g$ exist and the latter preserves the former. Then
    $$(j\plext g)\circ(f\plext j)\simeqq f\plext g$$
\end{lemma}

\begin{proof}
    Using \ref{lem:ff_res}, we have
    $$(j\plext g) \circ (f\plext j) \simeqq f\plext((j\plext g)\circ j)\simeqq f\plext g$$
\end{proof}

\begin{corollary}
    For $i : A \hookrightarrow K$ fully faithful and all $f : A \to B$ such that $i \plext f$ exists, $(i\plext f) \circ i \simeqq f$.
\end{corollary}

\begin{definition}
    \label{def:dense}
    \index{arrow!dense}\index{arrow!codense}
    \index{distributor!dense}\index{distributor!codense}
    A distributor $p : A \dfrom B$ is \textbf{dense} when $e_p : p \To p$ exhibits $B(1, 1) : B \dfrom B$ as the right lift $p\rift p$. An arrow is \textbf{dense} if its conjoint is.

    Dually, it is \textbf{codense} when $e_p : p \To p$ exhibits $A(1, 1)$ as the right extension $p\rext p$. An arrow is \textbf{codense} if its companion is.
\end{definition}

\begin{remark}
    \label{rem:vert_co_dense}
    The condition on conjoint and companion can be rewritten as $j$ being dense if and only if $j\plext j\simeqq \id_E$, exhibited by $e_j : j \To j$. Dually, it is codense if and only if $j\prext j \simeqq \id_A$, exhibited by $e_j$ once again.
\end{remark}

We think of a dense arrow as "essentially describing the codomain", the same way a dense continuous function "describes" the target space. In enriched category theory, this has many equivalent interpretations, notably that every object in the domain is a canonical weighted colimit over the dense functor. In a virtual equipment, we recover a characterization in terms of distributors, generalizing \cite[Lemma 3.20]{arkor2024formaltheoryrelativemonads}:

\begin{lemma}
    \label{lem:res_dense_fid}
    A distributor $p : A \dfrom E$ is dense if and only if the assignment of $\phi : \vec{r} \To E(g, h)$ to
    \[\begin{tikzcd}[sep=huge]
        \ar[d, equal] A & \ar[l, loose, "{p(1, g)}"', ""{anchor=center,name=t}] \cdot \ar[d, equal] & \ar[l, loose, "r_1"'] \cdots & \ar[l, loose, "r_n"'] \cdot \ar[d, equal]\\
        \ar[d, equal] A & \ar[l, loose, "{p(1, g)}", ""{anchor=center,name=m1}] \cdot & {} & \ar[ll, loose, "{E(g, h)}", ""{anchor=center,name=m2}] \cdot \ar[d, equal]\\
        A &&& \ar[lll, loose, "{p(1, h)}", ""{anchor=center,name=b}] \cdot
        \ar[from=2-3, to=2-2, phantom, ""{name=1}]
        \ar[from=t,to=m1, phantom, "="{anchor=center}] \ar[from=1-3, to=m2, phantom, "\phi"{anchor=center}] \ar[from=1, to=b, phantom, "{p\odot\smallsmile_g(1, h)}"{anchor=center}]
    \end{tikzcd}\]
    Induces a bijection of cells:
    \[\begin{iffseq}
        \vec{r}\To E(g, h)\\
        p(1, g), \vec{r}\To p(1, h)
    \end{iffseq}\]
    In particular, there is a bijection of cells:
    \[\begin{iffseq}
        g\To h\\
        p(1, g)\To p(1, h)
    \end{iffseq}\]
\end{lemma}

\begin{proof}
    The first bijection exhibits $E(g, h)$ as the right lift $p(1, h) \rift p(1, g)$, which Corollary \ref{cor:kan_2cat} gives when $j$ is dense, using the characterization in Remark \ref{rem:vert_co_dense}. Taking $g=h=\id_E$, we get $E(1, 1)\simeqq p\rift p$ hence the converse. Taking $n=0$, we get the second bijection.
\end{proof}

\begin{remark}
    In particular, when $p$ is dense, restriction of $p$ on the domain are fully faithful.
\end{remark}

Although dense arrows are not stable by composition, as long as the second arrow preserves certain weighted colimits, it is possible:

\begin{lemma}[{\cite[Lemma 7.9]{averyleinster}}]
    \label{lem:comp_dense}
    Let $A \xrightarrow{f} B \xrightarrow{g} C$ be dense arrows, and such that $g$ preserves $B(f, 1)$-weighted colimits. Then $gf$ is dense.
\end{lemma}

\begin{proof}
    $gf \plext gf \simeq g\plext(f\plext gf) \simeq g\plext(g\circ(f\plext f)) \simeq g\plext g \simeq 1_C$
\end{proof}

\begin{lemma}
    \label{lem:codense_lim_equiv}
    Let $\begin{tikzcd} A \ar[r, "f"] & B \ar[r, "g", bend left] \ar[r, "g'"', bend right] & C\end{tikzcd}$ such that $f$ is codense and $g, g'$ preserve $B(1, f)$-weighted limits. If $gf \simeq g'f$ then $g\simeq g'$.
\end{lemma}

\begin{proof}
    \begin{align*}
        g = g\circ 1_B &\simeq g(f\prext f)\\
                       &\simeq f \prext gf\\
                       &\simeq f \prext g'f\\
                       &\simeq g'(f\prext f)\\
                       &\simeq g'
    \end{align*}
\end{proof}

\begin{lemma}
    \label{lem:ff_comp_dense}
    Let $A \xrightarrow{f} B \xrightarrow{g} C$ with $g$ fully faithful and $gf$ dense. Then $g$ and $f$ are dense.
\end{lemma}

\begin{proof}
    \begin{align*}
        B(1, 1) &\simeq C(g, g) & (g\text{ f.f.})\\
                &\simeq (C(gf, 1) \rift C(gf, 1))(g, g) & (gf\text{ dense})\\
                &\simeq C(gf, g)\rift C(gf, g) & (\ref{lem:ext_lift:res})\\
                &\simeq B(f, 1)\rift B(f, 1) & (g\text{ f.f.})\\
        C(1, 1) &\simeq C(gf, 1)\rift C(gf, 1) & (gf\text{ dense})\\
                &\simeq C(gf, 1)\rift (B(f, 1)\odot C(g, 1)) & (\ref{th:comp_rest})\\
                &\simeq C(gf, 1)\rift (C(gf, g)\odot C(g, 1)) & (g\text{ f.f.})\\
                &\simeq (C(gf, 1)\rift C(gf, g))\rift C(g, 1) & (\ref{rem:ext_lift:comp})\\
                &\simeq (C(gf, 1)\rift C(gf, 1))(g, 1)\rift C(g, 1) & (\ref{lem:ext_lift:res})\\
                &\simeq C(g, 1)\rift C(g, 1) & (gf\text{  dense})
    \end{align*}
\end{proof}

Density arrows also interact nicely with weighted limits and absolute colimits:

\begin{lemma}
    \label{lem:ff_dense_prelim}
    Let $f$ be a dense and fully faithful vertical arrow. Then $f$ preserves weighted limits.
\end{lemma}

\begin{proof}
    Let $\vec{p}, g$ such that $\lim^{\vec{p}}g$ exists.
    \begin{align*}
        B(1, f\circ \lim{}^{\vec{p}}g) &= B(1, 1)(1, f\circ\lim{}^{\vec{p}}g)\\
                                     &\simeq (B(f, 1)\rift B(f, 1))(1, f\circ\lim{}^{\vec{p}}g) & (\text{density})\\
                                     &\simeq (B(f, f)\rift B(f, 1))(1, \lim{}^{\vec{p}}g) & (\ref{lem:ext_lift:res})\\
                                     &\simeq (A(1, 1)\rift B(f, 1))(1, \lim{}^{\vec{p}}g) & (\text{f.f.})\\
                                     &\simeq A(1, \lim{}^{\vec{p}}g)\rift B(f, 1) & (\ref{lem:ext_lift:res})\\
                                     &\simeq \vec{p}\rext (A(1, g)\rift B(f, 1)) & (\ref{def:lim}, \ref{lem:ext_lift:assoc_rift_rext})\\
                                     &\simeq \vec{p}\rext B(1, gf)
    \end{align*}
\end{proof}

\begin{lemma}
    \label{lem:dense_abs_iff_iso}
    Let $p : Z \dfrom \cdots \dfrom Y$ a weight, $j : A \dfrom E$ a dense distributor and $f : Z \to E$ an arrow. An arrow $c : Y \to E$ forms a $j$-absolute $\vec p$-weighted colimit of $f$ if and only if there is an isomorphism
    $$j(1, c) \simeq j(1, v) \odot_L \vec p$$
\end{lemma}

\begin{proof}
    The only if direction is immediate, so suppose the isomorphism exists. Then we have the following isomorphisms:
    \begin{align*}
        E(c, 1) &\simeqq j \rift j(1, c) & (\text{density})\\
                &\simeqq j \rift (j(1, f)\odot_L\vec p) & (\text{by hypothesis})\\
                &\simeqq (j \rift j(1, f))\odot_L \vec p & (\ref{lem:ext_lift:stab_wcomp}\&\ref{lem:ext_lift:stab_split})\\
                &\simeqq (j \rift j)(f, 1)\rift \vec p & (\ref{lem:ext_lift:res})\\
                &\simeqq E(f, 1)\rift \vec p & (\text{density})
    \end{align*}
    %FIXME: not finished ([AM26] 2.30)
\end{proof}

Absoluteness also characterises full faithfulness of weighted colimits:

\begin{lemma}
    Let $p : Y \dfrom Z, f : Y \to X$ such that $\colim^pf$ exists. Any three of the following conditions implies the fourth
\end{lemma}

\begin{proof}
    We have a chain of isomorphism, that can be translated to obtain every implication. The following is $(a)\land(b)\land(c)\implies(d)$:
    \begin{align*}
        Z(1, 1) &\simeqq p\rift p & (p\text{ dense})\\
                &\simeqq X(f, f)\odot_L p)\rift p & (f\text{ fully faithful})\\
                &\simeqq X(f, \colim^pf)\rift p & (\colim^pf\text{ is $f$-absolute})\\
                &\simeqq (X(f, 1)\rift p)(1, \colim^pf) & \text{(Lemma \ref{lem:ext_lift:res})}\\
                &\simeqq X(\colim^pf, \colim^pf)
    \end{align*}
\end{proof}

Finally, density and full faithfulness can characterize when left extensions witness equivalences.

\begin{lemma}
    \label{lem:left_witness_equiv}
    Let $j : A \to E, f: A \to B$ arrows such that $f\plext j$ and $j\plext f$ exist and $f\plext j \dashv j\plext f$. Then
    \begin{enumerate}
        \item The adjunction is coreflective if $j$ is fully faithful, $f$ is dense and $f\plext j$ is $j$-absolute.
        \item The adjunction is reflective if $f$ is fully faithful and $j$ is dense.
    \end{enumerate}
    In particular, the adjunction is an equivalence when $f$ and $j$ are dense and fully faithful and $f\plext j$ is $j$-absolute.
\end{lemma}

\begin{proof}
    $f\plext j$ is preserved by $j\plext f$ since it is $j$-absolute (Lemma \ref{lem:rift_absolute}) and by Lemma \ref{lem:ff_lext_comp} we have $(j\plext f)\circ(f\plext j) \simeqq f\plext f \simeqq 1$. Conversely, $f\plext j$ preserves all left extensions, being a left adjoint, so Lemma \ref{lem:ff_lext_comp} gives $(f\plext j)\circ(j\plext f) \simeqq j\plext j\simeqq 1$.
\end{proof}

\section{Adjunctions and Monads}

Not only do we have two kinds of arrows with cells between them in a virtual equipment, leading to two possible notions of adjunctions; there's also distributors, permitting the definition of adjunctions through the isomorphism on $\Hom$ distributors. We start with horizontal adjunctions and monads, as it leads to a few useful results in the study of (vertical) relative adjunctions and monads.

\subsection{Loose adjunctions and loose monads}
\label{subsec:hadj}

The definition and many results in this section make sense in an arbitrary virtual double category, as written in \cite[Subection 2.4]{arkor2024formaltheoryrelativemonads}.

\begin{definition}
    \index{adjunction!horizontal}
    \begin{wrapfigure}{r}{0.47\textwidth}
        \vspace{-\baselineskip}
        \[\begin{tikzcd}[row sep=small]
            C \ar[d, equal] & \ar[l, loose, "\ell"', ""{anchor=center,name=t1}] A \ar[d, equal] & \ar[l, equal, ""{anchor=center,name=t2}] A \ar[d, equal]\\
            C \ar[d, equal] & \ar[l, loose, "\ell", ""{anchor=center,name=b1}] A & \ar[l, loose, "r\odot\ell", ""{anchor=center,name=b2}] A \ar[d, equal]\\
            C && \ar[ll, loose, "\ell", ""{anchor=center,name=bb}] A
            \ar[from=t1, to=b1, phantom, "="{anchor=center}] \ar[from=t2, to=b2, phantom, "\eta"{anchor=center}, yshift=-0.4ex]
            \ar[from=2-2, to=bb, phantom, "{(1_\ell, \varepsilon)}"{anchor=center}]
        \end{tikzcd}
        =
        \begin{tikzcd}
            C \ar[d, equal] & \ar[l, loose, "\ell"', ""{anchor=center,name=t}] A \ar[d, equal]\\
            C & \ar[l, loose, "\ell", ""{anchor=center,name=b}] A
            \ar[from=t, to=b, phantom, "="{anchor=center}]
        \end{tikzcd}\]
        \[\begin{tikzcd}[row sep=small]
            A \ar[d, equal] & \ar[l, equal, ""{anchor=center,name=t1}] A \ar[d, equal] & \ar[l, loose, "r"', ""{anchor=center,name=t2}] C \ar[d, equal]\\
            A \ar[d, equal] & \ar[l, loose, "r\odot\ell", ""{anchor=center,name=b1}] A & \ar[l, loose, "r", ""{anchor=center,name=b2}] C \ar[d, equal]\\
            A && \ar[ll, loose, "r", ""{anchor=center,name=bb}] C
            \ar[from=t1, to=b1, phantom, "\eta"{anchor=center}, yshift=-0.4ex] \ar[from=t2, to=b2, phantom, "="{anchor=center}]
            \ar[from=2-2, to=bb, phantom, "{(\varepsilon, 1_r)}"{anchor=center}]
        \end{tikzcd}
        =
        \begin{tikzcd}
            A \ar[d, equal] & \ar[l, loose, "r"', ""{anchor=center,name=t}] C \ar[d, equal]\\
            A & \ar[l, loose, "r", ""{anchor=center,name=b}] C
            \ar[from=t, to=b, phantom, "="{anchor=center}]
        \end{tikzcd}\]
    \end{wrapfigure}
    A \textbf{horizontal adjunction} $\ell\dashv r$ comprises of the following data:
    \begin{enumerate}
        \item An object $A$, the \textbf{base}.
        \item An object $C$, the \textbf{apex}.
        \item A horizontal arrow $\ell : C \dfrom A$, the \textbf{left relative adjoint}.
        \item A horizontal arrow $r : A \dfrom C$ \textbf{right relative adjoint}.
        \item A cell $\eta : \To r\odot\ell$, the \textbf{unit}.
        \item A cell $\varepsilon : \ell, r \To e_C$, the \textbf{counit}.
    \end{enumerate}
    such that $r\odot\ell$ exists and we have the equalities on the right.
\end{definition}

\begin{example}
    Let $\ell : A \to C$ be a vertical arrow. Noting that $C(\ell, \ell) \simeq C(\ell, 1) \odot C(1, \ell)$ (Theorem \ref{th:comp_rest}), we have $C(1, \ell) \dashv C(\ell, 1)$ with $\smallfrown_\ell$ as unit and $\smallsmile_\ell$ as counit.
\end{example}

\begin{definition}
    \label{monad!loose}
    A \textbf{loose monad} in $\bb{X}$ comprises of the following data:
    \begin{enumerate}
        \item A \textbf{base} object $A$.
        \item A \textbf{underlying distributor} $t: A \dfrom A$.
        \item A \textbf{multiplication} $\mu: t, t \To t$
        \item A \textbf{unit} $\eta: \To t$
    \end{enumerate}
    Such that
    \begin{align*}
        \mu(\mu, e_t) &= \mu(e_t, \mu) & \mu(\eta, e_t) &= \mu(e_t, \eta) = e_t
    \end{align*}
    A \textbf{loose monad morphism} from $(t, \mu, \eta)$ to $(t', \mu', \eta')$ both with base $A$ is a cell $\tau : t \To t'$ such that
    \begin{align*}
        \tau\eta &= \eta' & \tau\mu &= \mu'(\tau, \tau)
    \end{align*}
\end{definition}

We will come back to loose monads in Subsection \ref{subsec:internal}.

Thankfully, loose adjunctions induce loose monads:

\begin{lemma}
    \label{lem:hadj_hmon}
    Every loose adjunction $\ell \dashv r$ induces a loose monad.
\end{lemma}

\begin{proof}
    Let $\ell \dashv r$ be an adjunction with unit $\eta$ and counit $\varepsilon$. The loose monad has $r\odot\ell$ as underlying distributor, $r\odot\varepsilon\odot\ell$ as multiplication and $\eta$ as unit. Associativity comes from associativity in composition of cells and unit laws from diagram chasing.
\end{proof}

Finally, a useful lemma characterizing certain adjunctions:

\begin{lemma}
    \label{lem:hadj_opcart}
    Let $\ell \dashv r$ be a horizontal adjunction. Then $r\odot\ell\simeqq A(1, 1)$ if and only if $\eta:\To r\odot\ell$ is opcartesian.
\end{lemma}

\begin{proof}
    One direction is trivial, we focus on the other: suppose $r\odot\ell\simeqq A(1, 1)$. The induced loose monad structure on A(1, 1) (by Lemma \ref{lem:hadj_hmon}). By Eckmann-Hilton, $\bb{X}\iset{A, A}(A(1, 1), A(1, 1))$ is equipped with a commutative monoidal structure, so the unit is a two-sided inverse. Transferring the structure back to $r\odot\ell$, the unit is invertible.
\end{proof}

\subsection{Relative adjunctions}

Intuitively, we have two possible definitions for adjunctions: either we use an isomorphism on the (restriction of) distributors, thinking of it as an isomorphism on the $\Hom$, or we use the unit/counit characterisation. The first is easily defined in an arbitrary virtual equipment, while the second can be expressed in the underlying 2-category $\underline{\bb{X}}$ (\cite[p.~762]{maranda} for a historic reference, \cite[Paragraph~1.1]{auderset_2-cat_adj} and \cite[Section 1.11]{kelly} for a modern treatment).

We will see later (Corollary \ref{cor:adj_defs_same}) that the two notions coincide. We delay the definition of an adjunction, which we will get as a special case of a relative adjunction.

% TODO: texte de transition

\begin{definition}
    \index{adjunction!relative}\index{root}\index{adjoint!right-}\index{adjoint!left-}
    \begin{wrapfigure}{r}{0.5\textwidth}
        \centering
        \begin{tikzcd}
            & C \ar[dr, "r", ""{anchor=north,name=radj}] &\\
            A \ar[ur, "\ell", ""{anchor=north,name=ladj}] \ar[rr, "j"'] && E
            \ar[from=ladj, to=radj, phantom, "\dashv"{anchor=center}]
        \end{tikzcd}
    \end{wrapfigure}
    A \textbf{relative adjunction} $\ell \dashv_j r$ comprises of the following data:
    \begin{itemize}
        \item A \textbf{root} $j : A \to E$.
        \item A \textbf{(relative) left adjoint} $\ell : A \to C$.
        \item A \textbf{(relatif) right adjoint} $r : C \to E$.
        \item An isomorphism $\sharp : C(\ell, 1) \simeq E(j, r) : \flat$. %FIXME: not great looking with Libertinus :c
    \end{itemize}
\end{definition}

\begin{remark}
    We also speak of \textit{$j$-relative adjunction}.
\end{remark}

The root $j$ breaks the symmetry of the adjunction. Dualizing, we get a different notion, that of relative coadjunction.

\begin{definition}
    \index{coadjunction!relative}\index{coroot}\index{coadjoint!right-}\index{coadjoint!left-}
    \begin{wrapfigure}{r}{0.5\textwidth}
        \centering
        \begin{tikzcd}
            Z \ar[dr, "r"', ""{anchor=south,name=radj}] \ar[rr, "i"] && V\\
            & X \ar[ur, "\ell"', ""{anchor=south,name=ladj}] &\\
            \ar[from=radj, to=ladj, phantom, "\vdash"{anchor=center}]
        \end{tikzcd}
    \end{wrapfigure}
    A \textbf{relative coadjunction} $r \vdash^i \ell$ comprises of the following data:
    \begin{itemize}
        \item A \textbf{coroot} $i : Z \to V$.
        \item A \textbf{relative left-coadjoint} $\ell : X \to V$.
        \item A \textbf{relative right-coadjoint} $r : Z \to X$.
        \item An isomorphism $\sharp : V(\ell, i) \simeq X(1, r) : \flat$. %FIXME: not great looking with Libertinus :c
    \end{itemize}
\end{definition}

It is possible to dualise results for relative coadjunctions. Note also that the distinction disappears when the (co)root is the identity, meaning $\ell\dashv_1 r$ if and only if $r\vdash^1\ell$.

As for classical adjunctions, there's multiple characterizations using unit and counit cells.

\begin{proposition}
    \label{prop:equiv_radj_car}
    Let $j : A \to E, \ell : A \to C, r : C \to E$ be vertical arrows. The following data are equivalent to characterize an $j$-relative adjunction $\ell \dashv_j r$:
    \begin{enumerate}[ref=\ref{prop:equiv_radj_car} (\alph*)]
        \item An isomorphism $\sharp : C(\ell, 1) \simeq E(j, r) : \flat$.
        \item An unit cell $\eta : j \To r\circ\ell$ and counit $\varepsilon : C(1, \ell), E(j, r) \To C(1 ,1)$ such that
            \[\begin{tikzcd}
                \ar[d, equal] C & \ar[l, loose, "{C(1, \ell)}"', ""{anchor=center,name=t1}] A \ar[d, equal] & \ar[l, equal, ""{anchor=center,name=t2}] A \ar[d, equal] \\
                \ar[d, equal] C & \ar[l, loose, "{C(1, \ell)}"', ""{anchor=center,name=m1}] A \ar[d, equal] & \ar[l, loose, "{E(j, j)}"', ""{anchor=center,name=m2}] A \ar[d, equal]\\
                \ar[d, equal] C & \ar[l, loose, "{C(1, \ell)}", ""{anchor=center,name=b1}] A & \ar[l, loose, "{E(j, r\ell)}", ""{anchor=center,name=b2}] A \ar[d, equal]\\
                C && \ar[ll, loose, "{C(1, \ell)}", ""{anchor=center,name=bb}] A
                \ar[from=t1, to=m1, phantom, "="{anchor=center}] \ar[from=t2, to=m2, phantom, "\smallfrown_j"{anchor=center}]
                \ar[from=m1, to=b1, phantom, "="{anchor=center}] \ar[from=m2, to=b2, phantom, "{E(j, \eta)}"{anchor=center}]
                \ar[from=3-2, to=bb, phantom, "{\varepsilon(1, \ell)}"{anchor=center}]
            \end{tikzcd}
            \hspace{1em} = \hspace{1em}
            \begin{tikzcd}
                \ar[d, equal] C & \ar[l, loose, "{C(1, \ell)}"', ""{anchor=center,name=t}] A \ar[d, equal]\\
                C & \ar[l, loose, "{C(1, \ell)}", ""{anchor=center,name=b}] A
                \ar[from=t, to=b, phantom, "="{anchor=center}]
            \end{tikzcd}\]
            \[\begin{tikzcd}
                \ar[d, equal] E & \ar[l, loose, "{E(1, j)}"', ""{anchor=center,name=t1}] A \ar[d, equal] & \ar[l, loose, "{E(j, 1)}"', ""{anchor=center,name=t2}] E \ar[d, equal] & \ar[l, loose, "{E(1, r)}"', ""{anchor=center,name=t3}] C \ar[d, equal]\\
                \ar[d, equal] E & \ar[l, loose, "{E(1, r\ell)}", ""{anchor=center,name=b1}] A & \ar[l, loose, "{E(j, 1)}", ""{anchor=center,name=b2}] E & \ar[l, loose, "{E(1, r)}", ""{anchor=center,name=b3}] C \ar[d, equal]\\
                E &&& \ar[lll, loose, "{E(1, r)}", ""{anchor=center,name=bb}] C
                \ar[from=t1, to=b1, phantom, "{E(1, \eta)}"{anchor=center}] \ar[from=t2, to=b2, phantom, "="{anchor=center}] \ar[from=t3, to=b3, phantom, "="{anchor=center}]
                \ar[from=b2, to=bb, phantom, "{E(1, r), \varepsilon}"{anchor=center}, yshift=-0.8ex]
            \end{tikzcd}
            \hspace{1em} = \hspace{1em}
            \begin{tikzcd}
                \ar[d, equal] E & \ar[l, loose, "{E(1, j)}"'] A & \ar[l, loose, "{E(j, 1)}"', ""{anchor=center,name=t}] E & \ar[l, loose, "{E(1, r)}"'] C \ar[d, equal]\\
                E &&& \ar[lll, loose, "{E(1, r)}", ""{anchor=center,name=b}] C
                \ar[from=t, to=b, phantom, "{\smallsmile_j(1, r)}"{anchor=center}]
            \end{tikzcd}\]
        \item A unit cell $\eta : j \To r\circ\ell$ and a cell $\flat : E(j, r) \To C(\ell, 1)$ such that the following diagrams commute
            \[\begin{tikzcd}
                \ar[d, "\smallfrown_j"'] A(1, 1) \ar[r, "\smallfrown_j"] & E(j, j) \ar[d, "{E(j, \eta)}"]\\
                C(\ell, \ell) & \ar[l, "{\flat(1, \ell)}"] E(j, r\ell)
            \end{tikzcd}
            \hspace{10pt}
            \begin{tikzcd}
                \ar[d, equal] E(j, r) \ar[r, "\flat"] & C(\ell, 1) \ar[d, "{\smallfrown_r(\ell, 1)}"]\\
                E(j, r) & \ar[l, "{E(\eta, r)}"] E(r\ell, r)
            \end{tikzcd}\]
        \item A cell $\sharp : C(\ell, 1) \To E(j, r)$ and a counit cell $\varepsilon : C(1, \ell), E(j, r) \To C(1, 1)$ such that
            \[\begin{tikzcd}
                \ar[d, equal] C & \ar[l, loose, "{C(1, \ell)}"', ""{anchor=center,name=t1}] A \ar[d, equal] & \ar[l, loose, "{C(\ell, 1)}"', ""{anchor=center,name=t2}] C \ar[d, equal]\\
                \ar[d, equal] C & \ar[l, loose, "{C(1, \ell)}", ""{anchor=center,name=b1}] A & \ar[l, loose, "{E(j, r)}", ""{anchor=center,name=b2}] C \ar[d, equal]\\
                C && \ar[ll, equal, loose, ""{anchor=center,name=bb}] C
                \ar[from=t1, to=b1, phantom, "="{anchor=center}] \ar[from=t2, to=b2, phantom, "\sharp"{anchor=center}]
                \ar[from=2-2, to=bb, phantom, "\varepsilon"{anchor=center}]
            \end{tikzcd}
            \hspace{1em} = \hspace{1em}
            \begin{tikzcd}
                \ar[d, equal] C & \ar[l, loose, "{C(1, \ell)}"'] A & \ar[l, loose, "{C(\ell, 1)}"'] C \ar[d, equal]\\
                C && \ar[ll, equal, loose, ""{anchor=center,name=b}] C
                \ar[from=1-2, to=b, phantom, "\smallsmile_\ell"{anchor=center}]
            \end{tikzcd}
            \hspace{3em}
            \begin{tikzcd}
                \ar[d, equal] A & \ar[l, equal, ""{anchor=center,name=t1}] \ar[d, equal] & \ar[l, loose, "{E(j, r)}"', ""{anchor=center,name=t2}] C \ar[d, equal]\\
                \ar[d, equal] A & \ar[l, loose, "{C(\ell, \ell)}", ""{anchor=center,name=b1}] A & \ar[l, loose, "{E(j, r)}", ""{anchor=center,name=b2}] C \ar[d, equal]\\
                A && \ar[ll, loose, "{E(j, r)}", ""{anchor=center,name=bb}] C
                \ar[from=t1, to=b1, phantom, "\smallfrown_\ell"{anchor=center}] \ar[from=t2, to=b2, phantom, "="{anchor=center}]
                \ar[from=2-2, to=bb, phantom, "{\sharp, \varepsilon}"{anchor=center}]
            \end{tikzcd}
            \hspace{1em} = \hspace{1em}
            \begin{tikzcd}
                \ar[d, equal] A & \ar[l, loose, "{E(j, r)}"', ""{anchor=center,name=t}] C \ar[d, equal]\\
                A & \ar[l, loose, "{E(j, r)}", ""{anchor=center,name=b}] C
                \ar[from=t, to=b, phantom, "="{anchor=center}]
            \end{tikzcd}\]
        \item A horizontal adjunction $C(1, \ell) \dashv E(j, r)$
    \end{enumerate}
\end{proposition}

\begin{proof}
    We just define $\eta$ from $\sharp$, $\varepsilon$ from $\flat$ and reciprocally.
    \[\eta :=
        \begin{tikzcd}
            \ar[d, equal] A & \ar[l, equal] A \ar[d, "\ell"'] \\
            \ar[d, equal] A & \ar[l, loose, "{C(\ell, 1)}"', ""{anchor=center,name=t}] C \ar[d, equal] \\
            \ar[d, "j"r] & \ar[l, loose, "{E(j, r)}", ""{anchor=center,name=b}] C \ar[d, "r"]\\
            E & \ar[l, loose, equal] E
            \ar[from=t, to=b, phantom, "\sharp"{anchor=center}]
        \end{tikzcd}
        \hspace{10pt}
        \varepsilon :=
        \begin{tikzcd}
            \ar[d, equal] C & \ar[l, loose, "{C(1, \ell)}"', ""{anchor=center,name=t1}] A \ar[d, equal] & \ar[l, loose, "{E(j, r)}"', ""{anchor=center,name=t2}] C \ar[d, equal] \\
            \ar[d, equal] C & \ar[l, loose, "{C(1, \ell)}", ""{anchor=center,name=b1}] A & \ar[l, loose, "{C(\ell, 1)}", ""{anchor=center,name=b2}] C \ar[d, equal] \\
            C && \ar[ll, loose, equal, ""{anchor=center,name=bb}] C
            \ar[from=t1, to=b1, phantom, "="{anchor=center}] \ar[from=t2, to=b2, phantom, "\flat"{anchor=center}]
            \ar[from=2-2, to=bb, phantom, "\smallsmile_\ell"{anchor=center}]
        \end{tikzcd}
    \]
    \[\sharp :=
        \begin{tikzcd}
            A \ar[dd, equal] & \ar[l, equal] A \ar[dd, "j"'] & \ar[l, equal, ""{anchor=center,name=t}] \ar[d, "\ell"] & \ar[l, loose, "{C(\ell, 1)}"'] C \ar[d, equal]\\
                             && C \ar[d, "r"] & \ar[l, equal, loose] C \ar[d, equal]\\
            A \ar[d, equal] & \ar[l, loose, "{E(j, 1)}"] E & \ar[l, loose, equal, ""{anchor=center,name=b}] E & \ar[l, loose, "{E(1, r)}"] C \ar[d, equal]\\
            A &&& \ar[lll, "{E(j, r)}"] C
            \ar[from=t, to=b, phantom, "\eta"{anchor=center}]
        \end{tikzcd}
        \hspace{10pt}
        \flat :=
        \begin{tikzcd}
            \ar[d, equal] A & \ar[l, equal] A \ar[d, "\ell"] & \ar[l, equal] A \ar[d, equal] & \ar[l, loose, "{E(j, r)}"', ""{anchor=center,name=t1}] C \ar[d, equal] \\
            \ar[d, equal] A & \ar[l, loose, "{C(\ell, 1)}"', ""{anchor=center,name=m1}] C \ar[d, equal] & \ar[l, loose, "{C(1, \ell)}"'] A & \ar[l, loose, "{E(j, r)}", ""{anchor=center,name=m2}] C \ar[d, equal]\\
            \ar[d, equal] A & \ar[l, loose, "{C(\ell, 1)}", ""{anchor=center,name=b1}] C && \ar[ll, loose, equal, ""{anchor=center,name=b2}] C \ar[d, equal]\\
            A &&& \ar[lll, loose, "{C(\ell, 1)}"] C
            \ar[from=t1, to=m2, phantom, "="{anchor=center}]
            \ar[from=m1, to=b1, phantom, "="{anchor=center}] \ar[from=2-3, to=b2, phantom, "\varepsilon"{anchor=center}]
        \end{tikzcd}
    \]
\end{proof}

\begin{corollary}
    \label{cor:adj_defs_same}
    Let $\ell : A \to C, r : C \to A$ be vertical arrows. Then they are $1_A$-relatively adjoint if and only if they are adjoint in $\underline{\bb{X}}$.
\end{corollary}

Considering relative adjunctions is not without consequences: the root distorts the adjunction, in particular the right adjoint. The usual properties are therefore modified to depend on the root and its properties. In particular, when the root is dense, everything works well.

\begin{lemma}
    If $\ell \dashv_j r$ and $\ell' \dashv_j r$ then $l \simeq l'$. If $j$ is dense, $\ell \dashv_j r$ and $\ell \dashv_j r'$ then $r\simeq r'$.
\end{lemma}

\begin{proof}
    Can be deduced from the isomorphism of horizontal arrows and the lemma \ref{lem:res_dense_fid}.
\end{proof}

We also recover a conservation of (co)limits.

\begin{proposition}
    \label{prop:radj_co_lim}
    Let $j : A \to E$ a root.
    \begin{itemize}
        \item If $\ell : A \to C$ is a left $j$-adjoint, then $\ell$ preserves the colimits that $j$ preserves.
        \item If $j$ is dense and $r : C \to E$ is a right $j$-adjoint, then $r$ preserves every limit.
    \end{itemize}
\end{proposition}

% FIXME: ajouter lemmes justifiant les isos
\begin{proof}
    Let $j: A \to E$ be a root and $\ell \dashv_j r$ a $j$-adjunction.
    \begin{itemize}
        \item Let $\vec{p} : Y \dfrom \cdots \dfrom X$ be a weight and $f : Y \to A$ a vertical arrow such that $\colim^{\vec{p}}f$ exists and is preserved by $j$:
            \begin{align*}
                C(\ell\circ\colim^{\vec{p}}f, 1) &\simeq E(j\circ\colim^{\vec{p}}f, r)\\
                                                 &\simeq E(\colim^{\vec{p}}(j\circ f), r)\\
                                                 &\simeq E(jf, r) \rift \vec{p}\\
                                                 &\simeq C(\ell\circ f, 1) \rift \vec{p}
            \end{align*}
        \item Suppose that $j$ is dense, let $\vec{p} : X \dfrom \cdots \dfrom Y$ be a weight and $f : Y \to C$ a vertical arrow such that $\lim^{\vec{p}}f$ exists:
            \begin{align*}
                E(1, r\circ\lim{}^{\vec{p}}f) &\simeq E(j \plext j, r\circ\lim{}^{\vec{p}}f)\\
                                            &\simeq E(j , r\circ\lim{}^{\vec{p}}f) \rift E(j, 1)\\
                                            &\simeq C(\ell, \lim{}^{\vec{p}}f) \rift E(j, 1)\\
                                            &\simeq (p\rext C(\ell, f)) \rift E(j, 1)\\
                                            &\simeq (p\rext E(j, rf)) \rift E(j, 1)\\
                                            &\simeq p \rext (E(j, rf) \rift E(j, 1))\\
                                            &\simeq p \rext E(j\plext j, rf)\\
                                            &\simeq p \rext E(1, rf)
            \end{align*}
    \end{itemize}
\end{proof}

\begin{corollary}
    Every dense and fully faithful arrow creates weighted limits.
\end{corollary}

\begin{proof}
    Let $f : A \to B$ a dense and fully faithful arrow, therefore it reflects limits by \ref{lem:ff_refl_co_lim}. By full faithfulness, $1\dashv_ff$, hence by density it preserves weighted limits by Proposition \ref{prop:radj_co_lim}.
\end{proof}

Better than that, when $j$ is dense the right adjoint also preserves a class of colimits:

\begin{lemma}
    \label{lem:radj_right_labs_colim}
    Let $j : A \to E$ be a dense arrow, $\ell\dashv_j r$ a relative adjunction. Then $r$ preserves $\ell$-absolute colimits
\end{lemma}

\begin{proof}
    We have
    \begin{align*}
        E(j, r\colim^{\vec p}f) &\simeqq C(\ell, \colim^{\vec p}f)\\
                                &\simeqq C(\ell, f)\odot_L\vec p\\
                                &\simeqq E(j, rf)\odot_L\vec p
    \end{align*}
    Hence $r\colim^{\vec p}f\simeqq\colim^{\vec p}(rf)$ by Lemma \ref{lem:dense_abs_iff_iso}.
\end{proof}

\begin{corollary}
    In particular, left $j$-adjoints preserve the left pointwise extensions that $j$ preserve and right $j$-adjoint preserve every right pointwise extension when $j$ is dense.
\end{corollary}

Relative adjoints, as usual adjoints, can be computed by absolute pointwise lifts and  extensions.

\begin{proposition}
    Let $j : A \to E, r : C \to E$ be arrows. An arrow $\ell : A \to C$ is a left $j$-adjoint to $r$ if and only if there is a cell $\eta : j \to r\ell$ exhibiting $\ell$ as the left lift $j\plift r$
    \[\begin{tikzcd}
        &  C \ar[dr, "r"] &\\
        A \ar[ur, dashed, "j\plift r"] \ar[rr, "j"'] &{}& E
        \ar[from=1-2, to=2-2, phantom, "\dashv"{anchor=center}]
    \end{tikzcd}\]
\end{proposition}

\begin{proof}
    Immediate from Definition \ref{def:pointwise}.
\end{proof}

Once again, we need an extra condition regarding the root when considering the right adjoint.

\begin{proposition}
    Let $j : A \to E$ and $\ell : A \to C$ be arrows with $j$ fully faithful.
    \[\begin{tikzcd}
        &  C \ar[dr, dashed, "\ell\plext r"] &\\
        A \ar[ur, "\ell"] \ar[rr, "j"'] &{}& E
        \ar[from=1-2, to=2-2, phantom, "\dashv"{anchor=center}]
    \end{tikzcd}\]
    \begin{enumerate}
        \item If $\ell\plext j$ exists and is $j$-absolute, $l\dashv_j\ell\dashv j$.
        \item If $j$ is dense and $\ell$ has a right $j$-adjoint $r$, $r$ is the pointwise left extension $\ell\plext j$ and it is $j$-absolute.
    \end{enumerate}
\end{proposition}

\begin{proof}
    For (a), we have by $j$-absoluteness $E(j, l\plext\ell, j)\simeq E(j, j)\odot_L C(\ell, 1)$ and by fullfaithfulness $E(j, j)\odot_L C(\ell, 1) \simeqq A(1, 1)\odot_L C(\ell, 1) \simeqq C(\ell, 1)$ as required. For (b), we use Lemma \ref{lem:dense_abs_iff_iso}.
\end{proof}

There are two notions of morphisms of relative adjunction corresponding to each adjoint. As most of the properties regarding relative adjunctions, the left kind is kinder than the right which requires the right properties. Relative adjoint morphisms will induce morphisms of relative monads in Subsection \ref{subsec:rmonad}.

\begin{definition}
    \index{adjunction!left morphism}
    Let $j : A \to E$ be an arrow, $\ell\dashv_j r$ and $\ell'\dashv_jr'$ $j$-relative adjunctions. A \textbf{left morphism} from $\ell\dashv_jr$ to $\ell'\dashv_j r'$ consists of
    \[\begin{tikzcd}
        &C'\ar[ddr, "r'", bend left]&\\
        &\ar[u, "c"{description}] C \ar[dr, "r"]\\
        A \ar[uur, bend left, "\ell'", ""{anchor=center,name=a}] \ar[ur, "\ell"] \ar[rr, "j"'] && E
        \ar[from=a, to=2-2, Rightarrow, "\lambda"', shorten=2mm]
    \end{tikzcd}\]
    \begin{enumerate}
        \item An arrow $c : C \to C'$ such that $r = r'c$.
        \item A cell $\lambda : \ell'\To \ell$.
    \end{enumerate}
    Rendering the following diagram commutative
    \[\begin{tikzcd}
        \ar[d, "{\smallfrown_c(\ell, 1)}"'] C(\ell, 1) \ar[r, "\sharp"] & E(j, r)\\
        C'(c\ell, c) \ar[r, "{C'(\lambda, c)}"'] & C'(\ell', c)\ar[u, "{\sharp'(1, c)}"']
    \end{tikzcd}\]
    It is \textbf{strict} when $\lambda$ is the identity. We write $\cat{RAdj}_L(j)$ the category of $j$-relative adjunctions and left morphisms.
\end{definition}

\begin{example}
    For any arrow $j : A \to E$ and $j$-relative adjunction $\ell\dashv_j r$, $(r, \eta)$ forms a unique left morphism from $\ell\dashv_j r$ to $j\dashv_j \id_E$. The latter is hence terminal in $\cat{RAdj}_L(j)$.
\end{example}

Thankfully \textit{for left morphisms}, the cell is uniquely determined by the arrow $c$. We write $\underline{\bb X}/E$ for the (1-)category of slices over an object $E$.

\begin{lemma}
    \label{lem:rradj_ff}
    Let $j : A \to E$ be an arrow. The functor $\cat{RAdj}_L(j) \to \underline{\bb X}/E$ sending a $j$-adjunction to its right adjoint and a left morphism to its arrow is fully faithful.
\end{lemma}

\begin{proof}
    Given a morphism $c : C \to C'$ such that $r=r'c$, we define $\lambda$ as the following cell
    \[\begin{tikzcd}
        \ar[d, equal] A & \ar[l, loose, "{C(\ell, 1)}"', ""{anchor=center, name=t1}] \ar[d, equal] C & \ar[l, loose, "{C'(c, 1)}"', ""{anchor=center, name=t2}] C' \ar[dd, equal]\\
        \ar[d, equal] A & \ar[l, loose, "{E(j, r)}"', ""{anchor=center, name=m}] C \ar[d, equal] & \\
        \ar[d, equal] A & \ar[l, loose, "{C'(\ell', c)}"', ""{anchor=center, name=b1}] C & \ar[l, loose, "{C'(c, 1)}"', ""{anchor=center, name=b2}] C' \ar[d, equal]\\
        A &{}& \ar[ll, loose, "{C'(\ell', 1)}"] C'
        \ar[from=t1, to=m, phantom, "\sharp"{anchor=center}] \ar[from=t2, to=b2, phantom, "="{anchor=center}]
        \ar[from=m, to=b1, phantom, "{\flat'(1, c)}"{anchor=center}]\\
        \ar[from=3-2, to=4-2, phantom, "{\smallsmile_c(\ell', 1)}"{anchor=center}]
    \end{tikzcd}\]
    This cell is induced by the compatibility conditions and is therefore unique.
\end{proof}

We now shift to right morphisms.

\begin{definition}
    \index{adjunction!right morphism}
    Let $j : A \to E$ be an arrow, $\ell\dashv_j r$ and $\ell'\dashv_jr'$ $j$-relative adjunctions. A \textbf{right morphism} from $\ell\dashv_jr$ to $\ell'\dashv_j r'$ consists of
    \[\begin{tikzcd}
        &C'\ar[ddr, "r'", bend left, ""{anchor=center,name=a}]&\\
        &\ar[u, "c"{description}] C \ar[dr, "r"]\\
        A \ar[uur, bend left, "\ell'"] \ar[ur, "\ell"] \ar[rr, "j"'] && E
        \ar[from=2-2, to=a, Rightarrow, "\rho", shorten=2mm]
    \end{tikzcd}\]
    \begin{enumerate}
        \item An arrow $c : C \to C'$ such that $\ell' = c\ell$.
        \item A cell $\rho : r \To r'c$.
    \end{enumerate}
    Rendering the following diagram commutative
    \[\begin{tikzcd}
        \ar[d, "{\smallfrown_c(\ell, 1)}"'] C(\ell, 1) \ar[r, "\sharp"] & E(j, r) \ar[d, "{E(j, \rho})"]\\
        C'(\ell', c) \ar[r, "{\sharp'(1, c)}"'] & E(j, r'c)
    \end{tikzcd}\]
    It is \textbf{strict} when $\rho$ is the identity. We write $\cat{RAdj}_R(j)$ the category of $j$-relative adjunctions and right morphisms.
\end{definition}

\begin{example}
    For any fully faithful arrow $j : A \to E$ and $j$-adjunction $\ell\dashv_j r$, the pair $(\ell, \eta)$ forms a unique right morphism from $\id_A\dashv_j j$ to $\ell\dashv_j r$. The latter is hence initial in $\cat{RAdj}_R(j)$.
\end{example}

As noted above, unicity of $\rho$ is not immediate and requires an hypothesis on the root. We write $A/\underline{\bb X}$ for the (1-)category of coslices under $A$.

\begin{lemma}
    \label{lem:lradj_ff}
    Let $j : A \to E$ be a dense arrow. Then the functor $\cat{RAdj}_R(j) \to A/\underline{\bb X}$ sending a $j$-adjunction to its left adjoint and a right morphism to its arrow is fully faithful.
\end{lemma}

\begin{proof}
    The compatibility condition for $\rho$ states that the following square commutes
    \[\begin{tikzcd}
        \ar[d, "{\smallfrown_c(\ell, 1)}"'] C(\ell, 1) & \ar[l, "\flat"'] E(j, r) \ar[d, "{E(j, \rho)}"]\\
        C'(\ell', c) \ar[r, "{\sharp'(1, c)}"'] & E(j, r'c)
    \end{tikzcd}\]
    The cell above uniquely determines $E(j, \rho)$. When $j$ is dense, this uniquely determines $\rho$ by \ref{lem:res_dense_fid}.
\end{proof}

Most of the properties regarding resolution of relative monads will rely on the intersection of the two notions, which we introduce here.

\begin{definition}
    A \textbf{strict morphism} of relative adjunction is a strict left (or equivalently right) morphism of relative adjunctions. We write $\cat{RAdj}(j)$ the category of $j$-adjunctions and strict morphisms.
\end{definition}

We conclude this Subsection by giving a few tools to construct new relative adjunctions from existing ones.

\begin{proposition}
    \label{prop:radj_precomp}
    Let $\ell\dashv_jr$ a relative adjunction and $\ell' : A\to B$ an arrow.
    \[\begin{tikzcd}
        &&C\ar[dr, "r"]&\\
        A \ar[r, "\ell'"] & B \ar[ur, "\ell"] \ar[rr, "j"'] &{}& D
        \ar[from=1-3, to=2-3, phantom, "\dashv"{anchor=center}]
    \end{tikzcd}
    \hspace{1em}
    \implies
    \hspace{1em}
    \begin{tikzcd}
        &C\ar[dr, "r"]&\\
        A \ar[ur, "\ell\ell'"] \ar[rr, "j\ell'"'] &{}& D
        \ar[from=1-2, to=2-2, phantom, "\dashv"{anchor=center}]
    \end{tikzcd}\]
    Then $(\ell\circ\ell')\dashv_{j\ell'}r$ and this assignment extends to functors
    \begin{align*}
        (-\circ\ell') &: \cat{RAdj}_L(j) \to \cat{RAdj}_L(j\ell')\\
        (-\circ\ell') &: \cat{RAdj}_R(j) \to \cat{RAdj}_R(j\ell')
    \end{align*}
\end{proposition}

\begin{proof}
    Using $\ell\dashv_j r$, we have $C(\ell\ell', 1) \simeqq D(j\ell', r)$.

    Given a left morphism $(c, \lambda)$ (resp. right morphism $(c, \rho)$) from $\ell_1\dashv_jr_1$ to $\ell_2\dashv_jr_2$, the pairs $(c, \lambda\ell')$ (resp. $(c, \rho)$) define left (resp. right) morphisms from $\ell_1\ell'\dashv_{j\ell'}r_1$ to $\ell_2\ell'\dashv_{j\ell'}r_2$. The compatibility condition follows from the condition on $(c, \lambda)$ (resp. $(c, \rho)$). Functoriality is immediate.
\end{proof}

The case of postcomposition on the right adjoint will be a corollary of the pasting law for relative adjunctions.

\begin{proposition}[Pasting law]
    \label{prop:radj_pasting}
    Consider the following configuration of arrows
    \[\begin{tikzcd}
        &C \ar[dr, "r"] &&\\
        && \ar[dr, "r'", ""{anchor=center,name=r_}] D &\\
        A \ar[uur, "\ell"] \ar[urr, "\ell'", ""{anchor=center,name=l_}] \ar[rrr, "j"'] &&& E\\
        \ar[from=l_, to=r_, phantom, "\dashv", shift right]
    \end{tikzcd}\]
    The left triangle is a relative adjunction $\ell\dashv_{\ell'}r$ if and only if the outer triangle is a relative adjunction $\ell\dashv_jr'r$. In this case, $(r, \eta)$ exhibits a left morphism of $j$-adjunction $(\ell\dashv_jr'r) \to (\ell'\dashv_jr)$ where $\eta$ is the unit of $\ell'\dashv_jr'$, and this assignment extends to functors
    \begin{align*}
        (r'\circ-) &: \cat{RAdj}_L(\ell') \to \cat{RAdj}_L(j)\\
        (r'\circ-) &: \cat{RAdj}_R(\ell') \to \cat{RAdj}_R(j)
    \end{align*}
\end{proposition}

\begin{proof}
    Using $\ell'\dashv_jr'$, we have $D(\ell', r) \simeqq E(j, r'r)$. The condition that the left (resp. outer) triangle be a relative adjunction asserts that the left (resp. right) side be isomorphic to $C(\ell, 1)$. $(r, \eta)$ exhibiting a left morphism follows by definition.

    Given a left morphism $(c, \lambda)$ (resp. right morphism $(c, \rho)$) from $\ell_1\dashv_{\ell'}r_1$ to $\ell_2\dashv_{\ell'}r_2$, the pairs $(c, \lambda)$ (resp. $(c, r'\rho)$) is a left (resp. right) from $\ell_1\dashv_{j'}r'r_1$ to $\ell_2\dashv_{j'}r'r_2$. The compatibility condition follows from the condition on $(c, \lambda)$ (resp. $(c, \rho)$). Functoriality is immediate.
\end{proof}

\begin{corollary}
    Let $\ell\dashv_j r$ a relative adjunction and $r' : E \to D$ a fully faithful arrow.
    \[\begin{tikzcd}
        &C\ar[dr, "r"]&&\\
        A \ar[ur, "\ell"] \ar[rr, "j"'] &{}& B \ar[r, "r'", hook] & D
        \ar[from=1-2, to=2-2, phantom, "\dashv"{anchor=center}]
    \end{tikzcd}
    \hspace{1em}
    \implies
    \hspace{1em}
    \begin{tikzcd}
        &C\ar[dr, "r'r"]&\\
        A \ar[ur, "\ell"] \ar[rr, "r'j"'] &{}& D
        \ar[from=1-2, to=2-2, phantom, "\dashv"{anchor=center}]
    \end{tikzcd}\]
    Then $\ell\dashv_{r'j}r'r$ and this assignment extends to functors
    \begin{align*}
        (r'\circ-) &: \cat{RAdj}_L(j) \to \cat{RAdj}_L(r'j)\\
        (r'\circ-) &: \cat{RAdj}_R(j) \to \cat{RAdj}_R(r'j)
    \end{align*}
\end{corollary}

\begin{proof}
    Apply Proposition \ref{prop:radj_pasting} with $\ell' := j$ and $j := r'j$.
\end{proof}

\begin{corollary}
    Let $\ell\dashv_j r$ and $j\ell'\dashv_{j'} r'$ be relative adjunctions as below:
    \[\begin{tikzcd}
        && C \ar[dr, "r"] &&\\
        & B \ar[ur, "\ell"] \ar[rr, "j"{description}] &{}& D \ar[dr, "r'", ""{anchor=center, name=r_}] &\\
        A \ar[ur, "\ell'"] \ar[urrr, "j\ell'"{description}, ""{anchor=center,name=jl_}] \ar[rrrr, "j'"'] &&&& E\\
        \ar[from=1-3, to=2-3, phantom, "\dashv"{anchor=center}]
        \ar[from=jl_, to=r_, phantom, "\dashv"{anchor=center}]
    \end{tikzcd}\]
    Then $\ell\ell'\dashv_{j'} r'r$ and $(r, \eta\ell')$ is a left morphism $(\ell\ell'\dashv_{j'}r'r) \to (j\ell'\dashv_{j'}r')$. This extends to functors
    \begin{align*}
        (r'\circ-\circ\ell') &: \cat{RAdj}_L(j) \to \cat{RAdj}_L(j')\\
        (r'\circ-\circ\ell') &: \cat{RAdj}_R(j) \to \cat{RAdj}_R(j')
    \end{align*}
\end{corollary}

\begin{proof}
    The result is a combination of Proposition \ref{prop:radj_precomp} by precomposing $\ell\dashv_jr$ by $\ell'$ and Proposition \ref{prop:radj_pasting} by pasting $\ell\ell'\dashv_{j\ell'} r$ along $j\ell'\dashv_{j'}r'$.
\end{proof}

\subsection{Relative monads}
\label{subsec:rmonad}

\begin{definition}
    \index{relative monad}\index{monad!relative}
    \begin{wrapfigure}{r}{0.47\textwidth}
        \vspace{-2\baselineskip}
        \[\begin{tikzcd}[sep=large]
            \ar[d, equal] A & \ar[l, loose, "{E(j, t)}"', ""{anchor=center, name=t}] A \ar[d, equal]\\
            \ar[d, equal] A & \ar[l, loose, "{E(t, t)}"', ""{anchor=center, name=m}] A \ar[d, equal]\\
            A & \ar[l, loose, "{E(j, t)}", ""{anchor=center, name=b}] A
            \ar[from=t, to=m, phantom, "\dag"{anchor=center}]
            \ar[from=m, to=b, phantom, "{E(\eta, t)}"{anchor=center}]
        \end{tikzcd}
        \hspace{1em}=\hspace{1em}
        \begin{tikzcd}[sep=large]
            \ar[d, equal] A & \ar[l, loose, "{E(j, t)}"', ""{anchor=center, name=t}] A \ar[d, equal]\\
            A & \ar[l, loose, "{E(j, t)}", ""{anchor=center, name=b}] A
            \ar[from=t, to=b, phantom, "="{anchor=center}]
        \end{tikzcd}\]
        \vspace{-7\baselineskip}
    \end{wrapfigure}
    A \textbf{relative monad} consists of:
    \begin{enumerate}
        \item An arrow $j : A \to E$, the \textbf{root}
        \item An arrow $t : A \to E$, the \textbf{underlying arrow}
        \item A cell $\dag : E(j, t) \To E(t, t)$, the \textbf{extension}
        \item A cell $\eta : j \To t$, the \textbf{unit}
    \end{enumerate}
    Satisfying the following equations:
    \[\begin{tikzcd}[sep=large]
        \ar[d, equal] A & \ar[l, equal, ""{anchor=center, name=t}] A \ar[d, equal]\\
        \ar[d, equal] A & \ar[l, loose, "{E(j, j)}", ""{anchor=center, name=mt}] A \ar[d, equal]\\
        \ar[d, equal] A & \ar[l, loose, "{E(j, t)}", ""{anchor=center, name=mb}] A \ar[d, equal]\\
        A & \ar[l, loose, "{E(t, t)}", ""{anchor=center, name=b}] A
        \ar[from=t, to=mt, phantom, "\smallfrown_j"{anchor=center}]
        \ar[from=mt, to=mb, phantom, "{E(\eta, t)}"{anchor=center}]
        \ar[from=mb, to=b, phantom, "\dag"{anchor=center}]
    \end{tikzcd}
    \hspace{1em}=\hspace{1em}
    \begin{tikzcd}[sep=large]
        \ar[d, equal] A & \ar[l, equal, ""{anchor=center, name=t}] A \ar[d, equal]\\
        A & \ar[l, loose, "{E(t, t)}", ""{anchor=center, name=b}] A
        \ar[from=t, to=b, phantom, "\smallfrown_j"{anchor=center}]
    \end{tikzcd}\]
    \[\begin{tikzcd}[sep=large]
        \ar[d, equal] A & \ar[l, loose, "{E(j, t)}"', ""{anchor=center, name=t1}] A \ar[d, equal] & \ar[l, loose, "{E(j, t)}"', ""{anchor=center, name=t2}] A \ar[d, equal]\\
        \ar[d, equal] A & \ar[l, loose, "{E(t, t)}"', ""{anchor=center, name=m1}] A & \ar[l, loose, "{E(t, t)}"', ""{anchor=center, name=m2}] A \ar[d, equal] \\
            A &{}& \ar[ll, loose, "{E(t, t)}"] A
            \ar[from=t1, to=m1, phantom, "\dag"{anchor=center}] \ar[from=t2, to=m2, phantom, "\dag"{anchor=center}]
            \ar[from=2-2, to=3-2, phantom, "{\smallsmile_t(t, t)}"{anchor=center}]
    \end{tikzcd}
    \hspace{1em}=\hspace{1em}
    \begin{tikzcd}[sep=large]
        \ar[d, equal] A & \ar[l, loose, "{E(j, t)}"', ""{anchor=center,name=t1}] A \ar[d, equal] & \ar[l, loose, "{E(j, t)}"', ""{anchor=center,name=t2}] A \ar[d, equal]\\
        \ar[d, equal] A & \ar[l, loose, "{E(j, t)}"', ""{anchor=center,name=m1}] A & \ar[l, loose, "{E(t, t)}"', ""{anchor=center,name=m2}] A\ar[d, equal]\\
        \ar[d, equal] A &{}& \ar[ll, loose, "{E(j, t)}"'] A \ar[d, equal]\\
        A &{}& \ar[ll, loose, "{E(t, t)}"] A
        \ar[from=t1, to=m1, phantom, "="{anchor=center}] \ar[from=t2, to=m2, phantom, "\dag"{anchor=center}]
        \ar[from=2-2, to=3-2, phantom, "{\smallsmile_t(j, t)}"{anchor=center}]
        \ar[from=3-2, to=4-2, phantom, "\dag"{anchor=center}]
    \end{tikzcd}\]
    A morphism of $j$-relative monads from $(t, \dag, \eta)$ to $(t', \dag', \eta')$ is a cell $\tau : t \To t'$ such that
    \[\begin{tikzcd}
        & \ar[dl, "\eta"'] j \ar[dr, "\eta'"] &\\
        t \ar[rr, "\tau"'] && t'
    \end{tikzcd}
    \hspace{1em}
    \begin{tikzcd}
        \ar[d, "\dag"'] E(j, t) \ar[rr, "{E(j, \tau)}"] && E(j, t') \ar[d, "\dag'"]\\
        E(t, t) \ar[dr, "{E(t, \tau)}"'] && \ar[dl, "{E(\tau, t')}"] E(t', t')\\
                                         & E(t, t')
    \end{tikzcd}\]
    We write $\cat{RMnd}(j)$ for the (1-)category of $j$-relative monads and morphisms.
\end{definition}

\begin{example}
    For any arrow $j : A \to E$, $(j, e_{E(j, j)}, e_j)$ forms a $j$-relative monad. %TODO: initiale dans RMod(j)
\end{example}

A large class of relative adjunctions is given by the ones induced by relative adjunctions. Note that, not every relative monad necessarily induce a relative adjunction as it requires an object acting as a category of algebras. This is one of the subject of Section \ref{sec:alg_idem_mon}.

\begin{theorem}
    \label{th:radj_rmon}
    Every relative adjunction $\ell\dashv_jr$ induces a relative monad $\land_j(\ell \dashv_jr)$ with underlying cell $rl$ and this assignment extends to functors
    \begin{align*}
        \obslash_j &: \cat{RAdj}_L(j) \to \cat{RMnd}_L(j)^\op\\
        \oslash_j  &: \cat{RAdj}_R(j) \to \cat{RMnd}_R(j)
    \end{align*}
\end{theorem}

\begin{proof}
    We define $\eta$ as in Proposition \ref{prop:equiv_radj_car} and $\dag$ as the cell
    \[\begin{tikzcd}[sep=huge]
        \ar[d, equal] A & \ar[l, loose, "{E(j, r\ell)}"', ""{anchor=center,name=t}] A \ar[d, equal] \\
        \ar[d, equal] A & \ar[l, loose, "{C(\ell, \ell)}", ""{anchor=center,name=m}] C \ar[d, equal] \\
        A & \ar[l, loose, "{E(r\ell, r\ell)}", ""{anchor=center,name=b}] A
        \ar[from=t, to=m, phantom, "{\flat(1, \ell)}"{anchor=center}]
        \ar[from=m, to=b, phantom, "{\smallfrown_r(\ell, \ell)}"{anchor=center}]
    \end{tikzcd}\]
    %TODO: Finir la preuve (?)
\end{proof}

We are in fact interested in when a relative monad is induced by a relative adjunction. Most of the results will come in Section \ref{sec:alg_idem_mon}, but we introduce a few not needing (op)algebras.

\begin{definition}
    Let $j : A \to E$ an arrow and $T$ a $j$-monad. A resolution of $T$ is a $j$-adjunction $\ell\dashv_jr$ for which $T$ is equal to $\land_j(\ell\dashv_jr)$. A morphism of resolutions of $T$ from $\ell\dashv_jr$ to $\ell'\dashv_jr'$ is an arrow $c : C \to C'$ such that the following diagram commutes
    \[\begin{tikzcd}
        & C \ar[dd, "c"{description}] \ar[dr, "r"] &\\
        A \ar[ur, "\ell"] \ar[dr, "\ell'"'] && E\\
        & C' \ar[ur, "r'"']
    \end{tikzcd}\]
    We write $\cat{Res}(T)$ for the category of resolutions of $T$ and their morphisms.
\end{definition}

\begin{lemma}
    Let $\ell\dashv_jr$ and $\ell'\dashv_jr'$ be relative adjunctions. An arrow $c : C \to C'$ satisfying $c\ell = \ell'$ and $rc = r'$ is a strict morphism from $\ell\dashv_jr$ to $\ell'\dashv_jr'$ if and only if $\land_j(\ell\dashv_jr) = \land_j(\ell'\dashv_jr')$.
\end{lemma}

\begin{proof}
    Both relative monads are equal on their underlying cells, the compatibility condition for a strict morphism is equivalent to $\eta = \eta'$ and $\dag = \dag'$ using Theorem \ref{th:radj_rmon}.
\end{proof}

\begin{lemma}
    \label{lem:reso_lmon}
    Let $\ell\dashv_jr$ be a resolution of a relative monad $T$. The loose-monad induced by $C(1, \ell)\dashv E(j, r)$ is isomorphic to $E(j, T)$.
\end{lemma}

\begin{proof}
    The underlying distributor of each is $E(j, r\ell)$, and the units and extension coincide by diagrammatic reasoning.
\end{proof}

\begin{corollary}
    Let $\ell\dashv_jr$ a resolution of $T$. The distributor $C(\ell, \ell)$ is isomorphic to $E(j, T)$.
\end{corollary}

\begin{proof}
    By Lemma \ref{lem:reso_lmon}, $E(j, T)$ is induced by $C(1, \ell) \dashv E(j, r)$, however $E(j, r) \simeqq C(\ell, 1)$. Therefore, $E(j, T)$ is isomorphic to the loose monad induced by $C(1, \ell)\dashv C(\ell, 1)$ which is $C(\ell, \ell)$.
\end{proof}

\begin{corollary}
    Let $(\ell_2\ell_1\dashv_j r$ a relative adjunction with unit $\eta$. The unit then exhibits a relative adjunction $\ell_1\dashv_jr\ell_2$ if and only if the identity $e_{\ell_2\ell_2}$ exhibits a relative adjunction $\ell_1\dashv_{\ell_2\ell_1}\ell_2$. In this case, the induced $j$-monads are identical.
    \[\begin{tikzcd}
        B \ar[r, "\ell_2"] & C \ar[dr, "r"] &\\
        A \ar[u, "\ell_1"] \ar[ur, "\ell_2\ell_1"{description}] \ar[rr, "j"'] &{}& E
        \ar[from=1-2, to=2-2, phantom, "\dashv"]
    \end{tikzcd}
    \hspace{1em}=\hspace{1em}
    \begin{tikzcd}
        & B \ar[dr, "r\ell_2"] &\\
        A \ar[ur, "\ell_1"] \ar[rr, "j"'] &{}& E
        \ar[from=1-2, to=2-2, phantom, "\dashv"]
    \end{tikzcd}\]
\end{corollary}

\begin{proof}
    We use Proposition \ref{prop:radj_pasting} with that fact that the unit of $\ell_1\dashv_jr\ell_2$ forming a left morphism between the two relative monads is the identity.
\end{proof}

As with relative adjunctions, we can precompose a relative monad by an arrow:

\begin{proposition}
    \label{prop:rmnd_precomp}
    Let $j : B \to D$ an arrow, $T = (t, \dag, \eta)$ a $j$-monad and $\ell :  A \to B$ an arrow. Then $t\ell$ may be equipped with the structure of a $j\ell$-monad, and this assignment extends to a functor $\cat{RMnd}(j) \to \cat{RMnd}(j\ell)$ rendering the following diagram commutative
    \[\begin{tikzcd}
        \ar[d, "(-\circ\ell)^\op"'] \cat{RAdj}_L(j)^\op \ar[r, "\obslash_j"] & \cat{RMnd}(j) \ar[d, "(-\circ\ell)"{description}] & \ar[l, "\oslash_j"'] \cat{RAdj}_R(j) \ar[d, "(-\circ\ell)"]\\
        \cat{RAdj}_L(j\ell)^\op \ar[r, "\obslash_{j\ell}"'] & \cat{RMnd}(j\ell) & \ar[l, "\oslash_{j\ell}"] \cat{RAdj}_R(j\ell)
    \end{tikzcd}\]
\end{proposition}

\begin{proof}
    The unit is given by $\eta\ell$ and the extension operator by $\dag(\ell, \ell)$ and the relative monad laws follow from those of $T$. Given a $j$-monad morphism $\tau : T \to T'$, the cell $\tau\ell$ forms a $j\ell$-monad morphism $T\ell\to T'\ell$, and the laws follows from those of $\tau$. Functoriality is immediate and commutativity follows from the definition.
\end{proof}

And paste a relative monad along a left adjoint:

\begin{proposition}
    \label{prop:rmnd_paste}
    Let $\ell\dashv_j r$ a relative adjunction and $T = (t, \dag, \eta)$ be a $\ell$-monad.
    \[\begin{tikzcd}%FIXME
        & D \ar[dr, "r'"] &\\
        A \ar[ur, "\ell"{description}, ""{anchor=center,name=l}] \ar[ur, "t", bend left=50, "t", ""{anchor=center,name=t}] \ar[rr, "j"'] &{}& E
        \ar[from=l, to=t, Rightarrow, "\eta"'] \ar[from=1-2, to=2-2, phantom, "\dashv"{anchor=center}]
    \end{tikzcd}\]
    Then $r't$ may be equipped with the structure of a $j$-monad for which $r'\eta$ is a $j$-monad morphism $\land_j(\ell\dashv_jr) \to r'T$ and this assignment extends to a functor $\cat{RMnd}(\ell) \to \cat{RMnd}(j)$ rendering the following diagram commutative
    \[\begin{tikzcd}
        \ar[d, "(r\circ-)^\op"'] \cat{RAdj}_L(\ell)^\op \ar[r, "\obslash_\ell"] & \cat{RMnd}(\ell) \ar[d, "(r\circ-)"{description}] & \ar[l, "\oslash_\ell"'] \cat{RAdj}_R(\ell) \ar[d, "(r\circ-)"]\\
        \cat{RAdj}_L(j)^\op \ar[r, "\obslash_{j}"'] & \cat{RMnd}(j) & \ar[l, "\oslash_{j}"] \cat{RAdj}_R(j)
    \end{tikzcd}\]
\end{proposition}

\begin{proof}
    We define the extension and unit operators as below
    \[\begin{tikzcd}[sep=huge]
        \ar[d, equal] A & \ar[l, loose, "{E(j, rt)}"', ""{anchor=center,name=t}] A \ar[d, equal]\\
        \ar[d, equal] A & \ar[l, loose, "{D(\ell, t)}"', ""{anchor=center,name=mt}] A \ar[d, equal]\\
        \ar[d, equal] A & \ar[l, loose, "{D(t, t)}"', ""{anchor=center,name=mb}] A \ar[d, equal]\\
        A & \ar[l, loose, "{E(rt, rt)}", ""{anchor=center, name=b}] A
        \ar[from=t, to=mt, phantom, "{\flat(1, t)}"{anchor=center}]
        \ar[from=mt, to=mb, phantom, "\dag"{anchor=center}]
        \ar[from=mb, to=b, phantom, "{\flat(t, t)}"{anchor=center}]
    \end{tikzcd}
    \hspace{3em}
    \begin{tikzcd}[sep=huge]
        A \ar[d, "j"'] \ar[r, equal, ""{anchor=center,name=t1}] & A\ar[d, "r\ell"] \ar[r, equal, ""{anchor=center,name=t2}] & A \ar[d, "rt"]\\
        A \ar[r, loose, equal, ""{anchor=center,name=b1}] & A \ar[r, loose, equal, ""{anchor=center,name=b2}] & A
        \ar[from=t1, to=b1, phantom, "\eta'"] \ar[from=t2, to=b2, phantom, "\eta"]
    \end{tikzcd}\]
    The unit and associativity laws follows from those of $T$ and the $\sharp-\flat$ isomorphism. Functoriality is immediate. The squares agree on objects and the following diagram commutes
    \[\begin{tikzcd}
        \ar[d, "\land_\ell"'] \cat{RAdj}_L(\ell) \ar[r, "{(r\circ-)}"] & \cat{RAdj}(j) \ar[d, "\land_j"]\\
        \cat{RMnd}(\ell) \ar[r, "{(r \circ -)}"] & \cat{RMnd}(j)
    \end{tikzcd}\]
    In fact, the assignments agree on the underlying cell and on the unit; and the $\sharp-\flat$ isomorphism gives the equality on the extension operators.
\end{proof}

Using both Propositions \ref{prop:rmnd_precomp} and \ref{prop:rmnd_paste}, we get:

\begin{corollary}
    Let $j\ell\dashv_{j'}r$ be a relative adjunction and let $(t, \dag, \eta)$ be a $j$-monad.
    \[\begin{tikzcd}
        & B \ar[r, "j"{description}] \ar[r, "t", bend left=50] & D \ar[dr, "r", ""{anchor=center,name=r_}] &\\
        A \ar[ru, "\ell"] \ar[rru, "j\ell"{description}, ""{anchor=center,name=jl}] \ar[rrr, "j"'] &&& E
        \ar[from=jl, to=r_, phantom, "\dashv"{anchor=center}]
    \end{tikzcd}\]
    Then $rt\ell$ may be equipped with the structure of a $j$-monad for which $r\eta\ell$ is a $j$-monad morphism $rj\ell \to rt\ell$. This assignment extends to a functor $\cat{RMnd}(j) \to \cat{RMnd}(j')$ rendering the following diagram commutative.
    \[\begin{tikzcd}
        \ar[d, "(r\circ-\circ\ell)^\op"'] \cat{RAdj}_L(j)^\op \ar[r, "\obslash_j"] & \cat{RMnd}(j) \ar[d, "(r\circ-\circ\ell)"{description}] & \ar[l, "\oslash_j"'] \cat{RAdj}_R(\ell) \ar[d, "(r\circ-\circ\ell)"]\\
        \cat{RAdj}_L(j')^\op \ar[r, "\obslash_{j'}"'] & \cat{RMnd}(j') & \ar[l, "\oslash_{j'}"] \cat{RAdj}_R(j')
    \end{tikzcd}\]
\end{corollary}

\begin{proof}
    Precompose $t$ by $\ell$ using Proposition \ref{prop:rmnd_precomp} and paste $t\ell$ along $j\ell\dashv_{j'}r$ using Proposition \ref{prop:rmnd_paste}.
\end{proof}


\section{Constructions and properties of relative monads}
\label{sec:alg_idem_mon}

In this Section, we focus of the properties of relative monads, following \cite[Section 6]{arkor2024formaltheoryrelativemonads}, \cite{arkor2025idempotencerelativemonads}, \cite{arkor2025relativemonadicity}. Starting with the construction of the right notions of algebras and categories of algebras in a virtual equipment, we continue with idempotency before concluding with a monadicity theorem.

\subsection{(op)Algebras and (op)algebras object}

In the classical settings, two categories of algebras are usually studied: the Eilenberg-Moore category ("all the algebras") and the Kleisli category ("free algebras"). Internal notions in virtual equipment exist: algebras and opalgebras. While the first ones are known to the reader, the latter are less so. Although in the classical setting free algebras and opalgebras coincide, in the formal setting the latter notion is of importance.

This subsection is based on \cite[Section 6]{arkor2024formaltheoryrelativemonads}.

\begin{definition}
    \label{def:alg}
    \index{$T$-algebra} \index{$T$-algebra!morphism} \index{$T-\Alg_B$}
    \begin{wrapfigure}{r}{0.47\textwidth}
        \vspace{-\baselineskip}
        \[\begin{tikzcd}[sep=large]
            \ar[d, equal] A & \ar[l, loose, "{E(j, e)}"', ""{anchor=center,name=t}] D \ar[d, equal]\\
            \ar[d, equal] A & \ar[l, loose, "{E(t, e)}"', ""{anchor=center,name=m}] D \ar[d, equal]\\
            A & \ar[l, loose, "{E(j, e)}", ""{anchor=center,name=b}] D
            \ar[from=t, to=m, phantom, "\rtimes"{anchor=center}] \ar[from=m, to=b, phantom, "{E(\eta, e)}"{anchor=center}]
        \end{tikzcd}
        \hspace{1em}=\hspace{1em}
        \begin{tikzcd}
            \ar[d, equal] A & \ar[l, loose, "{E(j, e)}"', ""{anchor=center, name=t}] D\ar[d, equal]\\
            A & \ar[l, loose, "{E(j, e)}", ""{anchor=center, name=b}] D
            \ar[from=t, to=b, phantom, "="{anchor=center}]
        \end{tikzcd}\]
        \vspace{-7\baselineskip}
    \end{wrapfigure}
    Let $T$ be a relative monad. An \textbf{algebra for $T$} consists of:
    \begin{enumerate}
        \item An object $D$, the \textbf{domain}
        \item An arrow $e : D \to E$, the \textbf{underlying arrow}
        \item A cell $\rtimes : E(j, e) \To E(t, e)$
    \end{enumerate}
    Satisfying the following equations:
    \vspace{2\baselineskip}
    \[\begin{tikzcd}[sep=large]
        \ar[d, equal] A & \ar[l, loose, "{E(j, t)}"', ""{anchor=center, name=t1}] A \ar[d, equal] & \ar[l, loose, "{E(j, e)}"', ""{anchor=center, name=t2}] D \ar[d, equal]\\
        \ar[d, equal] A & \ar[l, loose, "{E(t, t)}"', ""{anchor=center, name=m1}] A & \ar[l, loose, "{E(t, e)}"', ""{anchor=center, name=m2}] D \ar[d, equal]\\
        A &{}& \ar[ll, loose, "{E(t, e)}", ""{anchor=center, name=b}] D
        \ar[from=t1, to=m1, phantom, "\dag"{anchor=center}] \ar[from=t2, to=m2, phantom, "\rtimes"{anchor=center}]
        \ar[from=2-2, to=3-2, phantom, "{\smallsmile_t(t, e)}"{anchor=center}]
    \end{tikzcd}
    \hspace{1em}=\hspace{1em}
    \begin{tikzcd}[sep=large]
        \ar[d, equal] A & \ar[l, loose, "{E(j, t)}"', ""{anchor=center,name=t1}] A \ar[d, equal] & \ar[l, loose, "{E(j, e)}"', ""{anchor=center, name=t2}] D \ar[d, equal]\\
        \ar[d, equal] A & \ar[l, loose, "{E(j, t)}"', ""{anchor=center,name=m1}] A & \ar[l, loose, "{E(t, e)}"', ""{anchor=center, name=m2}] D\ar[d, equal]\\
        \ar[d, equal] A &{}& \ar[ll, loose, "{E(j, e)}"] D \ar[d, equal]\\
        A &{}& \ar[ll, loose, "{E(t, e)}"] D
        \ar[from=t1, to=m1, phantom, "="{anchor=center}] \ar[from=t2, to=m2, phantom, "\rtimes"{anchor=center}]
        \ar[from=2-2, to=3-2, phantom, "{\smallsmile_t(j, e)}"{anchor=center}]
        \ar[from=3-2, to=4-2, phantom, "\rtimes"{anchor=center}]
    \end{tikzcd}\]
    A $T$-algebra morphism from $(e, \rtimes)$ to $(e', \rtimes')$ with common domain $D$ is a cell $\varepsilon : e \To e'$ such that
    \[\begin{tikzcd}
        \ar[d, "{E(j, \varepsilon)}"'] E(j, e) \ar[r, "\rtimes"] & E(t, e) \ar[d, "{E(t, \varepsilon)}"]\\
        E(j, e') \ar[r, "\rtimes'"'] & E(t, e')
    \end{tikzcd}\]
    We write $T-\Alg_D$ for the category of $T$-algebras with domain $D$ and their morphisms. This category is contravariantely functorial in $D$ and $T$.
\end{definition}

\begin{remark}
    \label{rem:ralg_act}
    Note that, using Corollary \ref{cor:bending}, we can \textit{bend} the arrows and rewrite the extension operator as $\lambda : E(1, t), E(j, e) \To E(1, e)$. This gives a link between the description as action in the multiactegory used in \cite[Theorem 6.15]{arkor2024formaltheoryrelativemonads}.
\end{remark}

We now introduce the other notion:

\begin{definition}
    \label{def:opalg}
    \index{$T$-opalgebra} \index{$T$-opalgebra!morphism} \index{$T-\Opalg_B$}
    \begin{wrapfigure}{r}{0.47\textwidth}
        \vspace{-\baselineskip}
        \[\begin{tikzcd}[sep=large]
            \ar[d, equal] A & \ar[l, equal, ""{anchor=center,name=t}] A \ar[d, equal]\\
            \ar[d, equal] A & \ar[l, loose, "{E(j, j)}"', ""{anchor=center,name=mt}] A \ar[d, equal]\\
            \ar[d, equal] A & \ar[l, loose, "{E(j, t)}"', ""{anchor=center,name=mb}] A \ar[d, equal]\\
            A & \ar[l, loose, "{B(a, a)}", ""{anchor=center,name=b}] A
            \ar[from=t, to=mt, phantom, "\smallfrown_j"{anchor=center}]
            \ar[from=mt, to=mb, phantom, "{E(j, \eta)}"{anchor=center}]
            \ar[from=mb, to=b, phantom, "\ltimes"{anchor=center}]
        \end{tikzcd}
        \hspace{1em}=\hspace{1em}
        \begin{tikzcd}
            \ar[d, equal] A & \ar[l, equal, ""{anchor=center, name=t}] A\ar[d, equal]\\
            A & \ar[l, loose, "{B(a, a)}", ""{anchor=center, name=b}] A
            \ar[from=t, to=b, phantom, "\smallfrown_a"{anchor=center}]
        \end{tikzcd}\]
    \end{wrapfigure}
    Let $T$ be a relative monad. An \textbf{opalgebra for $T$} consists of:
    \begin{enumerate}
        \item An object $B$, the \textbf{codomain}
        \item An arrow $a : A \to B$, the \textbf{underlying arrow}
        \item A cell $\ltimes : E(j, t) \To B(a, a)$
    \end{enumerate}
    Satisfying the following equations:
    \vspace{3\baselineskip}
    \[\begin{tikzcd}[sep=large]
        \ar[d, equal] A & \ar[l, loose, "{E(j, t)}"', ""{anchor=center, name=t1}] A \ar[d, equal] & \ar[l, loose, "{E(j, t)}"', ""{anchor=center, name=t2}] A \ar[d, equal]\\
        \ar[d, equal] A & \ar[l, loose, "{B(a, a)}"', ""{anchor=center, name=m1}] A & \ar[l, loose, "{B(a, a)}"', ""{anchor=center, name=m2}] A \ar[d, equal]\\
        A &{}& \ar[ll, loose, "{B(a, a)}", ""{anchor=center, name=b}] A
        \ar[from=t1, to=m1, phantom, "\ltimes"{anchor=center}] \ar[from=t2, to=m2, phantom, "\ltimes"{anchor=center}]
        \ar[from=2-2, to=3-2, phantom, "{\smallsmile_a(a, a)}"{anchor=center}]
    \end{tikzcd}
    \hspace{1em}=\hspace{1em}
    \begin{tikzcd}[sep=large]
        \ar[d, equal] A & \ar[l, loose, "{E(j, t)}"', ""{anchor=center,name=t1}] A \ar[d, equal] & \ar[l, loose, "{E(j, t)}"', ""{anchor=center, name=t}] A \ar[d, equal]\\
        \ar[d, equal] A & \ar[l, loose, "{E(j, t)}"', ""{anchor=center,name=m1}] A & \ar[l, loose, "{E(t, t)}"', ""{anchor=center, name=m}] A\ar[d, equal]\\
        \ar[d, equal] A &{}& \ar[ll, loose, "{E(j, t)}"] A \ar[d, equal]\\
        A &{}& \ar[ll, loose, "{B(a, a)}"] A
        \ar[from=t1, to=m1, phantom, "="{anchor=center}] \ar[from=t2, to=m2, phantom, "\dag"{anchor=center}]
        \ar[from=2-2, to=3-2, phantom, "{\smallsmile_t(j, t)}"{anchor=center}]
        \ar[from=3-2, to=4-2, phantom, "\ltimes"{anchor=center}]
    \end{tikzcd}\]
    A $T$-opalgebra morphism from $(a, \ltimes)$ to $(a', \ltimes')$ with common codomain $B$ is a cell $\alpha : a \To a'$ such that
    \[\begin{tikzcd}
        \ar[d, "\ltimes'"'] E(j, e) \ar[r, "\ltimes"] & B(a, a) \ar[d, "{B(a, \alpha)}"]\\
        B(a', a') \ar[r, "{B(\alpha, a')}"'] & B(a, a')
    \end{tikzcd}\]
    We write $T-\Opalg_B$ for the category of $T$-opalgebras with codomain $B$ and their morphisms. This category is covariantely functorial in $B$ and $T$.
\end{definition}

\begin{remark}
    \label{rem:ropalg_act}
    Once again, using Corollary \ref{cor:bending} we can bend the arrows and write the extension operator as $\rho : B(1, a), E(j, t) \To B(1, a)$. This gives a link between the description as action in the multiactegory used in \cite[Theorem 6.22]{arkor2024formaltheoryrelativemonads}.
\end{remark}

We can express $T$-opalgebras in terms of loose monad morphism:

\begin{lemma}
    \label{lem:opalg_lmnd_mor}
    Let $j : A \to E$ be an arrow and $T$ a $j$-monad. A $T$-opalgebra is an arrow $a : A \to B$ and a loose monad morphism $E(j, T) \to B(a, a)$.
\end{lemma}

\begin{proof}
    Follows from $T$-opalgebra laws (compare the diagrams).
\end{proof}

The first immediate example of a $T$-algebra and a $T$-opalgebra is when a relative monad is induced by a relative adjunction.

\begin{proposition}
    \label{prop:radj_alg_opalg}
    Let $\ell\dashv_jr$ a $j$-relative adjunction and $T$ the induced $j$-monad. The left $j$-adjoint forms a $T$-opalgebra and the right $j$-adjoint forms a $T$-algebra.
\end{proposition}

\begin{proof}
    We define the two extensions $\ltimes : E(j, r\ell) \To C(\ell, \ell)$ and $\rtimes : E(j, r) \To E(r\ell, r)$ as $\flat(1, \ell)$ and $\smallfrown_r(\ell, 1)\circ\flat$ and the proofs of compatibility are as in \ref{th:radj_rmon}.
\end{proof}

As with relative adjunctions and relative monads, we can construct new algebras and opalgebras by pre and postcomposition. And as before, precomposition ("on the left") works without hypothesis while postcomposition ("on the right") requires one.

\begin{proposition}
    Let $j : A \to E$ an arrow and $T$ a $j$-monad. Let $\ell : X \to A$ an arrow. Every $T$-opalgebra induces a $T\ell$-opalgebra and every $T$-algebra induces a $T\ell$-algebra, and this assignment extends to functors
    \begin{align*}
        (-\circ\ell) &: T-\Opalg_C \to T\ell-\Opalg_C\\
        (-\circ\ell) &: T-\Alg_C \to T\ell-\Alg_C
    \end{align*}
    natural in $T$ and $C$. This assignment is also compatible with resolutions in the following sense:
    \[\begin{tikzcd}
        \ar[d, "{(-\circ\ell)}"'] T-\Opalg_C & \lar \cat{Res}(T)_C \ar[d, "{(-\circ\ell)}"] \rar & T-\Alg_C \ar[d, "{(-\circ\ell)}"]\\
        T\ell-\Opalg_C & \lar \cat{Res}(T\ell)_C \rar & T\ell-\Alg_C
    \end{tikzcd}\]
    Where the unlabelled  arrows are given by \ref{prop:radj_alg_opalg}
\end{proposition}

\begin{proof}
    Given a $T$-opalgebra $(b, \ltimes)$ (resp. $T$-algebra $(d, \rtimes)$) we have a cell $E(j\ell, t\ell=\To C(b\ell, b\ell)$ (resp. $E(j\ell, d) \To E(t\ell, d)$) by precomposing $\ell$ in the both (resp. the first) arguments. The fact that these define $T\ell$-algebra and opalgebra is a simple verification analogous to the proof in \ref{prop:rmnd_precomp}. Functoriality is immediate and the diagram commutes by definition of opalgebras and algebras.
\end{proof}

\begin{proposition}
    Let $\ell\dashv_jr$ and $T = (t, \dag, \eta)$ an $\ell$-monad. Every $T$-opalgebra induces a $rT$-opalgebra and every $T$-algebra induces a $rT$-algebra, and this assignment extends to functors
    \begin{align*}
        (r\circ-) &: T-\Opalg_C \to rT-\Opalg_C\\
        (r\circ-) &: T-\Alg_C \to rT-\Alg_C
    \end{align*}
    natural in $T$ and $C$. This assignment is also compatible with resolutions in the following sense:
    \[\begin{tikzcd}
        \ar[d, "{(r\circ-)}"'] T-\Opalg_C & \lar \cat{Res}(T)_C \ar[d, "{}(r\circ-)"] \rar & T-\Alg_C \ar[d, "{(r\circ-)}"]\\
        T\ell-\Opalg_C & \lar \cat{Res}(T\ell)_C \rar & T\ell-\Alg_C
    \end{tikzcd}\]
\end{proposition}

\begin{proof}
    Given a $T$-opalgebra $(a : A \to B, \ltimes)$, define a cell $E(j, rt) \simeqq D(\ell, t) \To C(a, a)$ by applying the relative adjunction. For a $T$-algebra $(e, \rtimes)$ we define $E(j, re) \To E(rt, re)$ by $\smallfrown_r\rtimes(\flat, =)$. Once again the verifications that these cells define $rT$-algebra and opalgebra are analogous to \ref{prop:rmnd_paste} and functoriality and commutativity of the diagram are immediate and by definition respectively.
\end{proof}

\begin{corollary}
    In fact, $(r\circ-) : T-\Opalg_C \to rT-\Opalg_C$ is invertible.
\end{corollary}

We now want to assemble algebras and opalgebras into their own internal objects. To write the definition, we need a more general definition of morphism between (op)algebras. In particular, we want to remove the constraint of the (co)domain and for that we parametrize ("grade") the morphism with distributors.

\begin{definition}
    \index{$T$-algebra!graded morphism} \index{$T$-algebra!ungraded morphism}
    Let $(e : D \to E, \rtimes), (e' : D'\to E, \rtimes')$ be $T$-algebras. A \textbf{$\vec p$-graded $T$-algebra morphism} from $(e, \rtimes)$ to $(e', \rtimes')$ is a cell
    $$\varepsilon : E(1, e), \vec p \To E(1, e')$$
    Such that
    \[\begin{tikzcd}
        \ar[d, equal] E & \ar[l, loose, "{E(1, t)}"'] A & \ar[l, loose, "{E(j, e)}"'] \ar[d, equal] D & \ar[l, loose, "p_1"'] \cdots & \ar[l, loose, "p_n"'] D' \ar[d, equal]\\
        \ar[d, equal] E &{}& \ar[ll, loose, "{E(1, e)}"] D & \ar[l, loose, "p_1"'] \cdots & \ar[l, loose, "p_n"'] D' \ar[d, equal]\\
        E &&{}&& \ar[llll, loose, "{E(1, e')}", ""{anchor=center,name=b}] D'
        \ar[from=1-2, to=2-2, phantom, "\lambda"{anchor=center}] \ar[from=1-4, to=2-4, phantom, "="{anchor=center}]
        \ar[from=2-1, to=2-5, phantom, ""{anchor=center,name=m}]
        \ar[from=m, to=b, phantom, "\varepsilon"{anchor=center}]
    \end{tikzcd}
    \hspace{1em}=\hspace{1em}
    \begin{tikzcd}
        \ar[d, equal] E & \ar[l, loose, "{E(1, t)}"', ""{anchor=center,name=t1}] A \ar[d, equal] & \ar[l, loose, "{E(j, 1)}"', ""{anchor=center,name=t2}] \ar[d, equal] E & \ar[l, loose, "{E(1, e)}"'] D & \ar[l, loose, "p_1"'] \cdots & \ar[l, loose, "p_n"'] D' \ar[d, equal]\\
        \ar[d, equal] E & \ar[l, loose, "{E(1, t)}", ""{anchor=center,name=m1}] A & \ar[l, loose, "{E(j, 1)}", ""{anchor=center,name=m2}] E &&& \ar[lll, loose, "{E(1, e')}", ""{anchor=center,name=m3}] D' \ar[d, equal]\\
            E &&&&& \ar[lllll, loose, "{E(1, e')}", ""{anchor=center,name=b}] D'
            \ar[from=1-3, to=1-6, phantom, ""{anchor=center,name=t3}]
            \ar[from=2-1, to=2-6, phantom, ""{anchor=center,name=m}]
            \ar[from=t1, to=m1, phantom, "="{anchor=center}] \ar[from=t2, to=m2, phantom, "="{anchor=center}] \ar[from=t3, to=m3, phantom, "\varepsilon"{anchor=center}]
            \ar[from=m,to=b,phantom,"\lambda'"{anchor=center}]
    \end{tikzcd}\]
    Where $\lambda$ is as in Remark \ref{rem:ralg_act}. When $n=0$, such a morphism is \textbf{ungraded}.
\end{definition}

In particular, an ungraded $T$-algebra morphism is exactly a $T$-algebra morphism as given in Definition \ref{def:alg}.

\begin{remark}
    Once again bending the arrows (using Corollary \ref{cor:bending}), an alternative description is as a cell
    $$\varepsilon : \vec p \To E(e, e')$$
    Such that
    \[\begin{tikzcd}[sep=large]
        \ar[d, equal] A & \ar[l, loose, "{E(j, e)}"', ""{anchor=center,name=t1}] D \ar[d, equal] & \ar[l, loose, "p_1"'] \cdots & \ar[l, loose, "p_n"'] D' \ar[d, equal]\\
        \ar[d, equal] A & \ar[l, loose, "{E(t, e)}"', ""{anchor=center,name=m1}] D &{}& \ar[ll, loose, "{E(e, e')}"] D' \ar[d, equal]\\
        A &&& \ar[lll, loose, "{E(t, e')}", ""{anchor=center,name=b}] D
        \ar[from=2-1, to=2-4, phantom, ""{anchor=center,name=m2}]
        \ar[from=t1, to=m1, phantom, "\rtimes"{anchor=center}] \ar[from=1-3, to=2-3, phantom, "\varepsilon"{anchor=center}]
        \ar[from=m2, to=b, phantom, "{\smallsmile_e(t, e')}"{anchor=center}]
    \end{tikzcd}
    \hspace{1em}=\hspace{1em}
    \begin{tikzcd}[sep=large]
        \ar[d, equal] A & \ar[l, loose, "{E(j, e)}"', ""{anchor=center,name=t1}] D \ar[d, equal] & \ar[l, loose, "p_1"'] \cdots & \ar[l, loose, "p_n"'] D' \ar[d, equal]\\
        \ar[d, equal] A & \ar[l, loose, "{E(t, e)}"', ""{anchor=center,name=m1}] D &{}& \ar[ll, loose, "{E(e, e')}"] D' \ar[d, equal]\\
        \ar[d, equal] A &&& \ar[lll, loose, "{E(j, e')}", ""{anchor=center,name=b}] D \ar[d, equal]\\
        A &&& \ar[lll, loose, "{E(t, e')}", ""{anchor=center,name=bb}]
        \ar[from=2-1, to=2-4, phantom, ""{anchor=center,name=m2}]
        \ar[from=t1, to=m1, phantom, "\rtimes"{anchor=center}] \ar[from=1-3, to=2-3, phantom, "\varepsilon"{anchor=center}]
        \ar[from=m2, to=b, phantom, "{\smallsmile_e(j, e')}"{anchor=center}]
        \ar[from=b, to=bb, phantom, "\rtimes'"{anchor=center}]
    \end{tikzcd}\]
\end{remark}

\begin{example}
    The extension operator of a $T$-algebra $(e, \rtimes)$ is a $E(j, e)$-graded $T$-algebra morphism from $(t, \dag)$ to $(e, \rtimes)$.
\end{example}

We now axiomatize the category of algebras. Note that the following definition is stronger than the 2-categorical one, using generalized objects.

\begin{definition}
    \label{def:alg_obj}
    Let $T$ be a relative monad. A \textbf{$T$-algebra} $(u_T : \Alg(T) \to E, \rtimes_T)$ is an \textbf{algebra object} for $T$ when
    \begin{enumerate}
        \item For every $T$-algebra $(e : D \to E, \rtimes)$ there is a unique arrow $\unit_{(e, \rtimes)} : D \to \Alg(T)$ such that $u_T\unit_{(e, \rtimes)} = e$ et $\rtimes_T(1, \unit_{(e, \rtimes)}) = \rtimes$.
        \item For every graded $T$-algebra morphism $\varepsilon : E(1, e), \vec p \To E(1, e')$ there is a unique cell $\unit_\varepsilon : \Alg(T)(1, \unit_{(e, \rtimes)}), \vec p \To \Alg(1, \unit_{(e', \rtimes')})$ such that $\varepsilon = E(1, u_T), \unit_\varepsilon$.
    \end{enumerate}
\end{definition}

Note that in the second condition we are using the fact that $u_T\unit_{(e, \rtimes)} = e$ from the first condition. Since $\rtimes_T(1, \unit_{(e, \rtimes)}) = \rtimes$, the second condition implies that there is a bijection between graded $T$-algebra morphisms $E(1, e), \vec p \To E(1, e')$ and cells $\Alg(T)(1, \unit_{(e, \rtimes)}), \vec p \To \Alg(T)(1, \unit_{(e', \rtimes')})$

\begin{remark}
    When $T$ has an algebra object, there is an arrow $f_T := \unit_{(t, \dag)} : A \to \Alg(T)$ induced by the $T$-algebra $(t, \dag)$.
\end{remark}

Algebra objects, as intended, induce resolutions:

\begin{lemma}
    \label{lem:rmnd_alg_obj_res}
    Let $T$ be a relative monad admitting an algebra object. Then $T$ admits a resolution.
    \[\begin{tikzcd}
        & \Alg(T) \ar[dr, "u_T"] &\\
        A \ar[ur, "f_T"] \ar[rr, "j"'] &{}& E
        \ar[from=1-2, to=2-2, phantom, "\dashv"{anchor=center}]
    \end{tikzcd}\]
\end{lemma}

\begin{proof}
    We denote $\eta : j \To u_Tf_T$ the unit of $T$ and note that $\rtimes_T$ induces a cell $\unit_\rtimes : E(j, u_T) \To \Alg(T)(f_T, 1)$. We prove that the data of those two forms a $j$-adjunction.
    \[
    \begin{tikzcd}
        \ar[d, equal] A & \ar[l, equal, ""{anchor=center,name=t1}] A \ar[d, "f_T"] & \ar[l, loose, "{E(j, u_T)}"', ""{anchor=center,name=t2}] \Alg(T) \ar[d, equal]\\
        A & \ar[l, loose, "{E(j, u_T)}", ""{anchor=center,name=b1}] \Alg(T) & \ar[l, loose, equal, ""{anchor=center,name=b2}] \Alg(T)
        \ar[from=t1, to=b1, phantom, "\eta"{anchor=center}] \ar[from=t2, to=b2, phantom, "\unit_\rtimes"{anchor=center}]
    \end{tikzcd}
    \hspace{1em}=\hspace{1em}
    \begin{tikzcd}
        \ar[d, equal] A & \ar[l, equal, ""{anchor=center,name=t1}] A \ar[d, equal] & \ar[l, loose, "{E(j, u_T)}"', ""{anchor=center,name=t2}] \Alg(T) \ar[d, equal] \\
        \ar[d, equal] A & \ar[l, loose, "{E(j, u_Tf_T)}", ""{anchor=center,name=m1}] A & \ar[l, loose, "{E(j, u_T)}"', ""{anchor=center,name=m2}] \Alg(T) \ar[d, equal] \\
        A &{}& \ar[ll, loose, "{E(j, u_T)}"] \Alg(T)
        \ar[from=t1, to=m1, phantom, "\eta"{anchor=center}] \ar[from=t2, to=m2, phantom, "="{anchor=center}]
        \ar[from=2-2, to=3-2, phantom, "{E(j, u_T),\unit_{\rtimes_T}}"]
    \end{tikzcd}
    \]
    $$
    =\hspace{1em}
    \begin{tikzcd}
        \ar[d, equal] A & \ar[l, equal, ""{anchor=center,name=t1}] A \ar[d, equal] & \ar[l, loose, "{E(j, u_T)}"', ""{anchor=center,name=t2}] \Alg(T) \ar[d, equal]\\
        \ar[d, equal] A & \ar[l, loose, "{E(j, u_Tf_T)}", ""{anchor=center,name=m1}] A & \ar[l, loose, "{E(j, u_T)}", ""{anchor=center,name=m2}] \Alg(T) \ar[d, equal]\\
        A &{}& \ar[ll, loose, "{E(j, u_T)}"] \Alg(T)
        \ar[from=t1, to=m1, phantom, "\eta"{anchor=center}] \ar[from=t2, to=m2, phantom, "="{anchor=center}]
        \ar[from=2-2, to=3-2, phantom, "{\rtimes_T(j, 1)}"{anchor=center}]
    \end{tikzcd}
    \hspace{1em}=\hspace{1em}
    \begin{tikzcd}
        \ar[d, equal] A & \ar[l, loose, "{E(j, u_T)}"', ""{anchor=center, name=t}] E \ar[d, equal]\\
        A & \ar[l, loose, "{E(j, f_T)}"', ""{anchor=center, name=b}] E
        \ar[from=t, to=b, phantom, "="{anchor=center}]
    \end{tikzcd}
    $$
    Where the first equality is by bending $f_T$, the second by definition of $\unit_{\rtimes_T}$ and the last is the unit law of $(u_T, \rtimes_T)$. Then, using the definition of $\unit_{\rtimes_T}$ we have
    \[\begin{tikzcd}[sep=large]
        \ar[d, equal] E & \ar[l, loose, "{E(1 ,u_Tf_T)}"', ""{anchor=center,name=t1}] A \ar[d, equal] & \ar[l, equal, ""{anchor=center,name=t2}] A \ar[d, equal]\\
        \ar[d, equal] E & \ar[l, loose, "{E(1, u_Tf_T)}", ""{anchor=center,name=m1}] A & \ar[l, loose, "{E(j, u_Tf_T)}", ""{anchor=center,name=m2}] A \ar[d, equal]\\
        A &{}& \ar[ll, loose, "{E(1, u_Tf_T)}"]
        \ar[from=t1, to=m1, phantom, "="{anchor=center}] \ar[from=t2, to=m2, phantom, "\eta"{anchor=center}]
        \ar[from=2-2, to=3-2, phantom, "{E(1, u_T),\unit_{\rtimes_T}}"{anchor=center}]
    \end{tikzcd}
    \hspace{1em}=\hspace{1em}
    \begin{tikzcd}[sep=large]
        \ar[d, equal] E & \ar[l, loose, "{E(1 ,u_Tf_T)}"', ""{anchor=center,name=t1}] A \ar[d, equal] & \ar[l, equal, ""{anchor=center,name=t2}] A \ar[d, equal]\\
        \ar[d, equal] E & \ar[l, loose, "{E(1, u_Tf_T)}", ""{anchor=center,name=m1}] A & \ar[l, loose, "{E(j, u_Tf_T)}", ""{anchor=center,name=m2}] A \ar[d, equal]\\
        E &{}& \ar[ll, loose, "{E(1, u_Tf_T)}"] A
        \ar[from=t1, to=m1, phantom, "="{anchor=center}] \ar[from=t2, to=m2, phantom, "\eta"{anchor=center}]
        \ar[from=2-2, to=3-2, phantom, "\rtimes_T"{anchor=center}]
    \end{tikzcd}
    \]
    Bending $t$ and using $\rtimes_T(1, f_T) = \dag$, the corresponding cell is
    \[\begin{tikzcd}[sep=large]
        \ar[d, equal] A & \ar[l, equal, ""{anchor=center,name=t}] A \ar[d, equal]\\
        \ar[d, equal] A & \ar[l, loose, "{E(j, j)}", ""{anchor=center,name=mt}] A \ar[d, equal]\\
        \ar[d, equal] A & \ar[l, loose, "{E(j, u_Tf_T)}", ""{anchor=center,name=mb}] A \ar[d, equal]\\
        A & \ar[l, loose, "{E(t, u_Tf_T)}", ""{anchor=center,name=b}] A
        \ar[from=t, to=mt, phantom, "\smallfrown_j"{anchor=center}]
        \ar[from=mt, to=mb, phantom, "{E(j, \eta)}"{anchor=center}]
        \ar[from=mb, to=b, phantom, "{\rtimes_T(1, f_T)}"{anchor=center}]
    \end{tikzcd}
    \hspace{1em}=\hspace{1em}
    \begin{tikzcd}[sep=large]
        \ar[d, equal] A & \ar[l, equal, ""{anchor=center,name=t}] A \ar[d, equal]\\
        \ar[d, equal] A & \ar[l, loose, "{E(j, j)}", ""{anchor=center,name=mt}] A \ar[d, equal]\\
        \ar[d, equal] A & \ar[l, loose, "{E(j, t)}", ""{anchor=center,name=mb}] A \ar[d, equal]\\
        A & \ar[l, loose, "{E(t, t)}", ""{anchor=center, name=b}] A
        \ar[from=t, to=mt, phantom, "\smallfrown_j"{anchor=center}]
        \ar[from=mt, to=mb, phantom, "{E(j, \eta)}"{anchor=center}]
        \ar[from=mb, to=b, phantom, "\dag"{anchor=center}]
    \end{tikzcd}
    \hspace{1em}=\hspace{1em}
    \begin{tikzcd}
        \ar[d, equal]A & \ar[l, equal, ""{anchor=center,name=t}] A \ar[d, equal]\\
        A & \ar[l, loose, "{E(t, t)}", ""{anchor=center,name=b}] A
        \ar[from=t, to=b, phantom, "\smallfrown_t"{anchor=center}]
    \end{tikzcd}\]
    Bending $t$ once again, we get $1_{E(1, u_Tf_T)}$. Since $1_T$ is an algebra endomorphism on $(t, \dag)$, by uniqueness the cell above s induced by $1_{\Alg(T)}$ by postcomposing $1_{E(1, u_T)}$; hence $f_T\dashv_ju_T$. Finally, the operator is given by
    \[\begin{tikzcd}[sep=huge]
        \ar[d, equal] A & \ar[l, loose, "{E(j, u_Tf_T)}"', ""{anchor=center,name=t}] A \ar[d, equal]\\
        \ar[d, equal] A & \ar[l, loose, "{\Alg(T)(f_T, f_T)}", ""{anchor=center,name=m}] A \ar[d, equal]\\
        A & \ar[l, loose, "{E(u_Tf_T, u_Tf_T)}", ""{anchor=center,name=b}] A
        \ar[from=t, to=m, phantom, "\unit_{\rtimes_T}"{anchor=center}]
        \ar[from=m, to=b, phantom, "{\smallfrown_{u_T}(f_T, f_T)}"{anchor=center}]
    \end{tikzcd}
    \hspace{1em}=\hspace{1em}
    \begin{tikzcd}[sep=huge]
        \ar[d, equal] A & \ar[l, equal] A \ar[d, "f_T"] & \ar[l, equal] A \ar[d, equal] & \ar[l, loose, "{E(j, u_Tf_T)}"', ""{anchor=center,name=t}] A \ar[d, equal]\\
        \ar[d, equal] A & \ar[l, loose, "{\Alg(T)(f_T, 1)}", ""{anchor=center,name=m1}] \Alg(T) \ar[d, equal] & \ar[l, loose, "{\Alg(T)(1, f_T)}"] A & \ar[l, loose, "{E(j, u_Tf_T)}", ""{anchor=center,name=m2}] A \ar[d, equal]\\
        A & \ar[l, loose, "{E(u_Tf_T, u_T)}", ""{anchor=center,name=b}] \Alg(T) &{}& \ar[ll, loose, "{\Alg(T)(1, f_T)}"] A
        \ar[from=t, to=m2, phantom, "="{anchor=center}]
        \ar[from=m1, to=b, phantom, "{\smallfrown_{u_T}(f_T, 1)}"] \ar[from=2-3, to=3-3, phantom, "{\unit_{\rtimes_T}(1, f_T)}"{anchor=center}]
    \end{tikzcd}\]
    Using the definition of $\unit_{\rtimes_T}$ and $\rtimes_T(1, f_T) = \dag$, we get
    \[
        \begin{tikzcd}[sep=huge]
        \ar[d, equal] A & \ar[l, equal] A \ar[d, "f_T"] & \ar[l, equal] A \ar[d, equal] & \ar[l, loose, "{E(j, u_Tf_T)}"', ""{anchor=center,name=t}] A \ar[d, equal]\\
        \ar[d, equal] A & \ar[l, loose, "{\Alg(T)(f_T, 1)}"] \Alg(T) \ar[d, "u_T"] & \ar[l, loose, "{\Alg(T)(1, f_T)}"] A & \ar[l, loose, "{E(j, u_Tf_T)}", ""{anchor=center,name=m2}] A \ar[d, equal]\\
        A & \ar[l, loose, "{E(u_Tf_T, 1)}", ""{anchor=center,name=b}] E &{}& \ar[ll, loose, "{\Alg(T)(1, u_Tf_T)}"] A
        \ar[from=t, to=m2, phantom, "="{anchor=center}]
        \ar[from=2-3, to=3-3, phantom, "{\rtimes_T(1, f_T)}"{anchor=center}]
    \end{tikzcd}
    \hspace{1em}=\hspace{1em} 
    \begin{tikzcd}
        \ar[d, equal] A & \ar[l, loose, "{E(j, t)}"', ""{anchor=center,name=t}] A \ar[d, equal]\\
        A & \ar[l, loose, "{E(t, t)}", ""{anchor=center,name=b}] A
        \ar[from=t, to=b, phantom, "\dag"{anchor=center}]
    \end{tikzcd}\]
    Therefore, we have a resolution $f_T\dashv_ju_T$ of $T$.
\end{proof}

We now state the main result of the subsection.

\begin{theorem}
    \label{th:partial_adj_radj_rmnd}
    Let $j : A \to E$ an arrow. The functor $\obslash : \cat{RAdj}_L(j) \to \cat{RMnd}(j)^\op$ admits a partial right adjoint section defined on the $j$-monads admitting algebra objects.
    $$\begin{tikzcd}[sep=huge]
        \cat{RAdj}_L(j) \ar[r, bend left=15, "\obslash_j", ""{anchor=center, name=t}] & \ar[l, bend left=15, dashed, "{f_{(-)}\dashv_ju_{(-)}}", ""{anchor=center,name=b}] \cat{RMnd}(j)^\op
        \ar[from=t, to=b, adj]
    \end{tikzcd}$$
    And a left morphism is strict if and only if its transpose is the identity $j$-monad morphism.
\end{theorem}

\begin{proof}
    It suffices to construct a section on objects and on homsets. Using \ref{lem:rmnd_alg_obj_res}, we get a partial assignment on objects. Let $T$ and $T'$ be $j$-monads such that $T$ admits a resolution $\ell\dashv_jr$ and $T'$ admits an algebra algebra object. We construct an inverse to
    $$(\obslash_j)_{(\ell\dashv_jr), (f_{T'}\dashv_ju_{T'})} : \cat{RAdj}_L(j)(\ell\dashv_jr, f_{T'}\dashv_ju_{T'}) \to \cat{RMnd}(j)(T', T)$$
    Since $r$ and $t$ form $T$-algebras, a $j$-monad morphism $\tau : T' \to T$ induces a $T$-algebra structure on each. The universal property of %FIXME: finir
\end{proof}

\begin{corollary}
    Let $j : A \to E$ an arrow. The partial functor $u_{(-)} : \cat{RMnd}(j)^\op \to \underline{\bb X}/E$ is fully faithful.
\end{corollary}

\begin{proof}
    Apply Lemma \ref{lem:rradj_ff} and Theorem \ref{th:partial_adj_radj_rmnd} and compose the functors.
\end{proof}

\begin{corollary}
    \label{cor:rmnd_res_term}
    For $T$ a $j$-relative monad admitting an algebra object, the resolution $f_T\dashv_ju_T$ is terminal in $\cat{Res}(T)$.
\end{corollary}

\begin{proof}
    Let $\ell\dashv_jr$ be another resolution of $T$. By the Theorem \ref{th:partial_adj_radj_rmnd}, strict morphisms from $\ell\dashv_jr$ to $f_T\dashv_ju_T$ correspond to the identity morphism on $T$, hence are unique.
\end{proof}

Finally, algebra objects for the trivial $j$-monad always exist, since they are trivial:

\begin{proposition}
    \label{prop:triv_rmnd_alg}
    Let $j : A \to E$ an arrow, then $(\id_E, 1_{E(j, 1)})$ exhibits an algebra object for the trivial $j$-monad.
\end{proposition}

\begin{proof}
    Let $(e : D \to E, \rtimes)$ a $j$⁻algebra. $e$ exhibits a unique arrow $\unit_{(e, \rtimes)} : D \to E$. Moreover, every graded algebra morphism factors through the unity on $E$.
\end{proof}

We shift to opalgebra objects and introduce the according notions.

\begin{definition}
    \index{$T$-opalgebra!graded morphism} \index{$T$-opalgebra!ungraded morphism}
    Let $(a : A \to B, \ltimes)$ and $(a : A \to B', \ltimes')$ be $T$-opalgebras. A \textbf{$\vec p$-graded $T$-opalgebra morphism} from $(a, \ltimes)$ to $(a', \ltimes')$ is a cell
    $$\alpha : \vec p, B(1, a) \To B'(1, a')$$
    Such that
    \[\begin{tikzcd}
        \ar[d, equal] B' & \ar[l, loose, "p_1"'] \cdots & \ar[l, loose, "p_n"'] B \ar[d, equal] & \ar[l, loose, "{B(1, a)}"'] A & \ar[l, loose, "{E(j, t)}"'] A \ar[d, equal] \\
        \ar[d, equal] & \ar[l, loose, "p_1"'] \cdots & \ar[l, loose, "p_n"'] B &{}& \ar[ll, loose, "{B(1, a)}"] A \ar[d, equal]\\
        B' &&&& \ar[llll, loose, "{B'(1, a')}"]
        \ar[from=1-2, to=2-2, phantom, "="{anchor=center}] \ar[from=1-4, to=2-4, phantom, "\rho"{anchor=center}]
        \ar[from=2-1, to=2-5, phantom, ""{anchor=center,name=m}] \ar[from=3-1, to=3-5, phantom, ""{anchor=center,name=b}]
        \ar[from=m, to=b, phantom, "\alpha"{anchor=center}]
    \end{tikzcd}
    \hspace{1em}=\hspace{1em}
    \begin{tikzcd}
        \ar[d, equal] B' & \ar[l, loose, "p_1"'] \cdots & \ar[l, loose, "p_n"'] B & \ar[l, loose, "{B(1, a)}"'] B \ar[d, equal] & \ar[l, loose, "{E(j, t)}"', ""{anchor=center,name=t2}] A \ar[d, equal] \\
        \ar[d, equal] B' &&& \ar[lll, loose, "{B'(1, a')}"] A & \ar[l, loose, "{E(j, t)}", ""{anchor=center,name=m2}] A \ar[d, equal]\\
        B' &&&& \ar[llll, loose, "{B'(1, a')}"]
        \ar[from=1-1, to=1-4, phantom, ""{anchor=center,name=t1}] \ar[from=2-1, to=2-4, phantom, ""{anchor=center,name=m1}]
        \ar[from=2-1, to=2-5, phantom, ""{anchor=center,name=m3}] \ar[from=3-1, to=3-5, phantom, ""{anchor=center,name=b}]
        \ar[from=t1, to=m1, phantom, "\alpha"{anchor=center}] \ar[from=t2, to=m2, phantom, "="{anchor=center}]
        \ar[from=m3, to=b, phantom, "\rho'"{anchor=center}]
    \end{tikzcd}\]
    Where $\rho$ is as in Remark \ref{rem:ropalg_act}. When $n=0$, such a morphism is \textbf{ungraded}.
\end{definition}

\begin{example}
    As for $T$-algebras, the extension operator $\ltimes : E(j, t) \To B(a, a)$ is a $(B(1, a), E(j, 1))$-graded $T$-opalgebra morphism from $(t, \dag)$ to $(a, \ltimes)$.
\end{example}

\begin{definition}
    \label{def:opalg_obj}
    Let $T$ be a relative monad. A \textbf{$T$-opalgebra} $(k_T : A \to \Opalg(T), \ltimes_T)$ is an \textbf{opalgebra object} for $T$ when
    \begin{enumerate}
        \item For every $T$-opalgebra $(a : A \to B, \ltimes)$ there is a unique arrow $[]_{(a, \ltimes)} : \Opalg(T) \to B$ such that $[]_{(a, \ltimes)}k_T= a$ et $[]_{(a, \ltimes)}\ltimes_T = \ltimes$.
        \item For every graded $T$-opalgebra morphism $\alpha : \vec p, B(1, a) \To B'(1, a')$ there is a unique cell $[]_\alpha : \vec p, \Opalg(T)(1, []_{(a, \ltimes)}) \To B'(1, []_{(a', \ltimes')})$ such that $\alpha = []_\alpha(1, k_T)$.
    \end{enumerate}
\end{definition}

\begin{remark}
    When $T$ has an opalgebra object, there is an arrow $v_T := []_{(t, \dag)} : \Opalg(T) \to E$ induced by the $T$-opalgebra $(t, \dag)$.
\end{remark}

Opalgebras also induce resolutions:

\begin{lemma}
    Let $T$ be a relative monad admitting an opalgebra object. Then $T$ admits a resolution.
    \[\begin{tikzcd}
        & \Opalg(T) \ar[dr, "v_T"] &\\
        A \ar[ur, "k_T"] \ar[rr, "j"'] &{}& E
        \ar[from=1-2, to=2-2, phantom, "\dashv"{anchor=center}]
    \end{tikzcd}\]
\end{lemma}

\begin{proof}
    We denote $\eta : j \To v_Tk_T$ the unit of $T$ and note that $\ltimes$ induces a cell $[]_{\ltimes_T} : \Opalg(T)(1, k_t), E(j, v_T)\To \Opalg(T)(1, 1)$. We prove that the data of those two forms a $j$-adjunction. First, using respectively the definition of $[]_{\ltimes_T}$, bending $v_T$, $v_t\ltimes_T=\dag$ and the unit law for $T$, we have:
    \[\begin{tikzcd}[sep=large]
        \ar[d, equal] A & \ar[l, equal, ""{anchor=center,name=t1}] A \ar[d, "k_T"'] & \ar[l, loose, "{E(j, v_Tk_T)}"', ""{anchor=center,name=t2}] A \ar[d, equal]\\
        A & \ar[l, loose, "{E(j, v_T)}", ""{anchor=center,name=b1}] \Opalg(T) & \ar[l, loose, "{\Opalg(T)(1, k_T)}", ""{anchor=center,name=b2}] A
        \ar[from=t1, to=b1, phantom, "\eta"{anchor=center}] \ar[from=t2, to=b2, phantom, "{[]_{\ltimes_T}(1, k_T)}"{anchor=center}]
    \end{tikzcd}
    \hspace{1em}=\hspace{1em}
    \begin{tikzcd}[sep=large]
        \ar[d, equal] A & \ar[l, equal, ""{anchor=center,name=t1}] A \ar[d, "k_T"'] & \ar[l, loose, "{E(j, v_Tk_T)}"', ""{anchor=center,name=t2}] A \ar[d, equal]\\
        A & \ar[l, loose, "{E(j, v_T)}", ""{anchor=center,name=b1}] \Opalg(T) & \ar[l, loose, "{\Opalg(T)(1, k_T)}", ""{anchor=center,name=b2}] A
        \ar[from=t1, to=b1, phantom, "\eta"{anchor=center}] \ar[from=t2, to=b2, phantom, "\ltimes_T"{anchor=center}]
    \end{tikzcd} \]
    \[\begin{tikzcd}[sep=huge]
        \ar[d, equal] A & \ar[l, loose, "{E(j, t)}"', ""{anchor=center, name=t}] A \ar[d, equal]\\
        \ar[d, equal] A & \ar[l, loose, "{\Opalg(T)(k_T, k_T)}"', ""{anchor=center, name=mt}] A \ar[d, equal]\\
        \ar[d, equal] A & \ar[l, loose, "{E(t, t)}", ""{anchor=center,name=mb}] A \ar[d, equal]\\
        A & \ar[l, loose, "{E(j, t)}", ""{anchor=center,name=b}]
        \ar[from=t, to=mt, phantom, "\ltimes_T"{anchor=center}]
        \ar[from=mt, to=mb, phantom, "{\smallfrown_{v_T}(k_T, k_T)}"{anchor=center}]
        \ar[from=mb, to=b, phantom, "{E(\eta, 1)}"{anchor=center}]
    \end{tikzcd}
    \hspace{1em}=\hspace{1em}
    \begin{tikzcd}[sep=huge]
        \ar[d, equal] A & \ar[l, loose, "{E(j, t)}"', ""{anchor=center,name=t}] A \ar[d, equal]\\
        \ar[d, equal] A & \ar[l, loose, "{E(t, t)}", ""{anchor=center,name=m}] A \ar[d, equal]\\
        A & \ar[l, loose, "{E(j, t)}", ""{anchor=center,name=b}] A
        \ar[from=t, to=m, phantom, "\dag"{anchor=center}]
        \ar[from=m, to=b, phantom, "{E(\eta, 1)}"{anchor=center}]
    \end{tikzcd}
    \hspace{1em}=\hspace{1em}
    \begin{tikzcd}[sep=huge]
        \ar[d, equal] A & \ar[l, loose, "{E(j, t)}"', ""{anchor=center,name=t}] A \ar[d, equal]\\
        A & \ar[l, loose, "{E(j, t)}", ""{anchor=center,name=b}] A
        \ar[from=t, to=b, phantom, "="{anchor=center}]
    \end{tikzcd}\]
    Then, bending $k_T$, using the definition of $[]_{\ltimes_T}$ and the unit law for the opalgebra $(k_T, \ltimes_T)$, we have:
    \[\begin{tikzcd}[sep=huge]
        \ar[d, equal] A & \ar[l, loose, "{\Opalg(T)(k_T, 1)}"', ""{anchor=center,name=t}] \Opalg(T) \ar[d, equal]\\
        \ar[d, equal] A & \ar[l, loose, "{E(j, v_T)}"', ""{anchor=center,name=m}] \Opalg(T) \ar[d, equal]\\
        A & \ar[l, loose, "{\Opalg(T)(k_T, 1)}", ""{anchor=center,name=b}] \Opalg(T)
        \ar[from=t, to=m, phantom, "\eta"{anchor=center}]
        \ar[from=m, to=b, phantom, "{[]_{\ltimes_T}}"{anchor=center}]
    \end{tikzcd}
    \hspace{1em}=\hspace{1em}
    \begin{tikzcd}[sep=huge]
        \ar[d, equal] A & \ar[l, equal] A \ar[d, "k_T"'] & \ar[l, equal] A \ar[d, equal] & \ar[l, loose, "{\Opalg(T)(k_T, 1)}"', ""{anchor=center,name=t3}] \Opalg(T) \ar[dd, equal] \\
        \ar[d, equal] A & \ar[l, loose, "{\Opalg(T)(k_T, 1)}"', ""{anchor=center,name=mt1}] \Opalg(T) \ar[d, equal] & \ar[l, loose, "{\Opalg(T)(1, k_T)}"', ""{anchor=center,name=mt2}] A \ar[d, equal] & \\
        \ar[d, equal] A & \ar[l, loose, "{E(j, v_T)}"', ""{anchor=center,name=mb1}] \Opalg(T) \ar[d, equal] & \ar[l, loose, "{\Opalg(T)(1, k_T)}", ""{anchor=center,name=mb2}] A & \ar[l, loose, "{\Opalg(T)(k_T, 1)}", ""{anchor=center,name=mb3}] \Opalg(T) \ar[d, equal] &&\\
        A & \ar[l, loose, "{\Opalg(T)(k_T, 1)}", ""{anchor=center,name=b1}] \Opalg(T) && \ar[ll, loose, equal, ""{anchor=center, name=b2}] \Opalg(T)
        \ar[from=t3, to=mb3, phantom, "="{anchor=center}]
        \ar[from=mt1, to=mb1, phantom, "\eta"{anchor=center}] \ar[from=mt2, to=mb2, phantom, "="{anchor=center}]
        \ar[from=mb1, to=b1, phantom, "{[]_{\ltimes_T}}"{anchor=center}] \ar[from=3-3, to=b2, phantom, "\smallsmile_{k_T}"{anchor=center}]
    \end{tikzcd}\]
    \[=\hspace{1em}
    \begin{tikzcd}[sep=huge]
        \ar[d, equal] A & \ar[l, equal, ""{anchor=center,name=t1}] A \ar[d, equal] & \ar[l, loose, "{\Opalg(T)(k_T, 1)}"', ""{anchor=center,name=t2}] \Opalg(T) \ar[dd, equal]\\
        \ar[d, equal] A & \ar[l, loose, "{E(j, t)}", ""{anchor=center,name=mt}] A \ar[d, equal] & \\
        \ar[d, equal] A & \ar[l, loose, "{\Opalg(T)(k_T, k_T)}"', ""{anchor=center,name=mb1}] A & \ar[l, loose, "{\Opalg(T)(k_T, 1)}"', ""{anchor=center,name=mb2}] \Opalg(T) \ar[d, equal]\\
        A && \ar[ll, loose, "{\Opalg(T)(k_T, 1)}", ""{anchor=center,name=b}] \Opalg(T)
        \ar[from=t1, to=mt, phantom, "\eta"{anchor=center}] \ar[from=t2, to=mb2, phantom, "="{anchor=center}]
        \ar[from=mt, to=mb1, phantom, "\ltimes_T"{anchor=center}]
        \ar[from=3-2, to=b, phantom, "{\smallsmile_{k_T}(k_T, 1)}"{anchor=center}]
    \end{tikzcd}
    \hspace{1em}=\hspace{1em}
    \begin{tikzcd}[sep=huge]
        \ar[d, equal] A & \ar[l, loose, "{\Opalg(T)(k_T, 1)}"', ""{anchor=center,name=t}] \Opalg(T) \ar[d, equal]\\
        A & \ar[l, loose, "{\Opalg(T)(k_T, 1)}", ""{anchor=center,name=b}] \Opalg(T)
        \ar[from=t, to=b, phantom, "="{anchor=center}]
    \end{tikzcd}\]
    Hence $k_T \dashv_j v_T$. The operator is given by
    \[\begin{tikzcd}[sep=large]
        \ar[d, equal] A & \ar[l, loose, "{E(j, v_T)}"', ""{anchor=center,name=t1}] \Opalg(T) \ar[d, equal] & \ar[l, loose, "{\Opalg(T)(1, k_T)}"', ""{anchor=center,name=t2}] A\ar[d, equal] \\
        A & \ar[l, loose, "{\Opalg(T)(k_T, 1)}", ""{anchor=center,name=b1}] \Opalg(T) & \ar[l, loose, "{E(v_T, v_Tk_T)}", ""{anchor=center,name=b2}] A
        \ar[from=t1, to=b1, phantom, "{[]_{\ltimes_T}}"{anchor=center}] \ar[from=t2, to=b2, phantom, "{\smallfrown_{v_T}(1, k_T)}"{anchor=center}]
    \end{tikzcd}
    \hspace{1em}=\hspace{1em}
    \begin{tikzcd}[sep=huge]
        \ar[d, equal] A & \ar[l, loose, "{E(j, t)}"', ""{anchor=center,name=t}] A \ar[d, equal]\\
        \ar[d, equal] A & \ar[l, loose, "{\Opalg(T)(k_T, v_T)}"', ""{anchor=center,name=m}] A \ar[d, equal]\\
        A & \ar[l, loose, "{E(v_Tk_T, v_Tk_T)}", ""{anchor=center,name=b}] A
        \ar[from=t, to=m, phantom, "\ltimes_T"{anchor=center}] \ar[from=m, to=b, phantom, "{\smallfrown_{v_T}(k_T, k_T)}"{anchor=center}]
    \end{tikzcd}
    \hspace{1em}=\hspace{1em}
    \begin{tikzcd}[sep=huge]
        \ar[d, equal] A & \ar[l, loose, "{E(j, t)}"', ""{anchor=center,name=t}] A \ar[d, equal]\\
        A & \ar[l, loose, "{E(t, t)}", ""{anchor=center,name=b}] A
        \ar[from=t, to=b, phantom, "="{anchor=center}]
    \end{tikzcd}\]
    Therefore, $k_T\dashv_jv_T$ is a resolution of $T$.
\end{proof}

\begin{corollary}
    Let $T$ be a relative monad admitting an opalgebra object $(k_T, \ltimes_T)$. Then $\ltimes_T$ is invertible.
\end{corollary}

\begin{proof}
    Since $k_T\dashv_jv_T$, $\Opalg(T)(k_T, k_T) \simeqq E(j, v_Tk_T) = E(j, t)$ and by construction this cell is $\ltimes_T$.
\end{proof}

\begin{theorem}
    \label{th:partial_radj_adj_rmnd}
    $\oslash_j : \cat{RAdj}_R(j) \to \cat{Rmnd}(j)$ admits a partial left adjoint section defined the $j$-monads admitting opalgebra objects
    \[\begin{tikzcd}[sep=huge]
        \cat{RAdj}_R(j) \ar[r, bend right=15, "\oslash_j"', ""{anchor=center,name=b}] & \ar[l, dashed, bend right=15, "{k_{(-)}\dashv_jv_{(-)}}"', ""{anchor=center, name=t}] \cat{RMnd}(j)^\op
        \ar[from=t, to=b, adj]
    \end{tikzcd}\]
    And a right morphism is strict if and only if its transpose is the identity $j$-monad morphism.
\end{theorem}

\begin{proof}
    %TODO
\end{proof}

\begin{corollary}
    The partial functor $k_{(-)} : \cat{RMnd}(j) \to A/\underline{\bb X}$ is fully faithful.
\end{corollary}

\begin{proof}
    Apply Lemma \ref{lem:lradj_ff} and Theorem \ref{th:partial_radj_adj_rmnd} and compose the functors.
\end{proof}

\begin{corollary}
    \label{cor:rmnd_res_init}
    For $T$ a relative monad admitting an opalgebra object, the resolution $k_T \dashv v_T$ is initial in $\cat{Res}(T)$.
\end{corollary}

\begin{proof}
    Let $\ell\dashv_jr$ be another resolution of $T$. By the Theorem \ref{th:partial_radj_adj_rmnd}, strict morphisms from $k_T\dashv_jv_T$ to $\ell\dashv_jr$ correspond to the identity morphism on $T$, hence are unique.
\end{proof}

Combining Corollaries \ref{cor:rmnd_res_termi} and \ref{cor:rmnd_res_init}, we get

\begin{corollary}
    For $T$ a relative monad admitting both an algebra and an opalgebra object, there is a unique arrow $i_T : \Opalg(T) \to \Alg(T)$ rendering the following diagrams commutative:
    \[\begin{tikzcd}
        {\Opalg(T)} \ar[r, "i_T"] & {\Alg(T)}\\
        A \ar[u, "k_T"] \ar[ur, "f_T"']
    \end{tikzcd}
    \hspace{3em}
    \begin{tikzcd}
        {\Opalg(T)} \ar[r, "i_T"] \ar[dr, "u_T"'] & {\Alg(T)} \ar[d, "u_T"]\\
                                                  & E
    \end{tikzcd}\]
\end{corollary}

\begin{proof}
    The arrow is $[]_{f_T\dashv_ju_T}$ or equivalently $\unit_{k_T\dashv_jv_T}$.
\end{proof}

As for algebra, we can describe opalgebra objects for the trivial relative monad... when the root is faithful.

\begin{proposition}
    \label{prop:triv_rmnd_opalg}
    For $j : A\to E$ a fully faithful arrow, $(1_A, \ltimes_j := \smallfrown_j^{-1})$ forms an opalgebra object for the trivial $j$-monad.
\end{proposition}

\begin{proof}
    By Lemma \ref{lem:opalg_lmnd_mor}, $(1_A, \ltimes_j)$ forms a $j$-opalgebra. Moreover, it is an opalgebra object: for a $j$-opalgebra $(a, \ltimes)$, define $[]_{(a, \ltimes)} := a$ and for a $\vec p$-graded algebra morphism $\alpha$, define $[]_\alpha := \alpha$. This data satisfies the required properties, thus $(1_A, \ltimes_j)$ is an opalgebra object.
\end{proof}

When the root isn't fully faithful however, the situation is a bit more complicated. And interesting: if an opalgebra object exists, it may be non trivial but its underlying arrow still has several useful "surjectivity" properties. This has yet to be researched, as \cite{arkor2024formaltheoryrelativemonads} notes. The article however states that, in $\Dist$, it is known that an opalgebra object is exactly the \textit{full image factorization} of the root into a bijective-on-objects arrow followed by a fully faithful arrow.

$\Dist$ also has the nice property of being \textit{exact}. We will not dive too deep on exact virtual equipments and refer the interested reader to \cite{arkor2024nervetheoremrelativemonads}. However, it implies that every loose monad is induced by an arrow and in that case we can reduce the search for an opalgebra object to the trivial monad:

\begin{proposition}
    Let $T$ be a $j$-monad such that $E(j, T)$ is induced by $\ell : A\to C$. The isomorphism $E(j, T) \simeqq C(\ell, \ell)$ induces a bijection between $T$-opalgebras and $\ell$-opalgebras and their graded morphisms. $T$ hence admits an opalgebra object if and only if the trivial $\ell$-monad admits one, in which case we have a commutative diagram
    \[\begin{tikzcd}
        \Opalg(\ell) \ar[rr, "\simeqq"] && \Opalg(T)\\
        & \ar[ul, "k_\ell"] A \ar[ur, "k_T"'] &\\
    \end{tikzcd}\]
\end{proposition}

\begin{proof}
    By Lemma \ref{lem:opalg_lmnd_mor}, $T$-opalgebras are arrows $a : A \to B$ equipped with a loose monad morphism $E(j, T) \To B(a, a)$ while $\ell$-opalgebras are arrows $a : A \to B$ with $C(\ell, \ell) \To B(a, a)$. Therefore, the isomorphism $E(j, T) \simeqq C(\ell, \ell)$ induces a bijection between opalgebras and their graded morphisms and the opalgebras object satisfy the same universal property, hence are isomorphic.
\end{proof}

\subsection{Idempotent relative monads}


\subsection{Relative monadicity}
\label{subsec:rel_monadicity}

In this subsection, we follow \cite{arkor2025relativemonadicity} and establish a monadicity theorem for relative monads with dense root. Since one of the hypothesis is the creation of pointwise extension, we have a strict and a non-strict version.







\section{(Co)presheaves objects and (co)completions}
\label{sec:presheaves}

\subsection{Presheaves in a virtual equipment}

We fix a class of horizontal arrows $\bb{P}$ in a virtual equipment $\bb{X}$.

\begin{definition}
    \index{presheaf!object}
    Let $A$ be an object of $\bb{X}$. A \textbf{$\PP$-presheaf object} for $A$ comprises of an object $\psh\PP A$ and a horizontal arrow $\pi_A^\PP : A \dfrom \psh\PP A$ such that:
    \begin{enumerate}
        \item $\pi_A^\PP : A \dfrom \psh\PP A$ is dense.
        \item For every horizontal arrow $p : A \dfrom X$ in $\PP$, there is a unique vertical arrow $\widearc{p} : X \to \psh\PP A$ such that $p = \pi_A^\PP(1, \widearc{p})$.
        \item For every vertical arrow $f : X \to \psh\PP A$, the restriction $\pi_A^\PP(1, f) : A \dfrom X$ is in $\PP$.
    \end{enumerate}
\end{definition}

This definition axiomatizes well the classification of distributors by a presheaf object, and we think of $\pi_A^\PP$ as the evaluation functor. We also note that the existence of a $\PP$-presheaf object for $A$ only depends on the subclass $\PP|_{A\dfrom}$ of $\PP$ of vertical arrows with codomain $A$, which we can also characterize if $\psh\PP A$ exists:

$$\PP|_{A\dfrom} = \set{\pi_A^\PP(1, f) \mid \cod(f) = \psh\PP A}$$

\begin{remark}
    \label{rem:pobj_deter}
    An immediate consequence of this equality is that, for two classes of horizontal arrows $\PP, \PP'$ such that $\pi_A : A \dfrom P$ exhibits both a $\PP$-presheaf and a $\PP'$-presheaf object, we have $\PP|_{A\dfrom} = \PP'|_{A\dfrom}$.
\end{remark}

We also have, since we suppose that $\bb{X}$ is strict, $\widearc{\pi_A^\PP} = 1_{\psh\PP A}$ and, for every horizontal arrow $p : A \dfrom X$ in $\PP$ and $f : Y \to X$, $\widearc{p(1, f)} = \widearc{p}\circ f$.

\begin{proposition}
    \label{prop:rift_psh}
    Let $A$ be an object admitting a $\PP$-presheaf object. For $p : A \dfrom X, q : A \dfrom Y$ in $\PP$, we have $\psh\PP A(\widearc{p}, \widearc{q}) \simeq q \rift p$.
\end{proposition}

\begin{proof}
    By density of $\pi_A^\PP$ and restriction of right lifts:
    \begin{align*}
        \psh\PP A(\widearc{p}, \widearc{q}) &\simeq (\pi_A^\PP \rift \pi_A^\PP)(\widearc{p}, \widearc{q})\\
                                            &\simeq \pi_A^\PP(1, \widearc{q})\rift\pi_A^\PP(1, \widearc{p})\\
                                            &= q \rift p
    \end{align*}
\end{proof}

This proposition has a few useful consequences:

\begin{corollary}
    \label{cor:psh}
    Let $A$ an object admitting a $\PP$-presheaf object and $p : A \dfrom X$ in $\PP$.
    \begin{enumerate}[ref=\ref{cor:psh} (\alph*)]
        \item \label{cor:psh:dense_iff_ff} $\widearc{p} : X \to \psh\PP A$ is fully faithful if and only if $p$ is dense.
        \item \label{cor:psh:fact} Every cell $\phi$ with frame as in the left factorizes uniquely as in the right:
            \[\begin{tikzcd}
                A \ar[d, equal] & \ar[l, "p"', loose] X & \ar[l, "p_1"'{name=p1}, loose] \cdots & \ar [l, "p_n"', loose] X' \ar[d, equal]\\
                A &&& \ar[lll, "q"{name=q}, loose] X'
                \ar[from=p1, to=q, phantom, "\phi"{description}]
            \end{tikzcd}
            \hspace{5em}
            \begin{tikzcd}
                A \ar[d, equal] & \ar[l, "\pi_A"'{name=t}, loose] \psh\PP A \ar[d, equal] & \ar[l, "{\psh\PP A(1,\widearc{p})}"', loose] X & \ar[l, "p_1"'{name=p1}, loose] \cdots & \ar[l, "p_n"', loose] X' \ar[d, equal]\\
                A & \ar[l, "\pi_A"{name=b}, loose] \psh\PP A &&& \ar[lll, "{\psh\PP A(1, \widearc{q})}"{name=rest}, loose] X'
                \ar[from=t, to=b, phantom, "="{anchor=center}] \ar[from=p1, to=rest, phantom, "\widearc{\phi}"{description}]
            \end{tikzcd}\]
    \end{enumerate}
\end{corollary}

We can deduce a second characterization of $\PP$-presheaf object.

\begin{proposition}
    A horizontal arrow $\pi_A^\PP : A \dfrom \psh\PP A$ exhibits a $\PP$-presheaf object for $A$ if and only if for every object $X$, it induces an isomorphism
    $$\pi_A^\PP(1, -) : \bb{X}[X, \psh\PP A] \rightleftarrows \PP\iset{X, A} : \widearc{(-)}$$
    Natural in $X$, between the category of vertical arrows between $X$ and $\psh\PP A$, and the category of horizontal arrows in $\PP$ between $X$ and $A$.
\end{proposition}

\begin{proof}\hfill
    \begin{itemize}
        \item Suppose that $\pi_A^\PP : A \dfrom \psh\PP A$ exhibits a $\PP$-presheaf object. The bijection between objects comes from the universal property of $\PP$-presheaf objects, while the bijection between morphisms comes from the corollary \ref{cor:psh:fact}.
        \item Suppose that $\pi_A^\PP : A \dfrom \psh\PP A$ induces for all $X$ an isomorphism $\bb{X}[X, \psh\PP A] \rightleftarrows \PP\iset{X, A}$. The universal property is induced by the isomorphism and the closure by representable horizontal arrows from the definition of the isomorphism.
    \end{itemize}
\end{proof}

Going back to the definition of a $\PP$-presheaf object, its universal property characterizes it not only up to equivalence, but to isomorphism. This also often useful but is has an important consequence: if we take the closure of $\PP$ under isomorphic distributors, called the \textit{repletion}, we usually don't have an isomorphism between a $\PP$ and a $\PP'$-presheaf object for $A$ (Remark \ref{rem:pobj_deter}).

This doesn't mean $\PP^\simeqq$ isn't useful, it can still be used to create arrows:

\begin{proposition}
    \label{prop:psh_colims}
    Let $A$ be an object admitting a $\PP$-presheaf object, $\vec{p} : X \dfrom \cdots \dfrom Y$ a weight and $q : A \dfrom X$ such that $q\in\PP$ and $q\odot_L \vec{p}$ exists. Then $\psh\PP A$ admits the colimit $\colim^{\vec{p}}\widearc{q}$ if and only if $(q\odot_L\vec{p})\in\PP^\simeqq$, and in that case the colimit is $\pi_A$-absolute.
\end{proposition}

\todo{ajouter \cite[Proposition 3.20]{arkor2026presheavescocompletionsformalcategory} et son corollaire ?}

\begin{proof}
    We have the chain of isomorphisms:
    $$\psh\PP A(\widearc{q}, 1) \rift \vec{p} \simeq(\pi_A\rift q)\rift \vec{p} \simeq \pi_A (q\odot_L\vec{p})$$
    When the colimit exist, we then have $\pi_A(1, \colim^{\vec{p}}\widearc{q}) \simeq q\odot_L\vec{p}\simeq\pi_A(1, \widearc{q})\odot_L\vec{p}$.
\end{proof}

We usually want to consider $\PP$-presheaf objects "containing the representables":

\begin{definition}
    An object $A$ is $\PP$-admissible if it admits a $\PP$-presheaf object and such that $A(1, 1)\in\PP^\simeqq$. We will write
    $$\yo_A^\PP \simeq \widearc{A(1, 1)} : A \to \psh\PP A$$
\end{definition}

The deviation in notation from \cite{arkor2026presheavescocompletionsformalcategory} is to emphasize that this arrow is the analogue of Yoneda embeddings. In fact, it has several similar propreties:

\begin{lemma}
    \label{lem:pp_admisible}
    Let $A$ be a $\PP$-admissible object.
    \begin{itemize}[ref=\ref{lem:pp_admissible} (\alph*)]
        \item \label{lem:pp_admissible:yoneda} (Yoneda Lemma) There's an isomorphism $\psh\PP A(\yo_A^\PP, 1) \simeq \pi_A$.
        \item $\yo_A^\PP$ is dense, fully faithful and creates weighted limits.
        \item For each weight $p\in\PP$, $\widearc{p}$ is the $\yo_A^\PP$-absolute $p$-weighted colimit of $\yo_A^\PP$.
    \end{itemize}
\end{lemma}

\begin{proof}\hfill
    \begin{itemize}
        \item Using \ref{prop:rift_psh}, we have $\psh\PP A(\yo_A^\PP, 1) \simeq \pi_A \rift A(1, 1) \simeq \pi_A$.
        \item Density comes from density of $\pi_A$, fully faithfulness from Corollary \ref{cor:psh:dense_iff_ff} and weighted limits from \ref{todo}
        \item Consequence of \ref{prop:psh_colims}
    \end{itemize}
\end{proof}

A useful consequence of the Yoneda Lemma (\ref{lem:pp_admissible:yoneda}) will be studied in the next section:

\begin{corollary}
    Let $\PP\subseteq\PP'$ two classes of distributors and $A$ a $\PP$ and $\PP'$-admissible object. The induced arrow $\widearc{\pi_A^\PP}$ is fully faithful and commutes with the Yoneda embeddings up to isomorphism:
    \[\begin{tikzcd}
        \psh\PP A \ar[rr, hook, "\widearc{\pi_A^\PP}"] && \psh{\PP'}A\\
                                                       & \ar[lu, "\yo_A^\PP"] A \ar[ru, "\yo_A^{\PP'}"]
    \end{tikzcd}\]
\end{corollary}

\begin{proof}
    Since $\pi_A^\PP$ is dense so the associated arrow is fully faithful by Corollary \ref{cor:psh:dense_iff_ff}. Furthermore, the chain of isomorphisms
    $$\pi_A^{\PP'}(1, \widearc{\pi_A^\PP}\yo_A^\PP) = \pi_A^\PP(1, \yo_A^\PP) \simeq A(1, 1) \simeq \pi_A^{\PP'}(1, \yo_A^{\PP'})$$
    We conclude with \ref{lem:res_dense_fid}.
\end{proof}

\subsection{Copresheaves and Isbell Duality}

We dualize the definition of $\PP$-presheaf objects to get $\PP$-copresheaf objects. One again, we fix a class of horizontal cells $\PP$ in a virtual equipment $\bb{X}$.

\begin{definition}
    \index{copresheaf!object}
    Let $A$ be an object of $\bb{X}$. A $\PP$-copresheaf object for $A$ comprises of an object $\copsh\PP A$ and a horizontal arrow $\ip_A^\PP : \copsh\PP A \dfrom A$. such that:
    \begin{enumerate}
        \item $\ip_A^\PP : \copsh\PP A \dfrom A$ is codense.
        \item For every horizontal arrow $p : X \dfrom A$ de $\PP$, there is a unique vertical arrow $\widecheck{p} : X \to \copsh\PP A$ such that$p = \ip_A^\PP(\widecheck{p}, 1)$.
        \item For every vertical arrow $f : X \to \copsh\PP A$, the restriction $\ip_A^\PP(f, 1) : X \dfrom A$ is in $\PP$.
    \end{enumerate}
\end{definition}

The results from the previous section can be dualized:

\begin{proposition}
    \label{prop:rext_copsh}
    Let $A$ be an object admitting a $\PP$-copresheaf object. For $p : Y \dfrom A, q : Z \dfrom A$ in $\PP$, we have $\copsh\PP A(\widebreve{q}, \widebreve{p}) \simeq p \rext q$.
\end{proposition}

\begin{corollary}
    \label{cor:copsh}
    Let $A$ be an object admitting a $\PP$-copresheaf object.
    \begin{enumerate}[ref=\ref{cor:copsh} (\alph*)]
        \item For $p : A \dfrom X$ in $\PP$, $\widearc{p} : X \to \psh\PP A$ is fully faithful if and only if $p$ is codense.
        \item Every cell $\phi$ with frame as in the left factorizes as in the right:
            \[\begin{tikzcd}
                Z \ar[d, equal] & \ar[l, "p_1"', loose] \cdots & \ar[l, "p_n"'{name=pn}, loose] Y & \ar [l, "p"', loose] A \ar[d, equal]\\
                Z &&& \ar[lll, "q"{name=q}, loose] A
                \ar[from=pn, to=q, phantom, "\phi"{description}]
            \end{tikzcd}
            \hspace{5em}
            \begin{tikzcd}
                Z \ar[d, equal] & \ar[l, "p_1"', loose] \cdots & \ar[l, "p_n"'{name=pn}, loose] Y & \ar[l, "{\copsh\PP A(\widecheck{p}, 1)}"', loose] \copsh\PP A \ar[d, equal] & \ar[l, "\ip_A^\PP"'{name=t}, loose] A \ar[d, equal]\\
                Z &&& \ar[lll, "{\copsh\PP A(\widearc{q}, 1)}"{name=rest}, loose] \copsh\PP A & \ar[l, "\ip_A^\PP"{name=b}, loose] A
                \ar[from=t, to=b, phantom, "="{anchor=center}] \ar[from=pn, to=rest, phantom, "\widearc{\phi}"{description}]
            \end{tikzcd}\]
    \end{enumerate}
\end{corollary}

\begin{proposition}
    A horizontal arrow $\ip_A^\PP : \copsh\PP A \dfrom A$ exhibits a $\PP$-copresheaf object for $A$ if and only if, for every object $X$, it induces an isomorphism:
    $$\ip_A^\PP(-, 1) : \bb{X}[\copsh\PP A, X] \rightleftarrows \PP\iset{A, X} : \widebreve{(-)}$$
    Between the category of vertical arrows between $\copsh\PP A$ and $X$, and the category of horizontal arrows of $\PP$ between $A$ and $X$.
\end{proposition}

\begin{lemma}[coYoneda Lemma]
    There's an isomorphism $\copsh\PP A(1, \oy_A^\PP) \simeq \ip_A$.
\end{lemma}

\todo{\footnotesize À partir d'ici, ce n'est pas dans la littérature sous cette forme. \textbf{À prendre avec des pincettes}}

We can study the relationship between a $\PP$-presheaf and a $\PP$-copresheaf object. Under some hypotheses, we can state Isbell duality.

\begin{theorem}
    \index{Isbell Duality}
    \label{th:isbell_duality}
    Let $A$ be an object $\PP$-admissible and $\PP$-coadmissible. Suppose that $\psh\PP A(1, \yo_A^\PP)$ and $\copsh\PP A(\oy_A^\PP, 1)$ are in $\PP^\cong$. Then there's an adjunction
    %FIXME repletion ?
    \[\begin{tikzcd}
        \ca{O} : \psh\PP A \ar[r, bend left=15, ""{name=O, below}] & \ar[l, bend left=15, ""{name=Spec, above}] \copsh\PP A : \Spec \ar[from=O, to=Spec, adj]
    \end{tikzcd}\]
\end{theorem}

\begin{proof}
    Let $\ca{O} := \widebreve{\psh\PP A(1, \yo_A^\PP)}$ and $\Spec := \widearc{\copsh\PP A(\oy_A^\PP, 1)}$.
    \begin{align*}
        \copsh\PP A(\ca{O}, 1) &= \copsh\PP A(\widebreve{\psh\PP A(1, \yo_A^\PP)}, 1)
                               &\text{(by definition)}\\
                               &\simeq \ip_A^\PP \rext \psh\PP A(1, \yo_A^\PP)
                               &\text{(by \ref{prop:rext_copsh})}\\
                               &\simeq \ip_A^\PP\rext\left(A(1, 1)\rift\pi_A^\PP\right)
                               &\text{(by \ref{prop:rift_psh})}\\
                               &\simeq \left(\ip_A^\PP\rext A(1, 1)\right)\rift\pi_A^\PP
                               &\text{(by \ref{lem:ext_lift:assoc_rift_rext})}\\
                               &\simeq \copsh\PP A(\oy_A^\PP, 1)\rift\pi_A^\PP
                               &\text{(by \ref{prop:rext_copsh})}\\
                               &\simeq \psh\PP A(1, \widearc{\copsh\PP A(\oy_A^\PP, 1)})
                               &\text{(by \ref{prop:rift_psh})}\\
                               &=\psh\PP A(1, \Spec)
                               &\text{(by definition)}
    \end{align*}
\end{proof}

\begin{remark}
    \label{rem:o_spec_plrext}
    The astute reader will see that, hidden in the proof, we also proved that $\ca{O} \simeq \yo_A^\PP \plext \oy_A^\PP$ and $\Spec \simeq \oy_A^\PP \prext \yo_A^\PP$. Another computation gives $\ca{O} \simeq \yo_A^\PP \prext \oy_A^\PP$ and $\Spec \simeq \oy_A^\PP \plext \yo_A^\PP$.
\end{remark}

\begin{definition}
    An object $A$ will be \textbf{Isbell $\PP$-admissible} if it verifies the conditions of Theorem \ref{th:isbell_duality}.
\end{definition}

\begin{proposition}[{\cite[Proposition 2.5]{diliberticodensity}}]
    The Isbell adjunction swaps adjunctions: Let $A$ be a Isbell $\PP$-admissible object and $i : A \hookrightarrow K$ be a dense fully faithful arrow such that $L := \yo_A^\PP \plext i \dashv i \plext \yo_A^\PP =: R$ exists (i.e. $K(i, 1)\in\PP^\simeqq$ and $\yo_A^\PP \plext i$ exists), then
    $$\ca{O}\circ R\dashv L\circ\Spec$$
\end{proposition}

\begin{proof}
    We have the following isomorphisms
    \begin{align*}
        L\circ\Spec &= (\yo_A^\PP \plext i) \circ (\oy_A^\PP \plext \yo_A^\PP)\\
                    &\simeq \oy_A^\PP \plext i\\
                    &\simeq (\ca{O} \circ \yo_A^\PP)\plext i\\
                    &\simeq (\ca{O} \circ R \circ i)\plext i\\
                    &\simeq (\ca{O}\circ R)\plext i \plext i\\
                    &\simeq (\ca{O}\circ R) \plext \id_K
    \end{align*}
\end{proof}

%FIXME: à enlever en compensant au dessus
This adjunction also restricts through inclusions of class of horizontal cells $\PP\subseteq\PP'$:

\begin{theorem}
    Let $A$ a Isbell $\PP$ and $\PP'$-admissible object. The adjunction of Theorem \ref{th:isbell_duality} for $\PP'$ restricts through the one for $\PP$.
\end{theorem}

\begin{proof}
    We consider the following diagram:
    \[\begin{tikzcd}
        \ar[dd, hook, "\widearc{\pi_A^\PP}"'] \psh\PP A \ar[rr, bend left=15, "\ca{O}", ""{anchor=center,name=O}] && \ar[ll, bend left=15, "\Spec", ""{anchor=center,name=Spec}] \copsh\PP A \ar[dd, hook, "\widebreve{\ip_A^\PP}"] \\
        & A \ar[ul, "\yo_A^\PP"] \ar[ur, "\oy_A^\PP"'] \ar[dl, "\yo_A^{\PP'}"'] \ar[dr, "\oy_A^{\PP'}"]&\\
        \psh{\PP'} A \ar[rr, bend left=15, "\ca{O}'", ""{anchor=center,name=O_}] && \ar[ll, bend left=15, "\Spec'", ""{anchor=center,name=Spec_}] \copsh{\PP'} A \\
        \ar[from=O, to=Spec, adj] \ar[from=O_, to=Spec_, adj]
    \end{tikzcd}\]
    We want to show that the inner and outer square commute up to isomorphism, that is: $\widearc{\pi_A^\PP} \circ \Spec \simeq \Spec' \circ \widebreve{\ip_A^\PP}$ et $\widebreve{\ip_A^\PP} \circ \ca{O} \simeq \ca{O}' \circ \widearc{\pi_A^\PP}$. We show that the first one is isomorphic to $\oy_A^\PP \plext \yo_A^{\PP'}$ and $\oy_A^\PP \prext \yo_A^{\PP'}$, and that the second is isomorphic to $\yo_A^\PP \prext \oy_A^{\PP'}$ and $\yo_A^\PP \plext \oy_A^{\PP'}$.
        \begin{align*}
            \copsh{\PP'} A(\ca{O}'\circ\widearc{\pi_A^\PP}, 1) &= \copsh{\PP'}A(\ca{O}', 1)(\widearc{\pi_A^\PP}, 1)
                                                               & \psh{\PP'}A(\Spec'\circ\widebreve{\ip_A^\PP}, 1) &= \psh{\PP'}A(\Spec', 1)(\widebreve{\ip_A^\PP}, 1)\\
                &\simeq (\ip_A^{\PP'}\rext\psh{\PP'} A(1, \yo_A^{\PP'}))(\widearc{\pi_A^\PP}, 1) & &\simeq (\pi_A^{\PP'}\rift\copsh{\PP'}A(\oy_A^{\PP'}, 1))(\widebreve{\ip_A^\PP}, 1)\\
                &\simeq \ip_A^{\PP'}\rext\psh{\PP'}A(\widearc{\pi_A^\PP}, \yo_A^{\PP'}) & &\simeq \ip_A^{\PP'}\rift\psh{\PP'}A(\oy_A^{\PP'}, \widebreve{\ip_A^\PP})\\
                &\simeq \ip_A^{\PP'}\rext(A(1, 1)\rift\pi_A^\PP) & &\simeq \ip_A^{\PP'}\rift(A(1, 1)\rext\pi_A^\PP)\\
                &\simeq (\ip_A^{\PP'}\rext A(1, 1))\rift\pi_A^\PP & &\simeq \pi_A^{\PP'}\rift\copsh\PP A(\oy_A^\PP, 1)\\
                &\simeq \copsh{\PP'}A(\oy_A^{\PP'}, 1)\rift\psh\PP A(\yo_A^\PP, 1) & &\simeq \psh{\PP'}A(\yo_A^{\PP'}, 1)\rift\copsh\PP A(\oy_A^\PP, 1)\\
            \copsh{\PP'}A(\widebreve{\ip_A^\PP}\circ\ca{O}, 1) &= \copsh{\PP'} A(\widebreve{\ip_A^\PP}, 1)(\ca{O}, 1)
                                                               & \psh{\PP'}A(\widearc{\pi_A^\PP}\circ\Spec, 1) &= \psh{\PP'}A(\widearc{\pi_A^\PP}, 1)(\Spec, 1)\\
                &\simeq (\ip_A^{\PP'}\rext\ip_A^\PP)(\ca{O}, 1) & &\simeq (\pi_A^{\PP'}\rift\pi_A^\PP)(\Spec, 1)\\
                &\simeq \ip_A^{\PP'}\rext\ip_A^\PP(\ca{O}, 1) & &\simeq \pi_A^{\PP'}\rift\pi_A^\PP(1, \Spec)\\
                &\simeq \ip_A^{\PP'}\rext\psh\PP A(1, \yo_A^\PP) & &\simeq \psh\PP A(\yo_A^\PP, 1)\rift\copsh{\PP'}A(\oy_A^\PP, 1)\\
                &\simeq \ip_A^{\PP'}\rext(A(1, 1)\rift\pi_A^\PP) & \psh{\PP'}A(1, \widearc{\pi_A^\PP}\circ\Spec) &= \psh{\PP'}A(1, \widearc{\pi_A^\PP})(1, \Spec)\\
                &\simeq (\ip_A^{\PP'}\rext A(1, 1))\rift\pi_A^\PP & &\simeq (\pi_A^\PP\rift\pi_A^{\PP'})(1, \Spec)\\
                &\simeq \copsh{\PP'}A(\oy_A^{\PP'}, 1)\rift\psh\PP A(\yo_A^\PP, 1) & &\simeq \pi_A^\PP(1, \Spec)\rift\pi_A^{\PP'}\\
            \copsh{\PP'} A(1, \widebreve{\ip_A^\PP}\circ\ca{O})
                &\simeq \copsh{\PP'}A(1, \widebreve{\ip_A^\PP})(1, \ca{O}) & &\simeq \copsh\PP A(\oy_A^\PP, 1)\rift\pi_A^{\PP'}\\
                &\simeq (\ip_A^\PP\rext\ip_A^{\PP'})(1, \ca{O}) & &\simeq (\ip_A^\PP\rext A(1, 1))\rift\pi_A^{\PP'}\\
                &\simeq \ip_A^\PP(\ca{O}, 1)\rext\ip_A^{\PP'} & &\simeq \ip_A^\PP \rext (A(1, 1)\rift\pi_A^{\PP'})\\
                &\simeq \psh\PP A(1, \yo_A^\PP)\rext\copsh{\PP'}A(1, \oy_A^{\PP'}) & &\simeq \copsh\PP A(1, \oy_A^{\PP'})\rext\psh{\PP'}A(1, \yo_A^{\PP'})
        \end{align*}
\end{proof}

\subsection{Reflexive completion}

We now want to study the fixed point of the Isbell adjunction. We will closely follow Avery and Leinster's article \cite{averyleinster}. One issue is that we don't have a notion of size nor of $\cat{V}$ or $\cat{Set}$ in a virtual equipment, so it is not possible to consider small functors. We chose to use $\PP$ as a measure for size. We will see in Subsection \ref{subsec:enriched} that, in the case of (possibly enriched) 1-categories, taking $\PP := \set{\text{small }\cat{V}\!-\!\text{presheaves on }A}$, we recover the results from \cite{averyleinster}.

We fix an Isbell $\PP$-admissible object $A$.

\begin{definition}[{\cite[Definition 4.6]{averyleinster}}]
    \index{arrow!$\PP$-small}\index{arrow!$\PP$-cosmall}
    A vertical arrow $f : A \to B$ is \textbf{$\PP$-small} if $N_f := B(f, 1) : A \dfrom B$ is in $\PP^\simeqq$. Dually, it is $\PP$-cosmall if $N^f := B(1, f) : B \dfrom A$ is in $\PP^\simeqq$. We will write $n_f^\PP := \widearc{N_f}$ and $n^f_\PP := \widebreve{N^f}$.
\end{definition}

\begin{lemma}[{\cite[Lemma 4.7]{averyleinster}}]
    \label{lem:small_comp_stable}
    If for every $p, q\in\PP, p\odot_L q\in\PP^\simeqq$, the collection of $\PP$-small (resp. $\PP$-cosmall) vertical arrows is closed under composition.
\end{lemma}

\begin{proof}
    Consequence of \ref{th:comp_rest}.
\end{proof}

\begin{remark}
    The minimal hypothesis is actually for the composition of representable (resp. corepresentables) to be in $\PP^\simeqq$.
\end{remark}

\begin{definition}\hfill
    \index{arrow!reflexive}\index{reflexive!arrow}\index{reflexive!completion}
    \begin{itemize}
        \item A vertical arrow $x_A : X \to \psh\PP A$ is \textbf{reflexive} if the unit $\eta : A \To \Spec\circ\ca{O}$ induces an isomorphism $\eta_{x_A} : x_A \To \Spec\circ\ca{O}\circ x_A$.
        \item A \textbf{reflexive completion} of $A$ is a fully faithful, reflexive arrow $\ca{r}_A^\PP : \ca{R}^\PP(A) \to \psh\PP A$ such that every reflexive arrow factorizes uniquely through $r_A^\PP$.
            \[\begin{tikzcd}
                \ar[dr, dashed] X \ar[rr, "x_A"] && \psh\PP A\\
                                                 & \ca{R}^\PP(A) \ar[ur, hook, "\ca{r}_a^\PP"']
            \end{tikzcd}\]
    \end{itemize}
\end{definition}

Since we are interested in the reflexive completion, we now suppose that a reflexive completion $\ca{r}_A^\PP$ exists for $A$.

\begin{lemma}
    $\yo_A^\PP$ is reflexive.
\end{lemma}

\begin{proof}
    \begin{align*}
        \Spec\circ\ca{O}\circ\yo_A^\PP &\simeq\widearc{\copsh\PP A(\oy_A^\PP, 1)}\circ\widebreve{\psh\PP A(1, \yo_A^\PP)}\circ\yo_A^\PP\\
                                       &\simeq\widearc{\copsh\PP A(\oy_A^\PP, 1)}\circ\widebreve{\psh\PP A(\yo_A^\PP, \yo_A^\PP)}
                                       \simeq\widearc{\copsh\PP A(\oy_A^\PP, 1)}\circ\widebreve{A(1, 1)}
                                       \simeq\widearc{\copsh\PP A(\oy_A^\PP, 1)}\circ\oy_A^\PP\\
                                       &\simeq\widearc{\copsh\PP A(\oy_A^\PP, \oy_A^\PP)}
                                       \simeq\widearc{A(1, 1)}
                                       \simeq \yo_A^\PP
    \end{align*}
\end{proof}

\begin{proposition}[{\cite[Definition 7.1]{averyleinster}}]
    A $\PP$-small (resp. $\PP$-cosmall) vertical arrow $f$ is dense (resp. codense) if and only if $n_f^\PP$ (resp. $n^f_\PP$) is fully faithful.
\end{proposition}

\begin{proof}
    $B(1, 1) \simeq \psh\PP A(n_f^\PP, n_f^\PP) \simeq B(f, 1) \rext B(f, 1)$. The dual is identical.
\end{proof}

\begin{lemma}[{\cite[Lemma 7.7]{averyleinster}}]
    \label{lem:ff_nerve_yo}
    $n_f^\PP \circ f \simeq \yo_A^\PP$ if $f$ is fully faithful.
\end{lemma}

\begin{proof}
    $n_f^\PP \circ f = \widearc{N_f} \circ f = \widearc{B(1, f)} \circ f = \widearc{B(f, f)} \simeq \widearc{A(1, 1)} = \yo_A^\PP$
\end{proof}

\begin{definition}[{\cite[Definition 7.10]{averyleinster}}]
    \index{arrow!adequate}\index{arrow!$\PP$-adequate}
    $f$ is \textbf{adequate} if it is dense, codense and fully faithful. It is \textbf{$\PP$-adequate} if it is adequate, $\PP$-small and $\PP$-cosmall.
\end{definition}

\begin{lemma}[{\cite[Lemma 7.13]{averyleinster}}]
    \label{lem:adequate_comp_closure}
    Adequate arrows and $\PP$-adequate arrows are respectively closed under composition.
\end{lemma}

\begin{proof}
    Let $A \xrightarrow{f} B \xrightarrow{g} C$ be adequate arrows. $g$ preserves weighted limits colimits thanks to the Lemma \ref{lem:ff_dense_prelim} and its dual, so the Lemma \ref{lem:comp_dense} (and its dual) gives the adequacy of $gf$. \ref{lem:small_comp_stable} gives the $\PP$-adequacy.
\end{proof}

\begin{remark}[{\cite[Lemma 7.14]{averyleinster}}]
    \label{rem:adequate_monic}
    A consequence of Lemma \ref{lem:codense_lim_equiv} is that if $f, g, g'$ are adequate, $gf \simeq g'f \implies g \simeq g'$.
\end{remark}

\begin{lemma}[{\cite[Lemma 7.15]{averyleinster}}]
    \label{lem:comp_small}
    Let $A \xrightarrow{f} B \xrightarrow{g} C$ with $g$ is fully faithful and $gf$ dense and $\PP$-small. Then $f$ is $\PP$-small.
\end{lemma}

\begin{proof}
    $\pi_A^\PP(1, n_{gf}^\PP\circ g) \simeq \pi_A^\PP(1, n_{gf}^\PP)(1, g) \simeq C(gf, g) \simeq B(f, 1) = N_f$
\end{proof}

\begin{corollary}[{\cite[Proposition 7.17]{averyleinster}}]
    \label{cor:comp_adequate}
    Let $A \xrightarrow{f} B \xrightarrow{g} C$ such that $g$ is fully faithful and $gf$ is $\PP$-adequate. Then $g$ is adequate and $f$ is $\PP$-adequate.
\end{corollary}

\begin{proof}
    Lemma \ref{lem:comp_small} and its dual.
\end{proof}

\begin{corollary}[{\cite[Corollary 7.16]{averyleinster}}]
    \label{cor:yo_small_adequate}
    $\ca{J}_A^\PP : A \to \ca{R}^\PP(A)$ induced by $\yo_A^\PP : A \to \psh\PP A$ is $\PP$-adequate.
\end{corollary}

\begin{proof}
    We consider the diagram
    \[\begin{tikzcd}[column sep=huge]
        && \psh\PP A\\
        A \ar[rrd, "\oy_A^\PP"', bend right, hook] \ar[rru, "\yo_A^\PP", bend left, hook] \ar[r, "\exists!\ca{J}_A^\PP", dashed] & \ca{R}^\PP(A)\ar[ur, "\ca{r}_A^\PP", hook] \ar[dr, hook] &\\
        && \copsh\PP A
    \end{tikzcd}\]
    Knowing that $\yo_A^\PP$ is $\PP$-small and dense and that $\oy_A^\PP$ is $\PP$-cosmall and codense, the corollary \ref{cor:comp_adequate} gives that $\ca{J}_A^\PP$ is $\PP$-adequate.
\end{proof}

\begin{corollary}[{\cite[Corollary 7.18]{averyleinster}}]
    Let $i : B\hookrightarrow\ca{R}^\PP(A)$ a fully faithful arrow such that there's an arrow $f: A \to B$ with $i\circ f = \ca{J}_A^\PP$. Then $f$ is $\PP$-adequate.
\end{corollary}

\begin{proof}
    Consequence of \ref{cor:yo_small_adequate} and \ref{cor:comp_adequate}.
\end{proof}

\begin{lemma}[{\cite[Lemma 8.1]{averyleinster}}]
    \label{lem:switch_nerves}
    Let $f : A\to B$ fully faithful, dense and $\PP$-small. Then $\ca{O}\circ n_f^\PP \simeq n^f_\PP$.
\end{lemma}

\begin{proof}
    \begin{align*}
        \ca{O}\circ n_f^\PP &= \widebreve{\psh\PP A(1, \yo_A^\PP)}\circ\widearc{B(f, 1)}\\
                            &= \widebreve{\psh\PP A(\widearc{B(f, 1)}, \yo_A^\PP)}\\
                            &\simeq \widebreve{A(1, 1)\rift B(f, 1)} & (\ref{prop:rift_psh})\\
                            &\simeq \widebreve{B(f, f)\rift B(f, 1)} & (f\text{ f.f})\\
                            &\simeq \widebreve{(B(f, 1)\rift B(f, 1))(1,f)} &(\ref{lem:ext_lift:res})\\
                            &\simeq \widebreve{B(1, f)} \simeq n^f_\PP & (f\text{ dense})
    \end{align*}
\end{proof}

\begin{proposition}[{\cite[Proposition 8.2]{averyleinster}}]
    \label{prop:refl_iff_adequate}
    Let $f : A \to B$ fully faithful, dense and $\PP$-small. $n_f^\PP$ is reflexive if and only if $f$ is $\PP$-adequate.
\end{proposition}

\begin{proof}\hfill
    \begin{itemize}
        \item[\underline{$\implies$}:] Suppose that $n_f^\PP$ is reflexive. Then $n_f^\PP\circ f\simeq \yo_A^\PP$ by Lemma \ref{lem:ff_nerve_yo}. We have the following diagram:
            \[\begin{tikzcd}[column sep=huge, row sep=large]
                &&\psh\PP A\\
                A \ar[urr, "\yo_A^\PP", bend left, hook] \ar[rr, "\ca{J}_A^\PP"', bend right, hook] \ar[r, "f", hook] & B \ar[ur, hook, "n_f^\PP"] \ar[r, dashed, "\exists!\ca{r}_f^\PP"'] & \ca{R}^\PP(A) \ar[u, hook, "\ca{r}_A^\PP"']
            \end{tikzcd}\]
            Where $\ca{r}_f^\PP$ is the factorization of $n_f^\PP$ through $\ca{r}_A^\PP$. We have $\ca{r}_A^\PP\circ\ca{J}_A^\PP = \yo_A^\PP$ and $\ca{r}_A^\PP\circ\ca{r}_f^\PP = n_f^\PP$ therefore $\ca{r}_A^\PP\circ\ca{r}_f^\PP\circ f \simeq \ca{r}_A^\PP\circ\ca{J}_A^\PP$ and by fully faithfulness of $\ca{r}_A^\PP$, we get $\ca{r}_f^\PP\circ f \simeq \ca{J}_A^\PP$. By Corollary \ref{cor:yo_small_adequate}, $\ca{J}_A^\PP$ is $\PP$-adequate so by \ref{cor:comp_adequate} $f$ is $\PP$-adequate.
        \item[\underline{$\impliedby$}:] Suppose that $f$ is $\PP$-adequate. Then Lemma \ref{lem:switch_nerves} and its dual gives $\Spec\circ\ca{O}\circ n_f^\PP \simeq \Spec\circ n^f_\PP \simeq n_f^\PP$.
    \end{itemize}
\end{proof}

\begin{corollary}[{\cite[Corollary 8.3]{averyleinster}}]
    \label{cor:rf_existence_ff}
    Let $f : A \to B$ be $\PP$-adequate arrow. Then there exists a unique $\ca{r}_f^\PP : B \to \ca{R}^\PP(A)$ up to isomorphism, such that the following diagram commutes up to isomorphism:
    \[\begin{tikzcd}[column sep=huge]
        &&\psh\PP A\\
        B \ar[urr, "n_f^\PP", hook, bend left] \ar[r, "\ca{r}_f^\PP", hook] \ar[drr, "n^f_\PP"', hook, bend right] & \ca{R}^\PP(A) \ar[ur, hook, "\ca{r}_A^\PP"'] \ar[dr, hook, "\ca{r}^A_\PP"] &\\
        && \copsh\PP A
    \end{tikzcd}\]
    Moreover, $\ca{r}_f^\PP$ is fully faithful.
\end{corollary}

\begin{theorem}[{\cite[Theorem 8.4]{averyleinster}}]
    \label{th:pu_reflexive_completion}
    Let $f : A \to B$ $\PP$-adequate. Then $\ca{r}_f^\PP : B \to \ca{R}^\PP(A)$ is adequate and the unique arrow up to isomorphism such that the following diagram commutes up to isomorphism:
    \[\begin{tikzcd}[row sep=large]
        B \ar[rr, "\ca{r}_f^\PP"] && \ca{R}^\PP(A)\\
        & A \ar[ul, "f"] \ar[ur, hook, "\ca{J}_A^\PP"'] &
    \end{tikzcd}\]
\end{theorem}

\begin{proof}
    The commutativity was checked in the proof of Proposition \ref{prop:refl_iff_adequate}, Corollaries \ref{cor:comp_adequate} and \ref{cor:yo_small_adequate} give the $\PP$-adequacy of $\ca{r}_f^\PP$. Remark \ref{rem:adequate_monic} and the last two corollaries give the unicity up to isomorphism.
\end{proof}

\begin{lemma}[{\cite[Lemma 8.8]{averyleinster}}]
    Let $f : A \to B$ $\PP$-adequate, such that $\ca{r}_f^\PP$ is $\PP$-adequate. Then for all $g : B \hookrightarrow C$ fully faithful such that $gf$ is $\PP$-adequate, then $g$ is $\PP$-adequate.
\end{lemma}

\begin{proof}
    Let $g : B \hookrightarrow C$ fully faithful such that $gf$ is $\PP$-adequate. We construct the diagram
    \[\begin{tikzcd}[row sep=large]
        C \ar[rrr, "\ca{r}_{gf}^\PP"] &&& \ca{R}^\PP(A)\\
                                      & \ar[ul, hook, "g"] B \ar[urr, "\ca{r}_f^\PP"'] &&\\
                                      && \ar[ul, hook, "f"] A \ar[uur, "\ca{J}_A^\PP"'] &\\
    \end{tikzcd}\]
    \begin{itemize}
        \item By Corollary \ref{cor:rf_existence_ff}, $\ca{r}_{gf}^\PP$ exists.
        \item By Corollary \ref{cor:comp_adequate}, $g$ is adequate.
        \item By Lemma \ref{lem:adequate_comp_closure}, $\ca{r}_{gf}^\PP\circ g$ is adequate.
        \item By Theorem \ref{th:pu_reflexive_completion}, $\ca{r}_f^\PP\circ f\simeq\ca{J}_A^\PP$.
    \end{itemize}
    And $\ca{r}_{gf}^\PP\circ g\circ f\simeq\ca{J}_A^\PP$ i.e. $\ca{r}_{gf}^\PP\circ g\circ f\simeq\ca{r}_f^\PP\circ f$, so $\ca{r}_{gf}^\PP\circ g\simeq\ca{r}_f^\PP$ using Theorem \ref{th:pu_reflexive_completion} and Remark \ref{rem:adequate_monic}. By hypothesis, $\ca{r}_f^\PP$ is $\PP$-adequate so, by Corollary \ref{cor:comp_adequate}, $g$ is $\PP$-adequate.
\end{proof}

We now study the functoriality of the reflexive completion. We fix $A$, $B$ two Isbell $\PP$-admissible objects having reflexive completions $\ca{r}_A^\PP, \ca{r}_B^\PP$.

\begin{lemma}[{\cite[Lemma 9.2]{averyleinster}}]
    \label{lem:refl_comp_func}
    Let $f : A \to B$ $\PP$-adequate, and write $\ca{R}^\PP(A) := \ca{r}_{\ca{J}_B^\PP\circ f}^\PP$. Then the squares commute up to isomorphism:
    \[\begin{tikzcd}
        \psh\PP A & \ar[l, "\widearc{\pi_B^\PP(f, 1)}"'] \psh\PP B\\
        \ca{R}^\PP(A) \ar[u, hook, "\ca{r}_A^\PP"] & \ar[l, "\ca{R}^\PP(f)"] \ca{R}^\PP(B) \ar[u, hook, "\ca{r}_B^\PP"']
    \end{tikzcd}
    \hspace{3em}
    \begin{tikzcd}
        \copsh\PP A & \ar[l, "\widebreve{\ip_B^\PP(1, f)}"'] \copsh\PP B\\
        \ca{R}^\PP(A) \ar[u, hook] & \ar[l, "\ca{R}^\PP(f)"] \ca{R}^\PP(B) \ar[u, hook]
    \end{tikzcd}\]
\end{lemma}

\begin{proof}
    \begin{align*}
        r_A^\PP\circ\ca{R}^\PP(f) &= \ca{r}_A^\PP\circ\ca{r}_{\ca{J}_B^\PP\circ f}^\PP\\
                                  &\simeq n_{\ca{J}_B^\PP\circ f}^\PP\\
                                  &\simeq \widearc{\ca{R}^\PP(B)(\ca{J}_B^\PP\circ f, 1)}\\
                                  &\simeq \widearc{\psh\PP A(\ca{r}_B^\PP\circ\ca{J}_B^\PP\circ f, \ca{r}_B^\PP)}\\
                                  &\simeq \widearc{\psh\PP A(\yo_B^\PP\circ f, \ca{r}_B^\PP)}\\
                                  &\simeq \widearc{\psh\PP A(\yo_B^\PP, 1)(f, 1)(1, \ca{r}_B^\PP)}\\
                                  &\simeq \widearc{\pi_B^\PP(f, 1)}\circ\ca{r}_B^\PP
    \end{align*}
\end{proof}

\begin{theorem}[{\cite[Theorem 9.3]{averyleinster}}]
    \label{th:small_ad_equiv}
    Let $f : A \to B$ $\PP$-adequate. The following are equivalent:
    \begin{enumerate}
        \item $\ca{r}_f^\PP$ is $\PP$-adequate.
        \item $\ca{R}^\PP(f)$ is $\PP$-adequate.
        \item $\ca{R}^\PP(f)$ is an equivalence.
    \end{enumerate}
    In this case, $\ca{r}_{\ca{r}_f^\PP}^\PP$ is defined and is a pseudo-inverse of $\ca{R}^\PP(f)$.
\end{theorem}

\begin{proof}
    $\ca{R}^\PP(f) \circ \ca{J}_B^\PP$ is adequate and such that $\ca{R}^\PP(f) \circ \ca{J}_B^\PP \simeq \ca{r}_f^\PP$.
    \begin{itemize}
        \item[\underline{(c)$\implies$(b)}:] Immediate.
        \item[\underline{(b)$\implies$(a)}:] Using the equation above and \ref{lem:adequate_comp_closure}.
        \item[\underline{(a)$\implies$(c)}:] We consider the diagram
            \[\begin{tikzcd}[row sep=huge,column sep=large]
                \ca{R}^\PP(A) \ar[r, bend right=10, "\ca{r}_{\ca{r}_f^\PP}^\PP"'] & \ar[l, bend right=10, "\ca{R}^\PP(f)"'] \ca{R}^\PP(B)\\
                A \ar[u, hook, "\ca{J}_A^\PP"] \ar[r, "f"'] & \ar[ul, "\ca{r}_f^\PP"] B \ar[u, hook, "\ca{J}_B^\PP"']
            \end{tikzcd}\]
            (a) defines $\ca{r}_{\ca{r}_f^\PP}\PP$, the bottom triangle commutes thanks to Theorem \ref{th:pu_reflexive_completion} whereas the top ones commute thanks to the Theorem and the equation above. The outer square commutes by definition of $\ca{R}^\PP(f)$. We have:
            \begin{align*}
                \ca{R}^\PP(f)\circ\ca{r}_{\ca{r}_f^\PP}^\PP\circ\ca{J}_A^\PP &\simeq\ca{R}^\PP(f)\circ\ca{J}_B^\PP\circ f\simeq\ca{r}_f^\PP\simeq\ca{J}_A^\PP\\
                \ca{r}_{\ca{r}_f^\PP}^\PP\circ\ca{R}^\PP(f)\circ\ca{J}_B^\PP &\simeq\ca{r}_{\ca{r}_f^\PP}^\PP\circ\ca{r}_f^\PP\simeq\ca{J}_B^\PP
            \end{align*}
            We conclude with \ref{rem:adequate_monic} that $\ca{r}_{\ca{r}_f^\PP}^\PP$ and $\ca{R}^\PP(f)$ are pseudo-inverses.
    \end{itemize}
\end{proof}

\begin{corollary}[{\cite[Corollary 9.4]{averyleinster}}]
    \label{cor:rr_equiv_r}
    If $\ca{R}^\PP(A)$ is Isbell $\PP$-admissible and has a reflexive completion, there's an equivalence
    \[\begin{tikzcd}[column sep=large]
        \ca{R}^\PP(A) \ar[r, bend left, "\ca{J}_{\ca{R}^\PP(A)}^\PP"] & \ar[l, bend left, "\ca{R}^\PP(\ca{J}_A^\PP)"] \ca{R}^\PP(\ca{R}^\PP(A))
    \end{tikzcd}\]
\end{corollary}

\begin{proof}
    We take $f = \ca{J}_A^\PP$ in Theorem \ref{th:small_ad_equiv} and notice that $\ca{r}_{\ca{J}_A^\PP}^\PP \simeq 1_{\ca{R}^\PP(A)}$.
\end{proof}

\begin{corollary}[{\cite[Corollary 9.5]{averyleinster}}]
    If $\ca{R}^\PP(A)$ is Isbell $\PP$-admissible and has a reflexive completion, the following are equivalent:
    \begin{enumerate}
        \item $\ca{J}_A^\PP : A \hookrightarrow \ca{R}^\PP(A)$ is an equivalence.
        \item Every reflexive $x : J \to \psh\PP A$ factorizes through $\yo_A^\PP$.
        \item There exists an object $B$ such that $A \simeq \ca{R}^\PP(B)$
    \end{enumerate}
\end{corollary}

\begin{proof}
    We first prove (a)$\iff$(b): using the diagram from Corollary \ref{cor:yo_small_adequate}, let $x : B \to \psh\PP A$ be a reflexive arrow, factorizing as $\ca{r}_A^\PP\circ\tilde{x}$ and $\ca{J}_A^{\PP-1}$ a pseudo inverse of $\ca{J}_A^\PP$. Then
    $$\yo_A^\PP\circ\ca{J}_A^{\PP-1}\circ\tilde{x} \simeq \ca{r}_A^\PP\circ\ca{J}_A^\PP\circ\ca{J}_A^{\PP-1}\circ\tilde{x}\simeq\ca{r}_A^\PP\circ\tilde{x}=x$$
    Now suppose that every reflexive $x : J \to \psh\PP A$ factorizes through $\yo_A^\PP$. Take $x=\ca{r}_A^\PP$, the factorization gives a pseudo-inverse of $\ca{J}_A^\PP$. Finally, Corollary \ref{cor:rr_equiv_r} gives the equivalence with (c).
\end{proof}

\begin{remark}
    The existence of a reflexive completion of the reflexive completion is not needed for (a)$\iff$(b).
\end{remark}

\begin{definition}
    An object $A$ is \textbf{reflexively complete} if it is equivalent to a reflexive completion.
\end{definition}
\todo{sert à rien, pas utilisé après}

We finish the section by studying the interaction with (co)limits.

\begin{lemma}[{\cite[Lemma 11.1]{averyleinster}}]\hfill
    \begin{itemize}
        \item $\ca{J}_A^\PP : A \hookrightarrow \ca{R}^\PP(A)$ preserves and reflects (weighted) limits and colimits.
        \item $\ca{r}_A^\PP : \ca{R}^\PP(A) \hookrightarrow \psh\PP A$ preserves limits and reflects (weighted) limits and colimits.
    \end{itemize}
\end{lemma}

\begin{proof}
    Reflection comes from the full faithfulness of both arrows.
    \begin{itemize}
        \item $\ca{J}_A^\PP$ is $\PP$-adequate by \ref{cor:yo_small_adequate}, so it preserves limits and colimits by \ref{lem:ff_dense_prelim} and its dual.
        \item $\ca{r}_A^\PP\circ\ca{J}_A^\PP = \yo_A^\PP$ is dense so $\ca{r}_A^\PP$ is dense by \ref{lem:ff_comp_dense} so by \ref{lem:ff_dense_prelim} it preserves limits.
    \end{itemize}
\end{proof}

\begin{proposition}[{\cite[Proposition 10.1 i]{averyleinster}}]
    $\ca{R}^\PP(A)$ is stable under absolute colimit that $\psh\PP A$ admits.
\end{proposition}

\begin{proof}
    Let $x : J \to \psh\PP A$ be a reflexive arrow, $\vec{p}$ a weight such that $\colim^{\vec{p}}x$ exists and is absolute. Then $\Spec\circ\ca{O}\circ\colim^{\vec{p}}x \simeq \colim^{\vec{p}}(\Spec\circ\ca{O}\circ x) \simeq \colim^{\vec{p}}x$ so $\colim^{\vec{p}}x$ is reflexive and factorizes through $\ca{R}^\PP(A)$.
\end{proof}


\section{Cocompleteness and cocompletions}

\subsection{Atomicity}

\begin{definition}
    Let $p : Y \dfrom \cdots \dfrom Z$ a weight. A distributor $j : A \dfrom E$ is a $\vec p$-atom if $E$ admits $\vec p$colimits and these colimits are $j$-absolute. It is a $\Phi$-atom for $\Phi$ a class of weights when it is a $\vec p$-atom for all $\vec p\in\Phi$. An arrow is a $\vec p$-atom (resp. $\Phi$-atom) when its conjoint is.
\end{definition}

\begin{theorem}
    Let $j : A \to E$ an arrow, and $\ell\dashv_j r$ a relative adjunction. Then
    \begin{enumerate}
        \item If $j$ is a $\Phi$-atom and $r$ is $\Phi$-cocontinuous, then $\ell$ is a $\Phi$-atom.
        \item If $j$ is dense and $\ell$ is a $\Phi$-atom then $r$ is $\Phi$-cocontinuous.
    \end{enumerate}
\end{theorem}

\begin{proof}
    Let $\vec p : X \dfrom \cdots \dfrom Z$ be a weight in $\Phi$ and $g : X \to C$ an arrow. Then
    \begin{enumerate}
        \item \begin{align*}
                C(\ell, \colim^{\vec p}g) &\simeqq E(j, r\colim^{\vec p}g)\\
                                          &\simeqq E(j, \colim^{\vec p}(rg))\\
                                          &\simeqq E(j, rg)\odot_L \vec p\\
                                          &\simeqq C(\ell, g)\odot_L \vec p
            \end{align*}
        \item Since $\ell$ is a $\Phi$-atom, $\Phi$-colimits are $\ell$-absolute and we conclude by Lemma \ref{lem:radj_right_labs_colim}.
    \end{enumerate}
\end{proof}

\begin{definition}
    \index{distributor!monic}
    \label{def:monic}
    A distributor $j : A \dfrom E$ is \textbf{monic} when for all $f, f' : X \to C$, $j(1, f) = j(1, f') \implies f = f'$.
\end{definition}

\begin{proposition}
    For $j : A \dfrom E$, the following are equivalent:
    \begin{enumerate}
        \item There exists a class $\PP$ of distributors such that $j(1, f)\odot_Lj\in\PP^\simeqq$ for each $f : A \to E$ and such that $j : A \dfrom E$ exhibits a $\PP$-presheaf object.
        \item $j$ is a dense and monic $j$-atom.
    \end{enumerate}
    In that case, $A$ is $\PP$-admissible if and only if $j$ is represented by a fully faithful arrow $A \to E$.
\end{proposition}

\begin{proof}
    By definition of a presheaf object, $j$ is dense and monic, and every $j$-colimit exists and is $j$-absolute by Proposition \ref{prop:psh_colims}.

    Conversely, define $\PP := \set{j(1, f) \mid f:X\to E}$. By construction, it is closed under precomposition by arrows into $E$ and for each $f : A \to E$ we have $j(1, \colim^jf)\simeqq j(1, f)\odot_L j$ because $\colim^jf$ is $j$-absolute. Then the repletion of $\PP$ is closed under the specified left-composites. Define $\pi_A^\PP := j : A \dfrom E$. It is dense by assumption and given $j(1, f) : A \dfrom X$, we define $\widearc{j(1, f)} = f$. Thus $\pi_A^\PP(1, \widearc{j(1, f)}) = j(1, f)$ and the choice is unique by monicity of $j$.

    If $A$ is $\PP$-admissible, then there exists a fully faithful arrow $A \to E$ representing $\pi_A^\PP$. If $j$ is represented by a fully faithful arrow $f$, then $j(1, f) \simeq E(f, f) \simeqq A(1, 1)$ and $j(1, f)\in\PP^\simeqq$.
\end{proof}

\subsection{Rank}

\begin{definition}
    Let $j : A \to E$ be a arrow.
    \begin{enumerate}
        \item An arrow $f : E \to I$ \textbf{has rank $j$} if the left extension $j\plext (fj)$ exits and the cell $j\plext(jf)\To f$ is invertible.
        \item A distributor $ p : E \dfrom I$ \textbf{has rank $j$} if the right lit $p(j, 1)\rift E(j, 1)$ exists and the cell $p \To p(1, j)\rift E(j, 1)$ is invertible.
    \end{enumerate}
    In particular, $f : E \to I$ has rank $j$ when its conjoint does.
\end{definition}

\begin{proposition}
    \label{prop:rank_conditions}
    Let $j : A \to E$ be an arrow and $f : E \to I$ (resp. $p : E \dfrom I$). Consider the following
    \begin{enumerate}
        \item $f$ has rank $j$ (resp. $p$ has rank $j$).
        \item $f$ is a left extension along $j$ (resp. $p$ is a right lift through $E(j, 1)$).
        \item $f$ preserves $j$-absolute colimits (resp. $p$ respects $j$-absolute colimits).
    \end{enumerate}
    Then $(a) \implies (b) \implies (c)$, if $j$ is fully faithful then $(b) \implies (a)$ and if it is also dense, then all three are equivalent.
\end{proposition}

\begin{proof}
    It suffices to prove the proposition for a distributor. $(a) \implies (b)$ is evident and $(b) \implies (c)$ follows from \ref{lem:rift_absolute}. Suppose $j$ fully faithful and write $p = p' \rift E(j, 1)$. Then
    \begin{align*}
        p &\simeqq p' \rift E(j, 1)\\
          &\simeqq (p' \rift A(1, 1))\rift E(j, 1)\\
          &\simeqq (p' \rift E(j, j))\rift E(j, 1)\\
          &\simeqq (p' \rift E(j, 1))(j, 1)\rift E(j, 1)\\
          &\simeqq p(j, 1)\rift E(j, 1)
    \end{align*}
    Suppose that $j$ is also dense, then $\id_E$ is the $j$-absolute left extension $j\plext j$, so
    $$E(j, j\plext j) \simeqq E(j, 1) \simeqq E(j, j) \odot_L E(j, 1)$$
    Hence $p$ preserves $j\plext j$ so
    $$p\simeqq p(j\plext j, 1)\simeqq p(j, 1)\rift E(j, 1)$$
\end{proof}

\subsection{Completions}

\begin{definition}
    \index{object!$\Phi$-cocomplete} \index{arrow!$\Phi$-cocontinuous} \index{distributor!$\Phi$-exact}
    Let $\Phi$ be a class of weights.
    \begin{enumerate}
        \item A \textbf{$\Phi$-colimit} is a $\vec p$-colimit for $\vec p\in\Phi$.
        \item An object $A$ is  \textbf{$\Phi$-cocomplete} if it admits every $\vec p$-colimit.
        \item An arrow $f$ with $\Phi$-cocomplete domain is \textbf{$\Phi$-cocontinuous} if it preserves every $\Phi$-colimit.
        \item A distributor $p$ with $\Phi$-cocomplete codomain is $\Phi$-exact if it respects every $\Phi$-colimit.
    \end{enumerate}
\end{definition}

\begin{remark}
    Note that an arrow with $\Phi$-cocomplete domain is $\Phi$-cocontinuous if and only if its conjoint is $\Phi$-exact.
\end{remark}

\begin{definition}
    \index{colimit equivalent}
    Two classes of weights $\Phi, \Phi'$ are \textbf{colimit equivalent} $\Phi \colimeq \Phi'$ if and only if $\Phi \colimleq \Phi'$ and $\Phi' \colimleq \Phi$ where $\Phi \colimleq \Phi'$ if and only if:
    \begin{enumerate}
        \item Every $\Phi'$-cocomplete object is $\Phi$-cocomplete.
        \item Every $\Phi'$-cocontinuous arrow with $\Phi'$-cocomplete domain is $\Phi$-cocontinuous.
        \item Every $\Phi'$-exact distributor with $\Phi'$-cocomplete domain is $\Phi$-exact.
    \end{enumerate}
\end{definition}

This relation defines a preorder on the collection of classes of weights. We are interested in the maximal elements.

\begin{definition}
    \index{class of weights!colimit saturation}
    Let $\Phi$ be a class of weights, the \textbf{colimit saturation} $\Phi^{\colim}$ of $\Phi$ is the largest class such that $\Phi^{\colim} \colimeq \Phi$. It is \textbf{colimit saturated} is $\Phi = \Phi^{\colim}$.
\end{definition}

We establish of few properties that will be used when linking cocompletions and presheaf objects.

\begin{lemma}
    $(-)^{\colim}$ defines a closure operator whose fixed points are the colimit saturated class of weights. For a class of weights $\Phi$, we have the following:
    \begin{enumerate}
        \item $\Phi^{\colim}$ is replete.
        \item $B(1, f)\in\Phi^{\colim}$ for each arrow $f : A \to B$.
        \item Given two compatible weights $\vec p, \vec q \in\Phi^{\colim}, (p_1, \cdots, p_m, q_1, \cdots, q_n)\in\Phi^{\colim}$.
        \item If $\vec p\in\Phi^{\colim}$ and $p_1\odot_L\cdots\odot_L p_i$ exists (for $1\leq i\leq m$) then $(p_1\odot_L \cdots\odot_L p_i, p_{i+1}, \cdots, p_m)\in\Phi^{\colim}$.
        \item If $\vec p\in\Phi^{\colim}$ and $p_i\odot\cdots\odot p_j$ exists (for $1\leq i\leq j\leq m$) then $(p_1, \cdots, p_{i-1}, p_i\odot \cdots\odot p_j, p_{j+1}, \cdots, p_m)\in\Phi^{\colim}$.
        \item If $p\in\Phi^{\colim}, p(1, f)\in\Phi^{\colim}$
    \end{enumerate}
    Finally, $\Phi\colimeq\Phi'$ if and only if $\Phi^{\colim} \colimeq\Phi'^{\colim}$.
\end{lemma}

\begin{proof}
    \begin{enumerate}
        \item Repletion comes from essential uniqueness of colimits.
        \item For every arrow $g : B \to C$, its $B(1, f)$-colimit is $g\circ f$.
        \item Follows from \ref{lem:ext_lift:stab_split}.
        \item Follows from \ref{lem:ext_lift:stab_wcomp}.
        \item Follows from \ref{lem:ext_lift:stab_comp}.
        \item Follows from $p(1, f) = p\odot B(1, f)$ and the previous items; or \ref{lem:co_lim:res}.
    \end{enumerate}
\end{proof}

\begin{definition}
    \index{$\Phi$-cocompletion}
    \label{def:cocomp}
    Let $\Phi$ be a class of weights and $A$ an object. A \textbf{$\Phi$-cocompletion} of $A$ consists of an object $\coc\Phi A$ and an arrow $\coce\Phi A \to \coc\Phi A$ such that:
    \begin{enumerate}
        \item $\coc\Phi A$ is $\Phi$-cocomplete.
        \item For each arrow $f : A \to X$ with $\Phi$-cocomplete codomain, the left extension $\coce\Phi A \plext f : \coc\Phi A \to X$ exists, the unit $f \To (\coce\Phi A \plext f) \circ \coce\Phi A$ is invertible and $\coce\Phi A\plext f$ is the essentially unique $\Phi$-cocontinuous arrow $\widetilde{f} : \coc\Phi A\to X$ such that $f\simeqq \widetilde{f} \circ \coce\Phi A$.
        \item For each distributor $p : A \dfrom X$, the right lift $p\rift\coc\Phi A (\coce\Phi A, 1) : \coc\Phi A \dfrom X$ exists, the cell $(p\rift\coc\Phi A(\coce\Phi A, 1))(\coce\Phi A, 1) \To p$ is invertible and $p\rift\coc\Phi A(\coce\Phi A, 1)$ is the essentially unique $\Phi$-exact distributor $\widetilde{p} : \coc\Phi A\dfrom X$ such that $p\simeqq \widetilde{p}(\coce\Phi A, 1)$.
    \end{enumerate}
    An object \textbf{$\Phi$-cocompletable} if it admits a $\Phi$-cocompletion.
\end{definition}

\begin{remark}
    Consequently, the conjoint of $\coce\Phi A$ is always lifting. Plus, since an arrow is $\Phi$-cocontinuous if and only if its conjoint is $\Phi$-exact, the condition of $\Phi$-cocontinuous arrows can be replaced with
    \begin{itemize}
        \item If $p : A \dfrom X$ is corepresentable and $X$ is $\Phi$-cocomplete, then $p\rift\coc\Phi A(\coce\Phi A, 1)$ is corepresentable.
    \end{itemize}
    The notion of cocompletion is compatible with the earlier discussion on colimit equivalence: for $\Phi$ and $\Phi'$ be classes of weights such that $\Phi\colimeq\Phi'$, an arrow $j : A \to E$ exhibits the $\Phi$-cocompletion of $A$ if and only if it exhibits the $\Phi'$-cocompletion of $A$.
\end{remark}

\begin{lemma}
    \label{lem:cocomp_ff}
    Every $\Phi$-cocompletion $\coce\Phi A : A \to \coc\Phi A$ is fully faithful.
\end{lemma}

\begin{proof}
    By definition there are isomorphisms
    $$p\rift\coc\Phi A(\coce\Phi A, \coce\Phi A) \simeqq (p\rift\coc\phi A(\coce\Phi A, 1))(\coce\Phi A, 1) \simeqq p \simeqq p\rift A(1, 1)$$
    Hence a natural bijection between cells $A(1, 1)\To p$ and $\coc\Phi A(\coce\Phi A, \coce\Phi A) \To p$, implying $A(1, 1)\simeqq\coc\Phi A(\coce\Phi A, \coce\Phi A)$, so $\coce\Phi A$ is fully faithful.
\end{proof}

\begin{lemma}
    Let $\Phi$ be a class of weights and $A$ a $\Phi$-cocompletable object. Then $\widetilde{A(1, 1)}\simeqq\coc\Phi A(1, \coce\Phi A)$.
\end{lemma}

\begin{proof}
    $\coc\Phi A(1, \coce\Phi A)(\coce\Phi A, 1) \simeqq \coc\Phi A(1, 1)$ by Lemma \ref{lem:cocomp_ff}, and we conclude by essential uniqueness.
\end{proof}

\begin{lemma}
    \label{lem:cocomp_leftext_abs}
    For $\coce\Phi A : A \to \coc\Phi A$ a $\Phi$-cocompletion, left extensions along $\coce\Phi A$ and $\Phi$-colimits in $\coc\Phi A$ are $\coce\Phi A$-absolute.
\end{lemma}

\begin{proof}
    %FIXME: exact or cocontinuous ?
    Every rift lift through $\coc\Phi A(\coce\Phi A, 1)$ is $\Phi$-exact, and we conclude by \ref{lem:lifting_abs_iff_lift_respect}.
\end{proof}

\begin{proposition}
    Let $\coce\Phi A : A \to \coc\Phi A$ a $\Phi$-cocompletion. For $f : \coc\Phi A \to X$, the following are equivalent:
    \begin{enumerate}
        \item $f$ is $\Phi$-cocontinuous.
        \item $f$ has rank $\coce\Phi A$
        \item $f$ is a pointwise left extension along $\coce\Phi A$.
        \item $f$ preserves $\coce\Phi A$-absolute colimits.
    \end{enumerate}
    For $p : \coc\Phi A\dfrom X$, the following are equivalent:
    \begin{enumerate}
        \item $p$ is $\Phi$-exact.
        \item $p$ has rank $\coce\Phi A$.
        \item $p$ is a right lift through $\coc\Phi A(\coce\Phi A, 1)$.
        \item $p$ respects $\coce\Phi A$-absolute colimits.
    \end{enumerate}
\end{proposition}

\begin{proof}
    $\coce\Phi A$ is fully faithful by Lemma \ref{lem:cocomp_ff}, so Proposition \ref{prop:rank_conditions} gives $(b) \iff (c) \implies (d)$. $(d) \implies (a)$ is given by Lemma \ref{lem:cocomp_leftext_abs} since all $\Phi$-colimits in $\coc\Phi A$ are $\coce\Phi A$-absolute.

    For $p$ a $\Phi$-exact distributor, the right lift $p(\coce\Phi A, 1)\rift\coc\Phi A(\coce\Phi A, 1)$ is the unique $\Phi$-cocontinuous $\widetilde{p(\coce\Phi A, 1)}$ such that $p(\coce\Phi A, 1) \simeqq \widetilde{p(\coce\Phi A, 1)}(\coce\Phi A, 1)$ exhibiting $p$ as a right lift through $\coc\Phi A(\coce\Phi A, 1)$, therefore $(a) \implies (c)$. For an arrow, take $p = X(f, 1)$ which exhibits $f$ as the pointwise left extension $\coce\Phi A \plext (f\coce\Phi A)$.
\end{proof}

\begin{corollary}
    \label{cor:cocomp_dense}
    Every $\Phi$-cocompletion $\coce\Phi A : A \to \coc\Phi A$ is dense.
\end{corollary}

\begin{proof}
    $\id_{\coc\Phi A}$ is $\Phi$-cocontinuous, hence has rank $\coce\Phi A$, hence $\coce\Phi A \plext \coce\Phi A \simeqq \id_{\coc\Phi A}$.
\end{proof}

\begin{lemma}
    \label{lem:cocomp_colimsat}
    Let $\coce\Phi A : A \to \coc\Phi A$ a $\Phi$-cocompletion. Then $\coc\Phi A(\coce\Phi A, 1)$ is in the colimit saturation of $\Phi$.
\end{lemma}

\begin{proof}
    $\Phi$-cocomplete objects admit left extensions along $\coce\Phi A$, $\Phi$-cocontinuous arrows preserve and $\Phi$-exact respect $\coce\Phi A$-absolute colimits, including left extensions along $\coce\Phi A$ by Lemma \ref{lem:cocomp_leftext_abs}.
\end{proof}

We deduce an alternative characterisation of cocompletions which we will use in Subsection \ref{subsec:psh_cocomp} to establish the link between presheaf objects and cocompletions.

\begin{theorem}
    An arrow $j : A \to E$ exhibits the $\Phi$-cocompletion of $A$ if and only if
    \begin{enumerate}
        \item $E$ is $\Phi$-complete.
        \item $E(j, 1)$ is in the colimit saturation of $\Phi$.
        \item $E(j, 1)$ is lifting and right lifts through $E(j, 1)$ are $\Phi$-exact.
        \item $j$ is dense and fully faithful.
    \end{enumerate}
\end{theorem}

\begin{proof}
    First assume that $j$ exhibits the $\Phi$-cocompletion of $A$, then (a) and (c) are by definition, (b) by Lemma \ref{lem:cocomp_colimsat} and (d) by Corollary \ref{cor:cocomp_dense} and Lemma \ref{lem:cocomp_ff}.

    Supposing (a) through (d) hold, we verify the conditions in Definition \ref{def:cocomp}. (a) is by assumption (a). For (b), since $E(j, 1)$ is in the colimit saturation of $\Phi$, the left extension $j \plext f : E \to X$ where $f$ is an arrow with $\Phi$-cocomplete codomain, exists and corepresents $X(f, 1)\rift E(j, 1)$. Finally for (c), given $p : A \dfrom X$, existence and $\Phi$-exactness of the right lifts exist by assumption (c) and the cell $(p\rift E(j, 1))(j, 1)\To p$ is invertible by full faithfulness of $j$ using Lemma \ref{lem:ff_res}. For $\widetilde p : E \dfrom E$ such that $p \simeqq \widetilde{p}(j, 1)$, and using density of $j$ (assumption (d)),  $\Phi$-exactness and $E(j, 1)$ being in the colimit saturation of $\Phi$, we have:
    $$\widetilde p \simeqq \widetilde{p}(j\plext j, 1) \simeqq \widetilde{p}(j, 1) \rift E(j, 1) \simeqq p\rift E(j, 1)$$
    Hence essential uniqueness.
\end{proof}

It is possible to get another characterisation of those objects that are equivalent to a $\Phi$-cocompletion with a useful criterion on adjointness.

\begin{proposition}
    Let $j : A \to E$ a dense arrow with $\Phi$-cocomplete codomain. An arrow $f : E \to X$ admits a right adjoint if and only if $(fj)\plext j$ exists.
\end{proposition}

\begin{proof}
    Since $j$ is dense, and by Lemma \ref{lem:pt_ext_props}, $(fj)\plext j \simeqq f \plext (j \plext j) \simeqq f \plext 1$. Being a $\Phi$-colimit, it is preserved by $f$ and we conclude by Corollary \ref{cor:adjoint_arrow_theorem}.
\end{proof}

\begin{corollary}
    \label{cor:cocomp_adj_iff_lext}
    Let $f : A \to X$ an arrow with $\Phi$-cocompletable domain and $\Phi$-cocomplete codomain. Then $\widetilde f : \coc\Phi A \to X$ admits a right adjoint if and only if $f\plext 1$ exists.
\end{corollary}

\begin{proposition}
    Let $A$ a $\Phi$-cocompletable object, $E$ an object. The following are equivalent:
    \begin{enumerate}
        \item $E\simeq\coc\Phi A$.
        \item There exists a dense and fully faithful $\Phi$-atom $j : A \to E$ with $\Phi$-cocomplete codomain for which $j \plext \coce\Phi A$.
    \end{enumerate}
\end{proposition}

\begin{proof}
    Take $j$ to be $\coce\Phi A$ which satisfies the hypotheses for Lemmas \ref{lem:cocomp_ff}, \ref{lem:cocomp_leftext_abs} and Corollary \ref{cor:cocomp_dense}. Conversely, by Corollary \ref{cor:cocomp_adj_iff_lext}, $\coce\Phi A\plext j\dashv j\plext\coce\Phi A$. Since the conditions of Lemma \ref{lem:left_witness_equiv} are satisfied, and $j$ is a $\Phi$-atom, $\coce\Phi A\plext j$ is $j$-absolute.
\end{proof}

More generally, we get a range a tests for $\Phi$-cocompleteness...

\begin{theorem}
    \label{th:cocomp_tests}
    The following are equivalent for an object $A$
    \begin{enumerate}
        \item $A$ is $\Phi$-cocomplete.
        \item $A$ is reflective in a $\Phi$-cocomplete object.
    \end{enumerate}
    And if $A$ admits a $\Phi$-cocompletion $\coce\Phi A : A \to \coc\Phi A$, the following are also equivalent
    \begin{enumerate}[resume]
        \item For every object $X$ with a $\Phi$-cocompletion $\coce\Phi X$ and each $g : X \to A$, the left extension $\coce\Phi X \plext g : X \to \coc\Phi X$.
        \item The left extension $\coce\Phi A \plext 1_A : \coc\Phi A \to C$ exists.
        \item For every weight $\vec p\in\Phi^{\colim}$, $A$ admits the colimit $\colim^{\vec p}A_1$.
        \item $\coce\Phi A : A \to \coc\Phi A$ admits a left adjoint.
        \item $A$ is reflective in a $\Phi$-cocompletion.
    \end{enumerate}
\end{theorem}

\begin{proof}
    $(a)\implies (b)\land(c)\land(e)$, $(c)\implies(d)$ and $(f)\implies(g)\implies(b)$ are immediate. $(b)\implies(a)$ by Lemma \ref{lem:refl_adj_colims}, $(e)\implies(d)$ by Lemma \ref{lem:cocomp_colimsat}, and $(d)\implies(e)$ because $\coce\Phi A\plext 1_A$ is the left adjoint.
    \begin{align*}
        A(\coce\Phi A\plext 1_A, 1) &\simeqq A(1, 1)\rift \coc\Phi A(\coce\Phi A, 1) &\\
                                    &\simeqq \coc\Phi A(\coce\Phi 1, \coce\Phi A)\rift\coc\Phi A(\coce\Phi A, 1) & \text{(Lemma \ref{lem:cocomp_ff})}\\
                                    &\simeqq \coc\Phi A(1, \coce\Phi A) & \text{(Corollary \ref{cor:cocomp_dense})}
    \end{align*}
\end{proof}

...And the corresponding characterisation for $\Phi$-cocontinuity:

\begin{corollary}
    Let $A$ and $B$ $\Phi$-cocompletable objects. The following are equivalent for an arrow $f : A \to B$:
    \begin{enumerate}
        \item $f$ is $\Phi$-cocontinuous.
        \item Given a pseudocommutative square of arrow on the left where the vertical arrows exhibit reflective inclusions, the mate on the right is invertible.
            \[\begin{tikzcd}
                B \ar[r, "b", ""{anchor=center,name=t}] & B'\\
                \ar[u, hook, "r"] A \ar[r, "a"', ""{anchor=center,name=b}] & A' \ar[u, hook, "r'"']
                \ar[from=t, to=b, phantom, "\simeqq"{anchor=center}]
            \end{tikzcd}
            \hspace{5em}
            \begin{tikzcd}
                \ar[d, "\ell"'] B \ar[r, "b"] & \ar[dl, Rightarrow] B' \ar[d, "\ell'"]\\
                A \ar[r, "a"'] & A'
            \end{tikzcd}\]
    \end{enumerate}
    If $A$ admits a $\Phi$-cocompletion $\coce\Phi A : A \to \coc\Phi A$ then the following are equivalent:
    \begin{enumerate}[resume]
        \item For every object $X$ admitting a $\Phi$-cocompletion $\coce\Phi X : X \to \coc\Phi X$ and each $g : X \to A$, the left extension $\coce\Phi X \plext g : \coc\Phi X\to A$ is preserved by $f$.
        \item The left extension $\coce\Phi A\plext 1_A$ is preserved by $f$.
        \item For every weight $\vec p\in\Phi^{\colim}$, the weighted colimit $\colim^{\vec p}1_A$ is preserved by $f$.
    \end{enumerate}
    Furthermore, if $B$ admits a $\Phi$-cocompletion $\coce\Phi B : B \to \coc\Phi B$ then the following are equivalent:
    \begin{enumerate}[resume]
        \item As in (b), but $r = \coce\Phi A$, $r' = \coce\Phi B$ and $b = \coc\Phi f := \coce\Phi A\plext (\coce\Phi B\circ a)$.
        \item As in (b), but $B$ and $B'$ are $\Phi$-cocompletion and $b\simeqq r\plext(r'\circ a)$.
    \end{enumerate}
\end{corollary}

\begin{proof}
    The proof is the similar as in Theorem \ref{th:cocomp_tests}. The hardest is $(d)\implies(f)$, we have
    $$(\coce\Phi A\plext1_B)\circ(\coce\Phi A\plext(\coce\Phi B\circ a)) \simeqq \coce\Phi A\plext ((\coce\Phi B\plext 1_B)\circ\coce\Phi B\circ a)\simeqq \coce\Phi A\plext a$$
    Thus the cell witnessing preservation of $\coce\Phi A\plext 1_A$ is the mate.
\end{proof}

\subsection{Presheaf objects as cocompletions}
\label{subsec:psh_cocomp}

\begin{theorem}
    \label{th:psh_cocomp}
    Let $\Phi$ be a class of weights, $\PP$ a class of distributors and $A$ a $\PP$-admissible object. Suppose that
    \begin{enumerate}
        \item For every weight $\vec p\in\Phi$ and distributor $q\in\PP$, the left composite $q\odot_L\vec p$ and is in $\PP^\simeqq$.
        \item $\pi_A^\PP$ is in $\Phi^{\colim}$.
        \item $\pi_A^\PP : A \dfrom \psh\PP A$ is lifting.
    \end{enumerate}
    Then $\yo_A^\PP : A \to \psh\PP A$ exhibits the $\Phi$-cocompletion of $A$.
\end{theorem}

\begin{proof}
    We use Theorem \ref{th:cocomp_chararct}:
    \begin{enumerate}
        \item $\psh\PP A$ is $\Phi$-cocomplete by Proposition \ref{prop:psh_colims}.
        \item $\psh\PP A(\yo_A^\PP, 1)$ is in the colimit saturation since it contains $\pi_A^\PP$ and is replete.
        \item $\pi_A^\PP$ is lifting and right lifts through $\pi_A^\PP$ are $\Phi$-exact by Lemma \ref{lem:lifting_abs_iff_lift_respect} where $\pi_A^\PP$-absoluteness of $\Phi$-colimits is obtained with Proposition \ref{prop:psh_colims}.
        \item $\yo_A^\PP$ is dense and fully faithful by Proposition \ref{prop:monic_dense_atom_equiv}
    \end{enumerate}
\end{proof}

Restricting to a class of unary weights, we get

\begin{corollary}
    Let $\PP$ be a class of distributors closed under left-composites and $A$ a $\PP$-admissible object. Then $\psh\PP A$ is $\PP$-cocomplete and every $\PP$-colimit in $\psh\PP A$ is $\PP$-absolute. If $\pi_A^\PP : A \dfrom \psh\PP A$ is lifting, then $\yo_A^\PP$ exhibits the $\PP$-cocompletion of $A$.
\end{corollary}

\begin{proof}
    Apply respectively Proposition \ref{prop:psh_colims} and Theorem \ref{th:psh_cocomp}.
\end{proof}

<++>

\section{Instantiations}
\label{sec:instances}

\subsection{Enriched categories}
\label{subsec:enriched}

%FIXME: Rewrite everything
Many constructions rely on the existence of additional structure of the homsets of categories: additive or linear categories, dg-categories, 2-categories... It is possible to axiomatize this approach: we \textit{enrich} a category. Although it is possible, as was done in the article \cite{arkor2026presheavescocompletionsformalcategory}, to enrich a category in a virtual bicategory (or even a virtual double category), we choose to stick to monoidal categories (as was done in \cite{arkor2024formaltheoryrelativemonads}) as it is more common although less general.

Most of the theory of enriched category theory can be found in \cite{kelly}.

\begin{definition}
    \index{category!monoidal}
    A \textbf{monoidal category} $(\cat{V}, \otimes, I, a, l, r)$, usually only written $\cat{V}$ comprises of the following data
    \begin{itemize}
        \item A category $\cat{V}$.
        \item A functor $\cat{V}\times\cat{V} \to \cat{V}$.
        \item An object $I$ of $\cat{V}$, the \textbf{unit}.
        \item Three natural isomorphisms $a, l, r$ with signatures $a_{XYZ} : (X \otimes Y) \otimes Z \to X \otimes (Y \otimes Z)$, $l_X : I \otimes X \to X$ and $r_X : X \otimes I \to I$ for $X, Y, Z$ objects in $\cat{V}$.
    \end{itemize}
    Such that the following diagrams commute
    \[\begin{tikzcd}[sep=scriptsize]
        {(X \otimes I) \otimes Y} \ar[rr, "a_{XIY}"]
        \ar[rdd, "r \otimes \id_Y"'] && {X \otimes (I \otimes Y)}
        \ar[ldd, "\id_X \otimes l"]\\
        \\
        & {X \otimes Y}
    \end{tikzcd}\]
    \[\begin{tikzcd}[sep=scriptsize]
        & {((W \otimes X) \otimes Y) \otimes Z} & \\
        \\
        {(W \otimes (X \otimes Y)) \otimes Z} && {(W \otimes X) \otimes (Y \otimes Z)} \\
        \\
        {W \otimes ((X \otimes Y) \otimes Z)} && {W \otimes (X \otimes (Y \otimes Z))}
        \arrow["{a_{WXY} \otimes \id_Z}", from=1-2, to=3-1]
        \arrow["{a_{(W \otimes X)YZ}}"', from=1-2, to=3-3]
        \arrow["{a_{W(X \otimes Y)Z}}", from=3-1, to=5-1]
        \arrow["{a_{WX(Y \otimes Z)}}"', from=3-3, to=5-3]
        \arrow["{\id_W \otimes a_{XYZ}}", from=5-1, to=5-3]
    \end{tikzcd}\]
\end{definition}

\begin{definition}
    \index{category!$\cat{V}$-enriched}
    A \textbf{$\cat{V}$-category} $\cat{C}$ (or \textbf{$\cat{V}$-enriched category}) comprises of the following data
    \begin{itemize}
        \item A class $\ob\cat{C}$, the \textit{objects} of $\cat{C}$.
        \item An object $\cat{C}(A, B)$ of $\cat{V}$ for $A, B$ objects of $\cat{C}$.
        \item A morphism $M : \cat{C}(B, C) \otimes \cat{C}(A, B)$, the \textit{composition law}.
        \item An \textit{identity} morphism $i_A : I \to \cat{C}(A, A)$.
    \end{itemize}
    Such that the associativity and gnitality diagrams are commutative:
    \[\begin{tikzcd}[sep=scriptsize]
        {(\cat{C}(C, D) \otimes \cat{C}(B, C)) \otimes \cat{C}(A, B)} && {\cat{C}(C, D) \otimes (\cat{C}(B, C) \otimes \cat{C}(A, B))} \\
        \\
        {\cat{C}(B, D)\otimes \cat{C}(A, B)} && {\cat{C}(C, D) \otimes \cat{C}(A, C)} \\
        \\
        & {\cat{C}(A, D)}
        \arrow["a", from=1-1, to=1-3]
        \arrow["{M\otimes \id}", from=1-1, to=3-1]
        \arrow["{\id \otimes M}", from=1-3, to=3-3]
        \arrow["M"', from=3-1, to=5-2]
        \arrow["M", from=3-3, to=5-2]
    \end{tikzcd}\]
    \[\begin{tikzcd}[sep=scriptsize]
        {\cat{C}(B, B) \otimes \cat{C}(A, B)} && {\cat{C}(A, B)} && {\cat{C}(A, B) \otimes \cat{C}(A, A)} \\
        \\
        {I \otimes \cat{C}(A, B)} &&&& {\cat{C}(A, B) \otimes I}
        \arrow["M", from=1-1, to=1-3]
        \arrow["M"', from=1-5, to=1-3]
        \arrow["{i_B \otimes \id}", from=3-1, to=1-1]
        \arrow["l"', from=3-1, to=1-3]
        \arrow["r", from=3-5, to=1-3]
        \arrow["{\id \otimes i_A}"', from=3-5, to=1-5]
    \end{tikzcd}\]
\end{definition}

\begin{example}
    Common examples include: $\cat{Set}$-enriched categories (known as locally small categories), $\cat{Ab}$-enriched categories (known as additive categories), $\cat{Vect}_k$ (known as $k$-linear categories), $\cat{Ch}_k$-enriched categories (known as dg-categories) and $\cat{Cat}$-enriched categories (known as 2-categories). More exotic examples include $\cat{2}$-enriched categories (posets) and $\overline{\bb{R}^+}$ (metric or ultrametric spaces).
\end{example}

\begin{definition}
    \index{$\cat{V}$-enriched!category}\index{$\cat{V}$-enriched!natural transformation}
    A \textbf{$\cat{V}$-functor} $F : \cat{A} \to \cat{B}$ comprises of a function $F : \ob\cat{A} \to \ob\cat{B}$ and for all $A, B$ objects of $\cat{A}$, a morphism $F_{AB} : \cat{A}(A, B) \to \cat{B}(FA, FB)$ de $\cat{V}$ such that $F$ is compatible with composition and units:
    \[\begin{tikzcd}[sep=small]
        \cat{A}(B, C) \otimes \cat{A}(A, B) \ar[rr, "M"] \ar[dd, "F \otimes F"']
            && \cat{A}(A, C) \ar[dd, "F"]\\
        \\
        \cat{B}(FB, FC) \otimes \cat{B}(FA, FB) \ar[rr, "M"] && \cat{B}(FA, FC)
    \end{tikzcd}
    \begin{tikzcd}[sep=small]
        && \cat{A}(A, A) \ar[dddd, "F"]\\
        \\
        I \ar[uurr, "i_A"] \ar[ddrr, "i_{FA}"]\\
        \\
        && \cat{B}(FA, FA)
    \end{tikzcd}\]
    A \textbf{$\cat{V}$-natural transformation} $\alpha : F \to G$ comprises of morphisms $\alpha_A : I \to \cat{B}(FA, GA)$ such that:
    \[\begin{tikzcd}
        & I \otimes \cat{A}(A, B) \ar[r, "\alpha_B\otimes F"] & \cat{B}(FB, GB) \otimes \cat{B}(FA, FB) \ar[rd, start anchor=east, "M"]\\
        \cat{A}(A, B) \ar[ru, "l^{-1}"] \ar[rd, "r^{-1}"] &&& \cat{B}(FA, GB)\\
        & \cat{A}(A, B) \otimes I \ar[r, "G \otimes \alpha_A"] & \cat{B}(GA, GB) \otimes \cat{B}(FA, GA) \ar[ru, start anchor=east, "M"]
    \end{tikzcd}\]
\end{definition}

It is common to add hypothesis on the enriching category: for example, asking for $\otimes$ to be symmetric, asking for $(-\otimes X)$ to have a right adjoint (viewed as the "internal hom" out of $X$) or completeness/cocompleteness. We refrain from doing so unless necessary.

\begin{definition}
    The virtual double category $\cat{V}-\bb{Cat}$ of $\cat{V}$-enriched categories has
    \begin{itemize}
        \item \textbf{objects:} $\cat{V}$-enriched categories.
        \item \textbf{arrows:} $\cat{V}$-functors.
        \item \textbf{distributors:} $\cat{V}$-distributors. A $\cat{V}$-distributor $p : X \dfrom Y$ comprises of an object $p(x, y)$ for all objects $x$ of $X$ and $y$ of $Y$, natural contravariantely in $x$ and covariantely in $y$ and compatible with composition and identities.
        \item \textbf{cells:} $\cat{V}$-natural transformations. A cell
            \[\begin{tikzcd}
                C_0 \ar[d, "f"'] & \ar[l, loose, "p_1"'] \cdots & \ar[l, loose, "p_n"'] C_n \ar[d, "g"]\\
                D && \ar[ll, loose, "q"] D'
            \end{tikzcd}\]
            Is the data of a $\cat{V}$-natural morphism
            $$p_1(x_0, x_1) \otimes \cdots \otimes p_n(x_{n-1}, x_n) \to q(fx_0, gx_n)$$
            With a few commutative diagrams that can be found in \cite[Definition 8.1]{arkor2024formaltheoryrelativemonads}.
    \end{itemize}
\end{definition}

\begin{proposition}
    This virtual double category is a virtual equipment.
\end{proposition}

\begin{proof}
    The loose identity on $A$ is given by $(x, y) \mapsto A(x, y)$ and the restriction by $p(f, g)(x, y) = p(fx, gy)$. The associated cell is the identity.
\end{proof}

We want to explicitly describe right lifts of $\cat{V}$-distributors.

\begin{definition}[{\cite[Definition 8.5]{arkor2024formaltheoryrelativemonads}}]
    A $\cat{V}$-presheaf on a $\cat{V}$-category $C$ is the data of an object $p(c)$ of $\cat{V}$ for each object $c$ of $C$ compatible with "morphisms" of $C$ : $C(d, c) \otimes p(c) \to p(d)$ for all $c, d$ in $C$ compatible with the identity and composition in the evident way.

    For $p, q$ $\cat{V}$-presheaves on $C$ and $v$ an object of $C$, a family of morphisms $\set{\phi_c : p(c)\otimes v \to q(c)}_{c\in\ob(C)}$ is $\cat{V}$-natural in $c$ when for all $c'$, the following diagram commutes
    \[\begin{tikzcd}
        C(c', c)\otimes p(c)\otimes v \dar \rar & C(c', c)\otimes q(c) \dar\\
        p(c')\otimes v \rar & q(c')
    \end{tikzcd}\]
    If there exists an object $\ca{P}C(p, q)$ of $\cat{V}$ with a universal $\cat{V}$-natural family, we call $\ca{P}C(p, q)$ the object of $\cat{V}$-natural transformations from $p$ to $q$.
\end{definition}

\begin{remark}
    When $\cat{V}$ is symmetric and closed, a $\cat{V}$-presheaf is just a $\cat{V}$-functor $C^\op \to \cat{V}$ and in that case, $\ca{P}C(p, q) \simeqq [C^\op, \cat{V}](p, q)$ (when it exists).
\end{remark}

\begin{lemma}[{\cite[Lemma 8.7]{arkor2024formaltheoryrelativemonads}}]
    Let $p : Z \dfrom Y, q : Y \dfrom X$ be $\cat{V}$-distributors. If the objects $\ca{P}C(p(-, x), q(-, y))$ for all $x\in\ob(X), y\in\ob(Y)$, they form the right lift $q\rift p : Y \dfrom X$.
\end{lemma}

\subsection{Internal and locally internal categories}
\label{subsec:internal}

We refer to \cite[Chapter 8]{borceux1} and \cite[Chapter 2]{johnstone} for an elementary treatment, and to \cite{betti} for a formal treatment. Overall, this subsection is heavily influenced by the work of Renato Betti, especially \cite{bettiwalters_closedbicats}, \cite{bettiwalters_completness}, \cite{betti}. We first use the characterization of internal categories as internal monads in the bicategories of spans.

\begin{definition}[{\cite[Definition 5.3.1]{leinster}}]
    Let $\mathbb{X}$ be a virtual double category. We construct the virtual double category of horizontal monads $\Mod(\bb{X})$:
    \begin{itemize}
        \item \textit{Objects:} $(X_0, X : X_0 \dfrom X_0, \mu_X : X, X \To X, \eta_X : \To X)$ such that $\mu(\mu, 1_X) = \mu(1_X, \mu)$ et $\mu(\eta, 1_X) = \mu(1_X, \eta) = 1_X$. We will usually just write $X$ for $(X_0, X, \mu_X, \eta_X)$.
        \item \textit{Arrows:} $f := (f_0 : X_0 \to Y_0, \omega_f) : X \to Y$ such that $\omega_f \circ \mu_X = \mu_Y(\omega_f, \omega_f)$ et $\omega_f \circ \eta_X = \omega_Y\circ f_0$.
            \[\begin{tikzcd}
                \ar[d, "f_0"'] X_0 & \ar[l, loose, "X"', ""{anchor=center,name=t}] X_0 \ar[d, "f_0"]\\
                Y_0 & \ar[l, loose, "Y", ""{anchor=center,name=b}] Y_0
                \ar[from=t, to=b, phantom, "\omega_f"{anchor=center}]
            \end{tikzcd}\]
        \item \textit{Distributors:} $p := (p : X_0 \dfrom Y_0, \chi, \varphi) : X \dfrom Y$ such that $\chi(1_p, \mu_Y) = \chi(\chi, 1_Y)$, $\chi(1_p, \eta_Y) = 1_p$ and dually $\varphi(\mu_X, 1_p) = \varphi(1_X, \varphi)$, $\varphi(\eta_X, 1_p) = 1_p$; et $\varphi(1_X, \xi) = \xi(\varphi, 1_Y)$.
            \[\begin{tikzcd}
                \ar[d, equal] X_0 & \ar[l, loose, "p"'] Y_0 & \ar[l, loose, "Y"'] Y_0 \ar[d, equal]\\
                X_0 & {} & \ar[ll, loose, "p"] Y_0
                \ar[from=1-2, to=2-2, phantom, "\chi"{anchor=center}]
            \end{tikzcd}
            \hspace{10pt}
            \begin{tikzcd}
                    \ar[d, equal] X_0 & \ar[l, loose, "X"'] X_0 & \ar[l, loose, "p"'] Y_0 \ar[d, equal]\\
                    X_0 & {} & \ar[ll, loose, "p"] Y_0
                    \ar[from=1-2, to=2-2, phantom, "\varphi"{anchor=center}]
            \end{tikzcd}\]
        \item \textit{Cell:} A cell in $\Mod(\bb{X})$ is given by one in $\bb{X}$ with frame
            \[\begin{tikzcd}
                \ar[d, "f_0"'] X_0 & \ar[l, loose, "p_1"'] \cdots & \ar[l, loose, "p_n"'] Y_0 \ar[d, "g_0"]\\
                X_0' & {} & \ar[ll, loose, "q"] Y_0'
                \ar[from=1-2, to=2-2, phantom, "\theta"{anchor=center}]
            \end{tikzcd}\]
            Such that $\theta(\varphi, 1\cdots 1) = \chi_q(\omega_f, \theta)$, $\theta(1\cdots 1, \chi) = \varphi_q(\theta, \omega_g)$ and for all $2 \leq i \leq n$, $\theta(1\cdots 1, \chi_{i-1}, 1\cdots 1) = \theta(1\cdots 1, \varphi_i, 1\cdots 1)$.
    \end{itemize}
\end{definition}

$\bb{X}$ having restrictions is a sufficient condition for $\Mod(\bb{X})$ to be a virtual equipment.

\begin{proposition}[{\cite[Proposition 5.5]{cs10}}]
    $\Mod(\bb{X})$ is normal and $A(1, 1)$ is given by $A$. That is, $\eta_A : \To A$ is opcartesian in $\Mod(\bb{X})$.
\end{proposition}

\begin{proposition}[{\cite[Proposition 7.4]{cs10}}]
    If $\bb{X}$ has all restrictions, then $\Mod(\bb{X})$ too, and the component of $p(f, g)$ is $p(f_0, g_0)$.
\end{proposition}

Returning to internal categories, fix $\cat{E}$ a category with pullbacks. We define $\Dist(\cat{E}) := \Mod(\Span(\cat{E}))$ and unravel the definitions:
\begin{itemize}
    \item An object is a monad in $\Span(\cat{E})$, or equivalently in the bicategory $\bSpan(\cat{E})$, that is, an \textbf{category internal to $\cat{E}$}.
    \item An arrow is a morphism preserving the monad structure, so an \textbf{internal functor}.
    \item A distributor $C \dfrom D$ is a diagram on $C^\op\times D$, so an \textbf{internal distributor}.
    \item A cell is an \textbf{internal natural transformation}.
\end{itemize}
 
Note that, when $\cat{E}$ has pullback stable reflexive coequalizers, we can compose horizontal arrows. As a sanity check, we will routinely check that the notions are the ones we know in $\Dist(E)$.

\begin{proposition}[{\cite[Proposition 4.1]{day_loccartclosed}}]
    When $\cat{E}$ has finite limits, $\bSpan\cat{E}$ has all right lifts and extensions (along/through unary weights) if and only if it is locally cartesian closed.
\end{proposition}

In particular, when $\cat{E}$ is an elementary topos, we can compose distributors and all rights and extensions exist.

\begin{remark}
    Note also that in the general case, since $\bSpan(\cat{E})$ is a symmetric bicategory, the existence of right extensions is equivalent to the existence of right liftings.
\end{remark}

\begin{definition}
    An internal category $C$ associated to a monad $\alpha$ is
    \begin{itemize}
        \item \textbf{discrete} if $\alpha \simeqq 1_{C_0}$.
        \item a \textbf{monoid} if $C_0$ is a terminal object.
        \item a \textbf{groupoid} if $\alpha$ is idempotent.
        \item a \textbf{group} if it is both a monoid and a groupoid.
    \end{itemize}
    Every object of $\bSpan(\cat{E})$, regarded with the trivial monad, is a \textbf{discrete internal category}. Switching the source and target arrows, we get the \textbf{opposite internal category} of $C$.
\end{definition}

Starting with something easy, for $f : C \to D$ an internal functor, the cell $\smallfrown_f$ looks like

\[\begin{tikzcd}
    C_0 \ar[d, equal] & \ar[l, "s"'] C_1 \ar[d, "{(s, f_1, t)}"] \ar[r, "t"] & C_0 \ar[d, equal]\\
    C_0 & \ar[l, "\pi_0"] C_0 \times_{D_0} D_1 \times_{D_0} C_0 \ar[r, "\pi_1"] & C_0
\end{tikzcd}\]

\begin{lemma}
    $f : C \to D$ is fully faithful if and only if $C_1$ is the iterated pullback
    \[\begin{tikzcd}
        \dar (C_0 \underset{D_0}{\times} D_1) \underset{D_1}{\times} (D_1 \underset{D_0}{\times} C_0) \rar \ar[dr, pullback] & \dar D_1 \underset{D_0}{\times} C_0 \rar \ar[dr, pullback] & C_0 \ar[d, "f_0"]\\
        \dar C_0 \underset{D_0}{\times} D_1 \rar \ar[dr, pullback] & \ar[d, "s"] D_1 \ar[r, "t"'] & D_0\\
        C_0 \ar[r, "f_0"'] & D_0
    \end{tikzcd}\]
\end{lemma}

To study internal categories, it is useful to introduce the notion of locally internal categories. We follow \cite{bettiwalters_closedbicats} and start with categories enriched in bicategories:

\begin{definition}[{\cite[Definition 1.2]{bettiwalters_closedbicats}}]
    \index{category!enriched in a bicategory}\index{$\ca B$-category}
    For $\ca{B}$ a bicategory, a \textbf{$\ca{B}$-enriched category} $X$ comprises of the following data:
    \begin{itemize}
        \item Objects $x, y, \cdots$
        \item For each object $x$, an underlying object $ex$ in $\ca B$.  We say that $x$ is an object over $ex$.
        \item For each pair $x, y$ of objects, an arrow $X(x, y) : ex \from ey$ in $\ca B$.
        \item For each object, a unit $1_{ex} \to X(x, x)$.
        \item A composition given by a cell $X(x, y) \cdot X(y, z) \to X(x, z)$ in $\ca B$.
    \end{itemize}
    and such that this data satisfy associativity and identity properties. A \textbf{$\ca{B}$-functor} $F : X \to Y$ is a function on objects preserving the underlying objects and a family of 2-cells
    $$X(x, y) \to Y(Fx, Fy)$$
    Satisfying the axioms of functors.
\end{definition}

\begin{remark}
    In the case of $\ca{B} := \bSpan(\cat{E})$, a $\ca{B}$-category with a single object is an internal category: the underlying object is the object of objects, the Hom arrow gives the source and target morphisms, the unit gives the identity morphism and the composition is the internal composition.
\end{remark}

\begin{definition}[{\cite[Definition 1.8]{bettiwalters_closedbicats}}]
    \label{def:intcat_res}
    For $X$ a $\ca B$-category, $x$ an object over $u$ and $f : v \to u$ a morphism in $\ca B$. A restriction $x_f$ of $x$ along $f$ is an object over $v$ such that
    \begin{align*}
        f \cdot X(x, y) &\simeqq X(x_f, y)\\
        X(y, x) \cdot f^\op &\simeqq X(y, x_f)
    \end{align*}
    $X$ has restrictions if for all $x$ and $f$, $x_f$ exists.
\end{definition}

We have two characterizations of a locally internal category.

\begin{definition}[{\cite[Definition 1.11, Proposition 1.13 and Proposition 4.3]{bettiwalters_closedbicats}}]
    A \textbf{locally internal category} is a $\bSpan(\cat{E})$-category with restrictions. This is equivalent to a fibration
    $$p : \cat F \to \cat E$$
    Such that the functor $[x, y] : \bSpan(\cat E)(px, py)^\op \to \cat{Set}$ taking $(f, g)$ to $\cat F_w(f^*x, g^*y)$ (where $w$ is the domain of $f$ and $g$) is representable for all $x$ and $y$.
\end{definition}

Note that in the case of $\bSpan\cat E$, the two equations are equivalent. The definition of module between $\ca B$-categories will coincide with the one of internal categories when viewed as one object $\ca B$-categories.

\begin{definition}
    Let $X$ and $Y$ be $\ca B$-categories. A \textbf{$\ca B$-distributor} $p : X \dfrom Y$ consists of an arrow $p(x, y) : ex \from ey$ in $\ca B$ for each $x$ object in $X$ and $y$ object in $Y$, with actions of $X$ and $Y$ :
    \begin{align*}
        X(x', x) \cdot p(x, y) &\to p(x', y)\\
        p(x, y) \cdot Y(y, y') &\to p(x, y')
    \end{align*}
    With associativity, unity and mixed associativity properties.
\end{definition}

Using this definition, we can see the isomorphisms of Definition \ref{def:intcat_res} as isomorphisms of distributors, seeing an arrow in $\ca B$ as a distributors between discrete internal categories.

\begin{lemma}
    Functors between $\bSpan(\cat E)$-categories preserve restrictions.
\end{lemma}

\begin{proof}
    Let $f : u \to v$ a map, $x$ an object over $v$ and $F : X \to Y$ a $\bSpan\cat E$-functor:
    $$Y(1, Fx_f) \simeqq Y(1, F) \cdot Y(1, x) \cdot f^\op \simeqq Y(1, (Fx)_f)$$
\end{proof}

As remarked in \cite{bettiwalters_completness}, when $\ca B$ is locally finitely complete and cocomplete, composites of the form $X \dfrom A \dfrom Y$ with $A$ having a finite number of objects (in particular, when $A$ is an internal category) always exist. Moreover, the composition of $p : X \dfrom Y$ and $q : Y \dfrom Z$, when it exists, is given by the coequalizer on the two actions in $\ca B(ex, ez)$:

\[\begin{tikzcd}
    \displaystyle{\sum_{y', y''}} p(x, y') \cdot Y(y', y'') \cdot q(y'', z) \ar[r, shift left] \ar[r, shift right]
    & \displaystyle{\sum_y} p(x, y) \cdot q(y, z)
\end{tikzcd}\]

When considering $\bSpan(\cat Set)$-categories (e.g. small categories), this is exactly the coend $\int^{y\in Y} q(y, z) \times p(x, y)$. We delay the explicit description of right extensions to the construction of presheaf objects.



\subsection{\texorpdfstring{$\infty$}{∞}-catégories}

\cite{riehl_verity}
\todo{problème des two-sided isofibrations discretes dans A\backslash Fib(K)/B}
\todo{remplacer par les \texorpdfstring{$\infty$}{∞}-équipements de \cite{ruit2025formalcategorytheoryinftyequipments} ?}

\subsection{Fibered Categories}
\todo{À jeter ?}

\cite{borceux2} \cite{benaboudistributorsfibrations}

\clearpage
\appendix

\section{Extension calculus cheatsheet}

We resume many results used in the calculus of extensions. We work in a virtual equipment $\bb{X}$, and every type is implicit.

\begin{itemize}
    \item $B(f, 1) \odot p \odot C(1, g) \simeq p(f, g)$
        \hfill Theorem \ref{th:comp_rest}
\end{itemize}

\begin{center}
    Definition \ref{def:colim}
\end{center}
\begin{multicols}{2}
    \begin{itemize}
        \item $X(\colim^{\vec{p}}f, 1) \simeq X(f, 1) \rift \vec{p}$
        \item $X(1, \lim^{\vec{p}}f) \simeq \vec{p} \rext X(1, f)$
    \end{itemize}
\end{multicols}

\begin{center}
    Definition \ref{def:pointwise}
\end{center}

\begin{multicols}{2}
    \begin{itemize}
        \item $Y(j \plext f, 1) \simeq Y(f, 1) \rift X(j, 1)$
        \item $Y(f\plift j, 1) \simeq Z(f, 1) \rext Z(j, 1) \simeq Z(j, f)$
        \item $Y(1, j\prext f) \simeq X(1, j) \rext Y(1, f)$
        \item $Y(1, f\prift j) \simeq Z(1, j) \rift Z(1, f) \simeq Z(f, j)$
    \end{itemize}
\end{multicols}

\begin{center}
    Lemma \ref{lem:ext_lift}
\end{center}

\begin{multicols}{2}
    \begin{itemize}
        \item $q \rift (p_1, \cdots, p_m) \simeq q \rift (p_1 \odot_L \cdots \odot_L p_i ,\cdots, p_m)$
        \item $q \rift (p_1, \cdots, p_m) \simeq q \rift (p_1, \cdots, p_i \odot \cdots \odot p_j, \cdots, p_m)$
        \item $(q \rift (p_1, \cdots, p_i)) \rift (p_{i+1}, \cdots, p_m) \simeq q \rift(p_1, \cdots, p_m)$
        \item $(q \rift \vec{p})(y, x) \simeq q(1, x) \rift (p_1, \cdots, p_m(1, y))$
        \item $(p_1, \cdots, p_m) \rext q \simeq (p_1, \cdots, p_i \odot_R \cdots \odot_R p_m) \rext q$
        \item $(p_1, \cdots, p_m) \rext q \simeq (p_1, \cdots, p_i \odot \cdots \odot p_j, \cdots, p_m) \rext q$
        \item $(p_1, \cdots, p_i) \rext ((p_{i+1}, \cdots, p_m) \rext q) \simeq (p_1, \cdots, p_m) \rext q$
        \item $(\vec{p} \rext q)(z, y) \simeq (p_1(y, 1), \cdots, p_m)\rext q(z, 1)$
    \end{itemize}
\end{multicols}
\vspace{-2\baselineskip+2\parsep}
$$\bullet\hspace{.5em}(\vec{p} \rext q) \rift \vec{r} \simeq \vec{p} \rext (q \rift \vec{r})$$

\begin{itemize}
    \item $\psh\PP A(\widearc{p}, \widearc{q}) \simeq q \rift p$
        \hfill Proposition \ref{prop:rift_psh}
    \item $\widearc{p(1, f)} = \widearc{p}\circ f$
\end{itemize}

\clearpage

\nocite{*}
\printbibliography[heading=bibintoc]

\clearpage
\printindex

\end{document}
